\documentclass[11pt,a4paper]{article}
\usepackage[margin=25mm]{geometry}
\usepackage{fontspec}
\setmainfont{Latin Modern Roman}
\setsansfont{Latin Modern Sans}
\setmonofont[Scale=0.87]{Latin Modern Mono}
\usepackage{amsmath,amssymb,amsthm,mathtools}
\usepackage{array,longtable,booktabs}
\usepackage{microtype}
\usepackage[dvipsnames]{xcolor}
\usepackage{hyperref}
\hypersetup{colorlinks=true,linkcolor=MidnightBlue,citecolor=MidnightBlue,urlcolor=MidnightBlue,pdftitle={AN GRH: Anabelian point separation and effective Galois bounds},pdfauthor={AN GRH research audit}}
\usepackage{enumitem}
\setlist{itemsep=2pt,topsep=5pt}
\setlength{\parindent}{1.2em}
\setlength{\parskip}{3pt}
\setlength{\emergencystretch}{2em}
\renewcommand{\arraystretch}{1.15}
\newcommand{\et}{\mathrm{\acute et}}
\title{\textbf{AN\_GRH}\\[7pt]Anabelian point separation, cyclotomic certificates,\\and effective Galois bounds}
\author{Research audit of the supplied programme sources\\GPT-6 Astra Pro; run of 21 November 2026}
\date{21 November 2026 (Europe/Berlin)}
\begin{document}
\maketitle
\begin{abstract}
We reconstruct and check an unconditional proof of the curve point-separation
assertion A1 from the recovered later run. The requested Hooley route has a
separate geometric premise: we prove that its universal primitive-root theta
condition is false by a six-prime obstruction on an explicit genus-two curve.
We prove weaker prime-power and composite primitive-factor sufficient
conditions, including an all-genus characteristic-two geometric criterion,
with exact rank-one examples of degrees $13$ and $15$ and an
ordinary example of degree $4159$. A Lang--Weil test and its fixed-prime
limitation give a curve-dependent numerical route. Under M, a Hooley--Lenstra
calculation supplies a positive-density family of composite candidate orders.
Generic curves on both positive-rank strata avoid every nonzero odd torsion
order. For each fixed finite prime support, all geometric theta torsion is
confined to an explicit finite field; one rank-one example has no nonzero
three-primary theta point. No fixed finite prime support works uniformly over
finite fields, even allowing all exponents and products. A separate ordinary
example defeats the older large-prime quotient test over every constant
extension, despite having a theta witness. Frobenius-orbit and Weil-pairing
arguments give further exact restrictions on candidate primes and successful
numerical tests. A parity improvement gives a separate rational cover test,
and a rank-one curve defeats that size test over every extension. A
prime-power canonical-class criterion supplies prescribed-point rational
primitive factors without size inequalities, and strictly extends the
recovered prime-level detector on an actual curve. A genus-three example
refutes a proposed universal rule for detecting quotients. An exact
degree-seven geometric witness exists on the rank-one numerical-obstruction
curve even though its original rational points all miss the canonical detector.
The residual canonical
classes and the universal
positive-rank finite-field incidence or primitive-orbit coverage question
remains open.
Under the stipulated Alp\"oge theorem, Bach's applicable general unit-generator
constant is $3$, while exact unconditional coefficient-$2$ certificates are
provided for $53$ moduli and two infinite families. The Masullo draft claim
supplies specified Dedekind-GRH effective Chebotarev consequences, including
explicit full adic semisimplification tests. The audit covers every supplied
flag and all $800$ recovered full register rows, with $591$ available proof
blocks. The missing $390$ full register statements and $209$ further LM7
proof blocks prevent an exhaustive audit of the stated $1190$-row source set.
No proof through Hooley, formal verification of this paper, or verification
of Masullo's unfinished main theorem is claimed.
\end{abstract}
\tableofcontents
\clearpage
\part{Checked findings and precisely scoped cited consequences}
\section{Scope, analytic premises, and what has been established}
\label{sec:premises}

The principal point-separation assertion A1 is proved here without
number-field GRH, by a complete checked reconstruction of the later
ANABELIAN\_A1 argument. The originally proposed passage from Hooley's
primitive-root density to theta-divisor incidence is not a theorem:
A04 identifies the logical gap, and A06 gives an explicit geometric
counterexample to its universal primitive-root version. This changes
the starting diagnosis. It does not turn the counterexample into a
counterexample to A1.

The exhaustive audit remains in progress; the
requested proof through Hooley is also not established. A1 itself
is proved by the different route just specified. In source coverage,
800 of the 1,190 full register statements were recovered, and 591
statement/proof blocks are available. The 210-row
ANABELIAN\_A1\_SUCCESSOR and 180-row GT\_TOWER\_CHECK registers were not
recovered; their visible excerpts are audited separately. Availability
of a proof is not a claim that its whole content was checked. Each
actual check is recorded in \texttt{CHECK\_EVENTS.tsv}. The positive-rank
theta-incidence variant and unrelated open hypotheses in programme
rows remain recorded as open, rather than being assigned consequences
of GRH.

\subsection{The two inputs are deliberately different}

\paragraph{Input A: the stipulated Alp\"oge theorem.}
We accept the cited Dirichlet GRH theorem as instructed. The exact
supplied declarations in \texttt{ALPOGE\_Main.lean} are:
\begin{quote}\small\raggedright
\texttt{riemannHypothesis}\\
\texttt{riemannZeta\_eq\_zero\_re}\\
\texttt{LFunction\_ne\_zero\_of\_half\_lt\_re}\\
\texttt{LFunction\_eq\_zero\_re}\\
\texttt{LFunction\_eq\_zero\_re\_of\_isPrimitive}.
\end{quote}
The definitions in \texttt{ALPOGE\_Statements.lean} quantify over
Dirichlet characters of modulus $N$; the critical-strip assertion
has $0<\Re(s)<1$ and concludes $\Re(s)=1/2$.
The note's summary GRH claim is step~20, p.~4 \cite{Alpoge}.
No other intermediate step of the four-page note is needed in the
proofs below. In particular we do not convert its phrases about
purity into new geometric theorems.

This run inspected those declarations, definitions, and the attached
scan report. It did not build the Lean archive or rescan 175,187
source files. Acceptance as a cited theorem is the convention of
this task, not a newly reported verification of its formal proof.

\paragraph{Input M: the Masullo draft claim.}
Write M, or $\mathrm H_M$, for the draft's claimed GRH for the Hecke
$L$-functions of every number field, as described in \S1.2
\cite{Masullo}. Every use is explicitly conditional on that claim.
The trivial Hecke character over $F$ gives $\zeta_F$, so M includes
Dedekind GRH for every field used below.

\paragraph{Finding S01 (the exact numbered-source boundary).}
The supplied draft does not contain a completed numbered GRH
statement at its intended culminating locator. On printed/PDF
page~522, Chapter~13, Theorem~13.3.1.1 is a trace formula, referring
back to Theorem~12.5.1.1. The next heading is
Theorem~13.3.1.2, \emph{Divisor of $L(s,\chi)$}, but it has no
statement. Its preceding cross-reference is unresolved, and the
footnotes ask for insertion of the statement and proof of the main
theorems. Therefore this report cites M as the claim of \S1.2, with
13.3.1.2 as its unfinished intended location; it never cites the
latter as an already stated and proved GRH theorem.

\begin{proof}
The page was inspected both by direct PDF text extraction and by
rendering the page. The displayed material under 13.3.1.1 is the
trace formula. After the heading 13.3.1.2 the page has only the
editorial notes, and Chapter~14 starts on the following page.
This proves the source-status assertion. It neither supplies a
mathematical refutation of M nor claims a review of all 940 pages.
\end{proof}

\paragraph{Finding S02 (the analytic dependency rule).}
Input A discharges Dedekind GRH for finite abelian extensions of
$\mathbb Q$. M discharges it for arbitrary number fields. Neither
input, as used here, supplies a missing geometric existence theorem,
a missing finite-ramification hypothesis, or a description of an
unspecified finite quotient. No Artin holomorphy conjecture is needed
for the Chebotarev bounds chosen in this report.

\begin{proof}
For an abelian extension $F/\mathbb Q$ the Dedekind zeta function
is the product of the primitive Dirichlet $L$-functions corresponding to the
characters of $\operatorname{Gal}(F/\mathbb Q)$. At an unramified
prime with Frobenius order $f$, the character product of Euler
factors is $(1-p^{-fs})^{-[F:\mathbb Q]/f}$, exactly the Dedekind
factor; the inertia-invariant character factors give the ramified
version. Their nontrivial
zeros therefore satisfy the asserted line condition under A.
For arbitrary $F$, the trivial character specialization of M is
exactly $\zeta_F$. The remaining sentences specify hypotheses not
included in either premise; the concrete failures to make the
corresponding inferences are proved in A04, A06, T02 and G04.
The exact Chebotarev theorem G-CHEB assumes only the indicated
Dedekind zeta RH statements.
\end{proof}

There is a necessary fourth status in the audit: \emph{already
unconditional}. A qualitative Chebotarev theorem, a finite-field
Weil bound, an elementary modular computation, or a previously proved
geometric result does not become unconditional for the first time
because A is added. Calling such a result ``unconditional by Alp\"oge''
would misstate its dependency. The row tables retain the requested
A-only/M-only/beyond-both distinctions for actual analytic upgrades,
and this fourth status for existing statements.

\subsection{The geometric framework and exact locators}

All finite number-field monodromy quotients below are finite etale
covers of an open subscheme of $\operatorname{Spec}\mathbb Z$ (or
$\operatorname{Spec}\mathcal O_K$). The cyclotomic examples are the
covers obtained by adjoining roots of unity. The underlying
fibre-functor and generic-point correspondences are SGA~1,
Expos\'e~V, Th\'eor\`eme~4.1, \S7 and Propositions~8.1--8.2.
For a geometrically connected curve, its arithmetic fundamental-group
exact sequence and point sections are SGA~1, IX, Th\'eor\`eme~6.1
and Corollaire~6.4. Purely inseparable invariance is IX,
Th\'eor\`eme~4.10. These are used with the arithmetic Galois quotient
unchanged.

The configuration argument in A05 uses smooth proper curve lifting
(SGA~1, III, Th\'eor\`eme~7.3 and Corollaire~7.4), with marked
points then lifted by smoothness over the complete henselian base,
proper smooth
specialization (X, Th\'eor\`eme~3.8 and Corollaire~3.9), and the
prime-to-characteristic comparison and inertia of smooth proper pairs
(XIII, Corollaire~2.9, \S2.10, Lemme~2.11). Its surface relation is
XIII, Corollaire~2.12, and its tame Kummer extension at a boundary
uses XIII, Proposition~5.2 \cite{SGA1}. The exact archive context is
cote~149, scans~33--36 for Kummer separation and its covering
criterion, and scans~37--38 for the surrounding geometric reductions.

For compactified curves over finite fields, A05 uses the trace
formula and purity of $H^1$, with Frobenius eigenvalues of weight one:
Deligne, \emph{La conjecture de Weil. I}, formula~(1.5.1) and
Th\'eor\`eme~1.6 \cite{WeilI}. The equivalent action on the
Tate module and its dual must retain the full integral adic module,
not just its reduction modulo a prime. The curve cohomology and
orientation used there are also explicitly given in
\cite[Proposition~14.2 and Remark~14.10]{MilneLEC}.

The coherent Frobenius calculation in A03 and A06 concerns the
Artin--Schreier sequence (\cite[Example~7.9(b)]{MilneLEC}) and the stable part of $H^1(C,\mathcal O_C)$.
It is a characteristic-$p$ calculation of $p$-rank. No theorem about
complex zeros of number-field $L$-functions changes that rank.
Finite quotient representations in G01--G03 are weight-zero Artin
realizations; the original Tate modules may have positive weight.
No new purity theorem, motivic standard conjecture, crystalline
comparison, divided-power identity, or $\lambda$-ring theorem is
asserted. Those additional structures are not inputs to the proofs
in this report.


Findings A07--A31 develop the covering criterion of cote~149,
scans~35--36, inside the Picard variety and finite \'etale covers of
curves (SGA~1, V, Th\'eor\`eme~4.1 and IX, Th\'eor\`eme~6.1).
A08's finite-type spreading uses EGA~IV$_3$, Theorems~8.8.2 and
9.7.7, and EGA~III$_2$, Theorems~7.7.5 and 7.9.4.
A12--A14 and A19--A20 combine the finite \'etale Lang quotient,
coherent Frobenius and the weight-one Jacobian polynomial; the
point-count and absolute-value formulas are
\cite[Theorem~19.1(a)--(c)]{MilneAV} and
\cite[Theorem~11.1]{MilneJV}, with the cohomological formulation
in Weil~I, formula~(1.5.1) and Theorem~1.6.
A15--A18 use cyclic character factors of the Jacobian, coherent
Frobenius, and the characteristic-two \'etale Tate module.
A17 remains a number-field statement about Frobenius conditions
on Kummer covers of an open subscheme of $\operatorname{Spec}\mathbf Z$;
it does not turn those covers into a theta-incidence theorem.

A21--A22 concern the finite \'etale symplectic level torsors of the
relative Picard scheme and their arithmetic descent (SGA~1, V,
Th\'eor\`eme~4.1 and IX, Th\'eor\`eme~6.1). The finite-field
specialization uses effective affine descent and the finite-field
trace estimates specified in A22; it is unrelated to Dedekind GRH.
A23 uses the same weight-one curve and Jacobian Frobenius polynomials
as A13 to test the numerical covering criterion of cote~149,
scans~35--36, in the exact recovered form A1 B.3.

A24 combines that integral Frobenius action with the genus-two
intersection bound; its separate density calculation lies on the
Kummer and quadratic covers of an open subscheme of
$\operatorname{Spec}\mathbf Z$. A25 uses the polarized finite
\'etale torsion sheaf, its Galois-equivariant Weil pairing, and
the same weight-one point bounds. The precise pairing locator is
\cite[\S16 and Lemma~16.2(e)]{MilneAV}.

A26--A29 concern the same pointed finite etale coverings of cote~149,
scans~35--36. Their group-ring sums are character Euler factors of
unramified finite-order sheaves. Genus-two Riemann--Roch truncates their
reduction modulo the characteristic to a point-character sum; the
canonical class controls its inverse-character symmetry. A29 uses
the Picard functor, Lang isogeny and cyclic prime-power norm projector,
with the fibre-functor and arithmetic exact-sequence framework of
SGA~1, V, Theorem~4.1 and IX, Theorem~6.1. A28 is the explicitly
credited recovered successor A.26, not a new prime-degree theorem.

A30 tests the exact genus boundary of that character argument by a
ramified double map to a supersingular elliptic curve and local Euler
factors. Its proposed preliminary double cover is excluded from the
prime-to-characteristic fibre functor. A31 gives an exact principal
divisor, its Frobenius-twisted theta class, and the resulting tame cyclic
cover on the same curve as A27. The geometry remains within the
pointed-cover question of cote~149, scans~35--36.


\section{A1: the question, the older route, and the recovered proof}
\label{sec:a1}

\subsection{What A1 says}
The archive contains no statement labelled ``A1'': this is the programme's
label. Cote~149, scan~35, distinguishes the full prime-to-characteristic
fundamental group from its abelian quotient; scan~36 asks whether the
point-separation assertion survives in characteristic $p$, and whether its
covering condition forces $P=Q$. The exact displayed maps denoted $(*)$ and
$(**)$ occur before the supplied scans. We do not reconstruct their omitted
wording. The following is the curve statement of the runs, with all
quantifiers restored from the recovered full manuscript
\cite[A.1, A.31--A.33]{A1run}.

\paragraph{Question A1.}
Let $K$ be finitely generated over $\mathbb F_p$ and let $X/K$ be a smooth
geometrically connected hyperbolic curve. Put
\[
 \Delta_X=\pi_1(X_{\overline K})^{(p')},\qquad
 \Pi_X=\pi_1(X)/\ker\bigl(\pi_1(X_{\overline K})\to\Delta_X\bigr).
\]
The arithmetic quotient is the \emph{whole} $G_K$; only the geometric
kernel is replaced. Is
\[
 X(K)\longrightarrow
 \operatorname{Sec}(\Pi_X\to G_K)/\Delta_X,
 \qquad P\longmapsto[s_P]
 \tag{A1}
\]
injective? In characteristic zero the corresponding assertion uses every
geometric prime. Hyperbolic means $2g-2+r>0$ for the geometric smooth
compactification and its $r$ missing points.

For a proper curve of genus at least two over a finite field, the equivalent
covering formulation is this: for every finite constant extension $L/k$
and every connected geometrically prime-to-$p$ Galois cover $D\to C_L$
with a chosen $P'\in D(L)$ above $P$, there is $Q'\in D(L)$ above $Q$
such that $[Q'-P']\in J_D(L)[p^\infty]$. The term ``$p$-torsion'' in
scan~36 has this primary meaning, as its preceding phrase about order a
power of $p$ confirms. The phrase ``a cover of $X$'' in that scan must be
read as a cover of $Y$ in its local notation; see the transcription table.

\paragraph{Finding A01 (Kummer and the covering formulation).}
The preceding covering condition is equivalent to equality of point-section
classes over a finite field. Equality of sections always implies the
stated primary torsion condition. No number-field GRH occurs here.

\begin{proof}
For a finite field $L$, the prime-to-$p$ Kummer kernel of $J_D(L)$ is its
$p$-primary subgroup: finite Kummer maps inject $J_D(L)/nJ_D(L)$ into
$H^1(L,J_D[n])$, and $J_D(L)$ is finite. Equal sections give the same
action on the fibre functor, hence compatible rational lifts to any
pointed cover, with equal abelian Kummer classes.

For the converse we first establish uniqueness of the specified lift.
If two lifts differ by a nonidentity deck element $a$ of order $m$ prime
to $p$, then $p^e(aR-R)$ is principal. Choose a function with this
divisor. Its norm for $\langle a\rangle$ is constant; over the algebraic
closure scale it to norm one. Hilbert~90 writes it as $a(u)/u$.
Thus $\operatorname{div}(u)-p^eR$ is $a$-invariant. The free deck action
has every orbit of length $m$, so the degree $-p^e$ of this divisor is
divisible by $m$, a contradiction.

Take the cofinal characteristic open normal subgroups $N$ of the finitely
generated geometric group $\Delta_C$. The section $s_P$ makes
$N\rtimes s_P(G_k)$ a pointed cover. Make its geometric deck group
constant by a finite extension and apply the assumed condition.
Uniqueness descends the lift of $Q$ to $k$; norm maps and uniqueness
make these lifts compatible as $N$ varies. They form a fixed point of
the profinite path torsor from $P$ to $Q$, which is exactly a geometric
conjugacy of the two sections. The prime-to-$p$ topology on each cover
is the full one: normal cores embed the additional quotient in a finite
product of prime-to-$p$ groups. This also justifies the use of its
Jacobian Kummer map. This checks the form in \cite[A.31]{A1run}.
\end{proof}

\subsection{The exact contribution of Hooley}

\paragraph{Finding A02 (the analytic consequence, conditional on M).}
For every integer $a\ne-1$ which is not a square, write
$K_n(a)=\mathbb Q(\zeta_n,a^{1/n})$. Under M, Hooley's primitive-root
theorem applies, since the trivial Hecke character of $K_n(a)$ has
$L$-function $\zeta_{K_n(a)}$. In particular,
\[
 \#\{\ell\le x:\operatorname{ord}_{\ell}(2)=\ell-1\}
 = A\operatorname{Li}(x)
   +O\!\left(\frac{x\log\log x}{\log^2x}\right),\qquad
 A=\prod_{r\ {
m prime}}\left(1-\frac1{r(r-1)}\right)>0.
 \tag{A02}
\]
This is a cited consequence of \cite[\S7, concluding theorem]{Hooley}; the precise modern
numbered formulation used for the asymptotic and error is
\cite[Theorem~1.1, specialized to $a=2$]{FanPollack}. We do not re-prove
Hooley's sieve. The needed fields have squarefree $n$; already $K_3(2)$
has nonabelian Galois group $S_3$. The stipulated Alp\"oge Dirichlet
theorems therefore do not supply this application of Hooley.

\begin{proof}
Apply M to the trivial character over each $K_n(2)$, and then the
cited theorem. Its density series is
$\sum_{n\ge1}\mu(n)/[K_n(2):\mathbb Q]=A$. The polynomial $T^3-2$
is irreducible by Eisenstein and has nonsquare discriminant $-108$;
its splitting field is $K_3(2)$ and has group $S_3$. This identifies
a specific nonabelian analytic input absent from the Dirichlet scope.
It does not assert that Dirichlet GRH logically contradicts, or can
never imply by some different argument, the needed Dedekind GRH.
\end{proof}

\paragraph{Finding A03 (what the theta route actually needs).}
For a genus-two curve $C/\overline{\mathbb F}_2$, a sufficient geometric
input is the following. The theta divisor is defined by the bundle of
locally exact differentials on the Frobenius twist $C^{(1)}$; we
transport its geometric line bundles to $C$ by the absolute
Frobenius-twist identification, and denote the resulting divisor by
$\Theta_B$. The required condition is
\[
 \begin{gathered}
 \exists\ell>2\text{ prime},\qquad \operatorname{ord}_{\ell}(2)=\ell-1,\\
 \exists L\in\operatorname{Pic}^0(C)[\ell]\setminus\{0\}
                  \quad\text{with }L\in\Theta_B.
 \end{gathered}
 \tag{T(C)}
\]
Together with the necessary covering implication of A01, this implies
point separation for any arithmetic model of $C$. Hooley supplies the
first condition on $\ell$, not the joint assertion $T(C)$.

\begin{proof}
We spell out the geometric implication, so that the exact missing input
is visible. The line bundle $L$ defines a connected cyclic etale cover
$D\to C$ of degree $\ell$. Its nontrivial character spaces
$H^1(C,L^i)$, $1\le i<\ell$, are one-dimensional by Riemann--Roch.
Absolute Frobenius sends the $i$-space to the $2i$-space semilinearly.
The locally exact differential sequence
$0\to\mathcal O_{C^{(1)}}\to F_*\mathcal O_C\to B\to0$,
tensored with $L^i$, identifies its kernel with
$H^0(C^{(1)},B\otimes L^i)$, with the usual Frobenius-twist
identifications. A theta incidence therefore makes one arrow zero.
The primitive-root assumption makes all nontrivial characters one
cycle, so Frobenius on their sum is nilpotent. The Artin--Schreier
sequence identifies its stable part with etale $\mathbb F_2$
cohomology. Thus the new factor $\operatorname{im}(\sigma-1)$ of
$J_D$ has two-rank zero.

Suppose distinct points collide. After a finite extension A01 gives
lifts $P',Q'$ with $T=[Q'-P']$ two-primary. Since
$(\sigma-1)T$ lies in a two-rank-zero abelian variety, it is zero.
Hence $\sigma Q'+P'\sim Q'+\sigma P'$. Equality of the divisors
is impossible by freeness of the deck action and $P'\ne Q'$.
They therefore give a degree-two pencil on $D$. But $D$ cannot be
hyperelliptic: an odd-order free deck automorphism would induce a
nontrivial automorphism of $\mathbb P^1$, which has a fixed point;
its fibre has at most two geometric points, each fixed by the
odd-order element. This contradicts freeness. The implication follows.
The same proof is the prime-degree instance of the older intake's
A.C2.7 and A.C2.19, without requiring its general composite-conductor
criterion.
\end{proof}

\paragraph{Finding A04 (the density inference is false).}
Infinitude, or positive density, of theta-supporting primes would not
by itself combine with Hooley to prove $T(C)$. Indeed the primes
$\ell\equiv1\pmod8$ have density $1/4$, while both $2$ and $-1$
are squares modulo every such prime. Consequently
$\langle2,-1\rangle\ne\mathbb F_\ell^\times$ throughout that set.
This refutes the proposed inference between properties of sets of
primes; it is not a curve counterexample to A1.

\begin{proof}
Quadratic reciprocity gives $(2/\ell)=1$ and $(-1/\ell)=1$ in the
specified progression. Their subgroup lies in the proper subgroup
of squares. Dirichlet's prime-distribution theorem gives density
$1/\varphi(8)=1/4$. The same obstruction persists at $\ell^a$ by
reduction modulo $\ell$. A fixed finite Chebotarev condition can be
combined with primitive roots using \cite[Theorem~3.1, and
Theorems~4.1--4.2 for positivity]{Lenstra1977}; no such encoding,
or positivity for the curve-dependent torsion condition, is supplied
by M. Goldfeld's average theorem \cite[Theorem, p.~223]{Goldfeld}
averages the integer base, and does not supply it for the fixed base~2.
\end{proof}

There is a second limitation in higher genus. The spaces $H^1(C,L^i)$
then have dimension $g-1$. One singular arrow round a Frobenius orbit
need not make its composite nilpotent: take every arrow to be the
identity except one $\operatorname{diag}(0,1,\ldots,1)$.
Its stable subspace has dimension $g-2$. Thus replacing ``ordinary''
by ``not ordinary'' cannot supply zero primitive $p$-rank. The prime-to-$p$ cover statements in
\cite[Propositions~2.3--2.4]{PopSaidi} do not repair this
step. In the latter source ``pairwise prime orders'' is justified by
a proof of pairwise \emph{coprime} orders, not of prime integer orders.

\subsection{A1 does have a different proof in the later run}

\paragraph{Finding A05 (A1, independently of A and M).}
The assertion A1 holds, using the full arithmetic quotient. The
following proof is a checked reconstruction of
\cite[A.2--A.5, A.14--A.19, A.30--A.33]{A1run}, recovered from the
exact checkpoint identified by the supplied summary. It is not a
proof through Hooley, and no originality or formal verification is
claimed. Its decisive input is cuspidal transport with integral
Frobenius information, together with the already known finite
exceptional-set theorem of Sa\"idi--Tamagawa.

\subsubsection*{Step 1. The relative configuration fibre is retained}
For the finite-field assertion let $k$ be finite, let $C/k$ be proper
of genus at least two, and let $P,Q\in C(k)$. To prove the required
geometric group comparison first work with a complex proper curve
$C$ of genus $g\ge2$, and write
\[
 \Lambda=\pi_1^{\rm top}(C-P),\quad
 \Gamma=\pi_1^{\rm top}(C^2-\operatorname{diag}),\quad
 B_0=\pi_1^{\rm top}(C),\quad \Delta=B_0^{(p')}.
\]
The second projection is a locally trivial fibre bundle, obtained
locally by moving the missing point with an isotopy. Since
$\pi_2(C)=0$, it gives $1\to\Lambda\to\Gamma\to B_0\to1$.
Put $F=\Lambda^{(p')}$, $E=\Gamma^{(p')}$,
$K_F=\ker(F\to\Delta)$ and
$K_E=\ker(E\to\Delta\times\Delta)$. Fix $r\ne p$ and replace
each $K$ by its maximal pro-$r$ quotient. Denote the resulting
groups by $F_c,E_c$. We prove
\[
 1\longrightarrow F_c\longrightarrow E_c\longrightarrow\Delta
 \longrightarrow1. \tag{A05.1}
\]

We need a free-group lemma. If $1\to L\to H\to B\to1$ has $L$
free of finite rank at least two, and $H$ acts on
$L/[L,L]L^r$ through an $r$-group, then $L^{(r)}\to H^{(r)}$
is injective. For a finite $r$-group $Q$, the kernel of
$\operatorname{Aut}(Q)\to\operatorname{Aut}(Q/\Phi(Q))$ is an
$r$-group: it acts freely on the $|\Phi(Q)|^d$ ordered generating
lifts of a fixed basis. Characteristic Frattini quotients therefore
show that the closure of the action on $L^{(r)}$ is pro-$r$.
It extends to $H^{(r)}$ and restricts on $L^{(r)}$ to the inner
action. This action is faithful, because a nonabelian free pro-$r$
group has trivial centre. One direct verification uses its Magnus
embedding into $1+(X_1,\ldots,X_d)$ in the noncommutative formal
power-series algebra over $\mathbb F_r$. Injectivity follows by
substitution into each finite $r$-group algebra, whose augmentation
ideal is nilpotent. A series commuting with both $X_1$ and $X_2$
is scalar, so a central Magnus unit of augmentation one is one.
Compactness makes the injection a closed embedding; hence every
finite $r$-quotient of $L$ is dominated by one of $H$.

Now take any finite quotient $Q$ of $F_c$, with normal $r$-subgroup
$R$ the image of $K_F$, and put $A=Q/R$. The corresponding proper
Galois cover $C'\to C$ gives
$Y=(C'\times C')-(a\times a)^{-1}(\operatorname{diag})$.
The removed components are disjoint labelled graphs. Its fibre
$V'=C'-a^{-1}(P)$ has free fundamental group $L$. In
$H_1(V',\mathbb F_r)$, monodromy is trivial both on boundary
classes and on the quotient $H_1(C',\mathbb F_r)$, since the
compactified family is a product. It is therefore an $r$-group of
upper block unipotent transformations. The lemma supplies a finite
$r$-quotient of $H=\pi_1^{\rm top}(Y)$ dominating the restriction
of $Q$ to $L$. Taking its core under the finite quotient
$\Gamma/H=A\times A$ gives a quotient of $\Gamma$ which is an
$r$-group extended by $A\times A$, still prime to $p$.
Its image of $K_E$ is an $r$-group, so it factors through $E_c$.
On the fibre its kernel is contained in $\ker(F_c\to Q)$.
Every finite quotient of $F_c$ is thus detected in $E_c$, proving
injectivity. Right exactness follows from the universal property of
pro-prime-to-$p$ completion; the further killed subgroup lies in
the fibre kernel. This proves (A05.1).

Lift the smooth proper curve in characteristic $p$ to characteristic
zero (SGA~1, III, Corollary~7.4); its marked points lift by smoothness
over the complete henselian base. Apply the
prime-to-$p$ comparison for smooth proper pairs to the pointed
curve and to its square with diagonal boundary (SGA~1, XIII,
Corollary~2.9, \S2.10 and Lemma~2.11). Functoriality identifies
the groups, the relative kernels, and boundary inertia. Thus
(A05.1) holds for the etale groups. Arithmetic fibre comparison
over a finite field is an isomorphism on geometric kernels and
on the unchanged $G_k$, hence on the whole extensions.

\subsubsection*{Step 2. Equal proper sections transport the cusp}
In the arithmetic configuration quotient the diagonal boundary has
a common decomposition group $\mathcal D$ with
\[
 1\to\mathbb Z_r(1)\to\mathcal D\to\Pi_C\to1,
 \tag{A05.2}
\]
mapping diagonally to $\Pi_C\times_{G_k}\Pi_C$. Here is why no
boundary injection is assumed without proof. The boundary of the
$n$th root construction along the diagonal is the gerbe of roots
of its normal line bundle. Finite tame covers extend on its
Kummer charts by SGA~1, XIII, Proposition~5.2. The inverse-limit
boundary group therefore has kernel the image of tame inertia.
After the relative quotient, this is a quotient of $\mathbb Z_r(1)$.
It is the full group: the boundary word of the punctured fibre is
the product of handle commutators, and sending one handle pair
to $1+e_{12},1+e_{23}$ in $\operatorname{UT}_3(\mathbb Z_r)$
sends it to $1+e_{13}$ of infinite order. Step~1 retains this
inertia. The transverse fibre has multiplicity one, so its cusp
group is the full inverse image of $s_P(G_k)$ in $\mathcal D$.

If $s_P,s_Q$ are geometrically conjugate, lift their conjugator
$\gamma\in\Delta_C$ through $\mathcal D$. Its conjugation
identifies the two fibre groups and cusps, and acts on their first
proper coordinate by $\operatorname{Inn}(\gamma)$. A geometric
element of the target fibre mapping to $\gamma^{-1}$ cancels this
inner automorphism. Consequently we obtain
\[
 \beta:G_P\xrightarrow\sim G_Q
 \quad\text{over }\operatorname{id}_{\Pi_C},
 \tag{A05.3}
\]
where $G_P,G_Q$ are the maximal cuspidally pro-$r$ quotients of
the actual punctured curves. It preserves cusp conjugacy classes.

\subsubsection*{Step 3. Corresponding open covers have the same counts}
If $K\triangleleft F$ is closed and
$N=\ker(K\to K^{(r)})$, every open $I\le F$ containing $N$
satisfies $I^{(r)}\simeq(I/N)^{(r)}$. Indeed take a finite
$r$-quotient $I/J$, put $I_0=\operatorname{core}_F(I)$, and take
the core $M$ of $J\cap I_0$. Then $I_0/M$ is an $r$-group.
The image of $K$ in $F/I_0$ is also an $r$-group because
$N\subset I_0$. Thus its image in $F/M$ is an $r$-group, and
the defining property of $N$ gives $N\subset M\subset J$.

Apply this to a geometric open cover of $G_P$ or $G_Q$. Killing
its actual cusp inertias and taking its pro-$r$ quotient recovers
the full pro-$r$ group of its smooth compactification. This is
equivariant for arithmetic Frobenius. Thus $\beta$ identifies
the integral $r$-adic $H^1$ of corresponding compactifications.
The trace formula gives identical compactified point counts over
every finite extension; the cusp double cosets give identical
degrees of the missing points. Subtracting them gives identical
counts for the open covers. The inputs are the usual curve trace
formula and purity, not analytic RH over $\mathbb Q$.

\subsubsection*{Step 4. Compactness recognizes actual point sections}
Fix $R\in(C-P)(L)$ and let $S=\beta(D_R)$. For every open
subgroup $H$ of $(G_Q)_L$ containing $S$, the source cover for
$\beta^{-1}H$ has an $L$-point above $R$. Step~3 shows that the
corresponding target cover $X_H$ has an $L$-point. Its finite
sets $X_H(L)$ have nonempty inverse limit: any finite collection
of compatibility conditions is satisfied by a point on the cover
for the intersection of the corresponding subgroups. Compactness
then supplies a compatible family, and a lift of its base point
$T\in(C-Q)(L)$ to the full profinite fibre. Its decomposition
group lies in every $H$ (rationality at each level, and injectivity
of the decomposition group's projection to $G_L$), hence in
$\bigcap_{H\supset S}H=S$. Both project isomorphically to $G_L$,
so they are equal. Thus $\beta$ transports every point section
to a genuine point section.

Sa\"idi--Tamagawa \cite[Definitions~1.3, 1.5, 1.10 and
Proposition~1.8(vi)]{ST2009} give one finite exceptional image
$E_C\subset C^{\rm cl}$ for the cofinite prime set $p'$.
Outside it the proper decomposition group uniquely determines
the geometric point, even up to commensurability. Remark~1.2
makes this one exceptional set valid over every constant extension.
Since $\beta$ is over the identity of $\Pi_C$, the transported
point is $R$ whenever $R\notin E_C\cup\{P,Q\}$.

\subsubsection*{Step 5. The Kummer difference is impossible}
The proper abelianization gives $[P-Q]\in J_C(k)[p^\infty]$.
Choose $f\in k(C)^\times$ with
$\operatorname{div}(f)=p^e(P-Q)$; descent of a geometric
generator follows from Hilbert~90. Put $M=\mathbb Z_r(1)$.
The geometrically cuspidally central quotients $E_P,E_Q$ have
kernel $M$. Infinite inertia follows from the same Heisenberg
test, and their geometric extension classes are the degree-one
point classes in $H^2(\Delta_C,M)=\mathbb Z_r$.
The map induced by $\beta$ acts on $M$ by a unit $u$.
Since it fixes the proper quotient and the degree-one class,
$u=1$.

For completeness the orientation identification uses the
punctured surface presentation (SGA~1, XIII, Corollary~2.12).
Equivalently it is the Kummer first Chern class of
$\mathcal O_C(P)$, with residue $+1$. The group/etale
cohomology comparison follows from the tower of all geometric
prime-to-$p$ covers: cyclic torsors kill $H^1$ classes, covers
of degree divisible by $r^a$ kill $H^2(-,\mu_{r^a})$, and a
smooth proper curve has no higher such cohomology. Taking the
associated spectral sequence and then inverse limits gives
the asserted $\mathbb Z_r(1)$ identification.

On $E_P\times_{\Pi_C}E_Q$ the difference
$\chi(u,v)=\overline\beta(u)v^{-1}$ is an $M$-valued
crossed homomorphism. Pull it back to the group of
$C-\{P,Q\}$. Its residues are $(1,-1)$, so
$\kappa_r(f)-p^e\chi$ has zero residues and descends to
$H^1(\Pi_C,M)$. Inflation-restriction identifies this group
with $H^1(G_k,M)$: the Frobenius eigenvalues on
$\operatorname{Hom}(T_rJ_C,\mathbb Z_r(1))\otimes\mathbb Q_r$
have absolute value $\sqrt{\#k}$, hence no invariant, and its
integral module is torsion free. Finite-field Kummer theory
therefore gives $a\in k^\times$ with
\[
 \kappa_r(f)=p^e\chi+\kappa_r(a). \tag{A05.4}
\]

At a good closed point $R$ put $L=k(R)$. Step~4 makes
$\beta(D_R^P)$ conjugate to $D_R^Q$. Its conjugator projects
to the normalizer in $\Pi_{C_L}$ of the proper section $D_R$, which is $D_R$
itself by \cite[Lemma~1.7]{ST2009}. Multiplying by an element
of $D_R^Q$ puts it in the cuspidal kernel; in $E_Q$ it lies
in $M$. Thus $\chi$ evaluates to a coboundary, and (A05.4)
gives $\kappa_{r,L}(f(R)/a)=0$.

If $P\ne Q$, the function $f$ is nonconstant. The excluded
points and their $f$-values form a finite set. Choose a
nontrivial $r$-power root of unity $\zeta\in\overline k$
with $a\zeta$ outside that set, and a point $R$ above it.
Then $\zeta=f(R)/a\in k(R)^\times$. The $r$-adic Kummer
map of this finite cyclic group is injective on its $r$-primary
subgroup, contradicting its asserted zero value. This proves
injectivity over every finite field, including every constant
extension.

\subsubsection*{Step 6. Finitely generated fields and open curves}
For a finitely generated field $K/\mathbb F_p$, spread a proper
curve and two distinct colliding points to a smooth connected
finite-type base, shrinking so that the sections are disjoint.
At a closed point use its henselian local ring $R$ and fraction
field $L$. Smooth proper specialization (SGA~1, X,
Theorem~3.8 and Corollary~3.9) identifies the geometric
prime-to-$p$ groups. The arithmetic map $G_L\to G_{k(s)}$
is onto, since finite etale covers of the normal henselian
domain have connected generic fibres. The same conjugator
would therefore give colliding distinct points on the finite
residue-field curve, a contradiction.

We record also the Kummer fact needed for open curves. For a
semiabelian $G$ over $K$ the prime-to-$p$ Kummer kernel is
$G(K)[p^\infty]$; in characteristic zero, using all primes,
it is zero. After a finite splitting extension of its torus,
the divisor sequence on a normal proper model makes
$K^\times$ an extension of constant units by a free abelian
group. Constant units are finite over finite constants, and
finite torsion plus free abelian over number-field constants.
Lang--Neron (\cite[Theorem~2.1]{Conrad}, with the perfect constant
field algebraically closed in $K$), together with Mordell--Weil on the trace over
number fields, makes the abelian quotient's rational points
finitely generated. Lifting its finitely many generators and
relations shows that $G(K)$ is free abelian plus a finitely
generated group; the intersections of its prime-to-$p$
multiples and of all its multiples are respectively its
$p$-primary torsion and zero. This is the argument of
\cite[A.30]{A1run}; it does not falsely assume $K^\times$
finitely generated.

Make a finite algebraic extension to split the compactification
and boundary; purely inseparable base change leaves the etale
category and Galois group unchanged (SGA~1, IX,
Theorem~4.10). If the compactified genus is at least two,
map a collision to the proper curve. In genus zero, choose
three boundary points as $0,1,\infty$; their coordinate embeds
the curve into $\mathbb G_m$, and equal Kummer classes have
ratio of $p$-power order in $K^\times$, hence ratio one.
In genus one use a rational point as origin of the compact
elliptic curve $E$. A collision difference has order $p^a$.
Pull back the open curve under $[m]$, with $(m,p)=1$ and
$m^2>p^a$, and split its boundary and deck group by finite
extension. Lift the collision with first point zero. Its class
in the generalized Jacobian has exact order $p^a$: the torus
kernel has no nonzero geometric $p$-primary points. Thus a
function with divisor $p^a(Q'-P')$ equals one at all points
of the pulled-back reduced boundary. There are at least
$m^2$ of them, whereas $f-1$ has pole degree $p^a$, impossible.
In characteristic zero the all-prime Kummer kernel is zero
and the Abel map to the generalized Jacobian separates points.
This proves A1 with the stated scope. It does not assert the
full higher-dimensional or section-surjectivity conjecture.

\paragraph{Consequence for the two inputs.}
A05 needs neither A nor M. Therefore it remains true if only
Alp\"oge's stated theorems are added to the existing background.
Hooley's number-field application is still conditional on M;
the theta-support and primitive-cover questions are separate
geometric assertions. The older intake's unresolved A1 status
has been superseded by the recovered later proof, not by an
unproved implication from primitive-root density.


\subsection{An explicit obstruction to the proposed Hooley bridge}

\paragraph{Finding A06 (the universal primitive-root theta condition is false).}
The condition $T(C)$ of A03 does not hold for every genus-two curve in
characteristic two, even after arbitrary finite extension of its field of
definition. Take the smooth proper model of
\[
 \begin{split}
 k&=\mathbb F_2[t]/(M(t)),\quad
 M(Z)=Z^{38}+Z^{14}+Z^{10}+Z^9+Z^8+Z^5+Z^2+Z+1,\\
 d&=t^3+1,\qquad
 C:\ y^2+y=A(u)^3,\qquad
 A(u)=(u^2+t^2u+t)/d.
 \end{split} \tag{A06.1}
\]
For its locally exact differential theta divisor $\Theta$, there is no
odd prime $\ell$ with $\operatorname{ord}_\ell(2)=\ell-1$ and
$\Theta\cap J_C[\ell](\bar k)\ne\{0\}$. This holds geometrically;
number-field GRH cannot change it.

\begin{proof}
The curve and scalar Frobenius are those of \cite[B.23]{A1run}.
We recall the algebraic data used here and give the new small
primitive-root obstruction in full. Rabin's irreducibility test gives
$Z^{2^{38}}=Z$ modulo $M$ and
$\gcd(M,Z^{2^2}-Z)=\gcd(M,Z^{2^{19}}-Z)=1$.
Thus $k$ is the stated field and $d\ne0$. Reduction of the
Artin--Schreier polynomial by even powers leaves leading odd term
$(t^2/d^3)u^5$. There is a unique point $P$ at infinity, the genus
is two, and $K_C\sim2P$. The one-pole Artin--Schreier curve has
two-rank zero (equivalently, Cartier on the basis $du,u\,du$ has
stable rank zero). Hence $J_C[2](\bar k)=0$ and the unique geometric
theta characteristic is $\mathcal O(P)$. In characteristic two
the bundle of locally exact differentials is this theta characteristic,
so its theta divisor is $\Theta=C-P$. With relative Frobenius this
statement is on $C^{(1)}$; the same equations with coefficients squared
and the same scalar Frobenius give exactly the argument below.
In particular it is symmetric
and contains zero.

Put $B(u)=(tu^2+u+t^2)/d$. There are two maps to
$E:y^2+y=x^3$, given by $(A(u),y)$ and
$(B(u),y+(u^3+1)/d)$. Their equations follow from
\[
 (u^2+t^2u+t)^3+(tu^2+u+t^2)^3
 =d(u^3+1)^2+d^2(u^3+1).
\]
Their joint map is birational onto its image, since $u=B+tA$.
The image, whose normalization has genus two, cannot lie in an
elliptic translate in $E^2$; thus $J_C\to E^2$ is an isogeny.
Both maps take $P$ to zero. The duplication formula on $E$ is
$[2](x,y)=(x^4,y^4+1)$, so its fourth-power Frobenius is $[-2]$.
It follows that the $k$-Frobenius of $J_C$ is $[-2^{19}]$.
Indeed composing the difference with the isogeny gives zero,
and a homomorphism from a connected abelian variety into a finite
kernel is zero. These geometric identities and the irreducibility
calculation were separately recomputed in this run.

A nonzero difference in $J_C$ has at most two ordered representations
as a difference of points on $\Theta$. For if
$R_1-S_1=R_2-S_2\ne0$, then
$R_1+S_2\sim R_2+S_1$. Either the effective divisors coincide,
which gives the same representation, or they belong to the unique
degree-two hyperelliptic pencil. The latter determines the sole
alternative representation $(\iota S_1,\iota R_1)$.

Suppose a primitive-root prime $\ell$ and a nonzero
$L\in\Theta[\ell]$ existed. Since the theta divisor is defined over
$k$ and symmetric, it contains $hL$ for
$h\in H=\langle2^{19}\rangle\subset\mathbb F_\ell^\times$.
Write $e=\gcd(\ell-1,19)$ and $m=(\ell-1)/e$.
It contains these $m$ distinct points and zero. Counting their
ordered nonzero differences in the cyclic subgroup gives
\[
 m(m+1)\le2(\ell-1). \tag{A06.2}
\]
If $e=1$, this forces $\ell\le2$, impossible. If $e=19$,
it forces $m\le37$, hence $\ell\le704$ and
$\ell\equiv1\pmod{19}$. Exhaustive trial division of the eighteen
integers $19m+1$ with $m=2,4,\ldots,36$ gives exactly
\[
\begin{array}{c|rrrrrr}
 \ell&191&229&419&457&571&647\\ \hline
 \operatorname{ord}_\ell(2)&95&76&418&76&114&323.
\end{array}
\]
Only $419$ remains. Here $g=2^{19}\bmod419=119$, and $H$ contains
\[
 (g^0,g^{11})=(1,418),\quad
 (g^{19},g^{20})=(71,69),\quad
 (g^9,g^8)=(350,348).
\]
These are three distinct ordered pairs with the same nonzero
difference $2$ modulo $419$, contradicting the two-representation
bound. The included standard-library certificate
\texttt{check\_theta\_obstruction.py} verifies every finite calculation.
\end{proof}

\paragraph{Scope of the refutation.}
The stronger source statement \cite[B.23]{A1run} excludes every
nontrivial odd cyclic etale cover with zero primitive two-rank on this
curve. Its exhaustive conductor reduction and its complete polynomial
B\'ezout certificates at $191$ and $229$ were rerun successfully;
the verifier, data and logs are included. The shorter A06 proof does
not depend on those large certificates. In particular the weaker
arithmetic condition $\langle2,-1\rangle=\mathbb F_\ell^\times$
cannot repair the universal geometric assertion on this example,
by that checked B.23 result.

This curve already has two-rank zero, so ordinary Kummer theory
separates its points. It is not a counterexample to A1, nor a
refutation of a strategy which first removes all zero-rank curves
and then seeks a different theorem for the remaining curves.
The positive-two-rank theta-incidence problem has not been solved
by this argument. Its unresolved status does not affect A05.


% Follow-up supplement. Standalone finding IDs reserved by root: A07,A09,A10,A11.
% A08 is supplied separately by effective_galois/finite_prime_avoidance.tex.
% Bibliography additions are in POSITIVE_RANK_BIB.tex.

\section{Positive rank and a weaker geometric bridge}
\label{sec:positive-rank-followup}

\paragraph{Finding A07 (the generic positive-rank bridge fails).}
For each $f\in\{1,2\}$ there is a smooth proper genus-two curve over a
field finitely generated over $\mathbb F_2$, of geometric $2$-rank $f$,
whose Raynaud theta divisor has no nonzero point of any odd prime order.
Consequently the assertion $T(C)$ used in A03 is false on this larger
class of ground fields, even after imposing positive $2$-rank.

\begin{proof}
Choose an irreducible component $Z$ of the nonempty stratum
$\mathcal M_2^f$ and a curve $C_\Omega$ representing its geometric
generic point. Ozman--Pries, Theorem~1.1(2), equivalently
Theorem~8.3, applies with $p=2$, $g=2$, and either $f=1$ or $f=2$:
for every odd prime $\ell$, its theta divisor contains no point of
order $\ell$ \cite[pp.~456,473]{OzmanPries2019}.
These assertions apply to the same geometric generic curve for all
$\ell$; they do not require choosing a different curve for each prime.
The curve and a projective embedding have finitely many coefficients,
so they descend to a finitely generated subfield $K\subset\Omega$.
Smoothness and geometric connectedness descend under this field
extension. The $2$-rank and the absence of geometric theta torsion
are preserved by algebraically closed extension, so this curve over
$K$ has the claimed properties. Any required rational base point can
also be included among the finitely many coefficients.
\end{proof}

This conclusion does not give a counterexample over
$\overline{\mathbb F}_2$: the generic curve has transcendental moduli.
Finding A08 gives the separate finite-field consequence with its
correct quantifiers, namely avoidance of each prescribed finite set
of primes by a curve that may depend on that set.
For the ordinary generic curve, the stronger statement that all
prime-to-$2$ cyclic etale covers are ordinary is the Nakajima result
recorded as Proposition~4 in Yang's Theorem~1.1
\cite[Theorem~1.1]{Yang2021}. No restriction $p>2$ occurs there.
For a generic positive-rank stratum, Yang's Section~3.6.2 and
Corollary~3.6 explicitly retain prime degree; they do not supply a
prime-power version for rank one.

\paragraph{Source qualification.}
The original Nakajima paper is \emph{On generalized Hasse--Witt
invariants of an algebraic curve}, Adv. Stud. Pure Math.
\textbf{2} (1983), 69--88. Its original Proposition~4/Theorem~2 was
not recovered in this follow-up. The generic assertion above is
verified in Yang's paper and, for prime degree, directly in
Ozman--Pries. The latter's Theorem~8.2(3) is not used: its summary of
Pop--Saidi omits their finite-field hypothesis and describes prime
degrees, whereas the cited argument establishes pairwise coprime
orders, not prescribed prime orders. The finite-primary projection
argument in Pop--Saidi, Section~2, is the evidence used below.
\cite[Propositions~2.1,2.3 and their proofs, pp.~112--114]{PopSaidi}
In particular, Yang's stated all-cyclic-degree result also rules
out the sufficient prime-power incidence condition of A11 on a
generic ordinary curve. This stronger consequence is attributed to
the verified statement in Yang, not to an inspection of Nakajima's
unavailable original text. No corresponding all-prime-power claim
is made for the generic rank-one stratum.
There is also a direct modern theorem to cite: Yang's
Theorem~3.5 allows any integer $m$ prime to $p$ and says that a
cyclic cover of a geometric generic pointed curve is new-ordinary
exactly when the associated ramification divisors are Frobenius
stable. In the unpointed etale case every such divisor is zero,
which satisfies Definition~3.3 and Lemma~3.2 automatically.
This verifies the all-odd-degree conclusion at $p=2$ without
claiming access to Nakajima's original paper
\cite[Definition~3.3, Theorem~3.5]{Yang2021}.



\subsection{No uniform finite list of primes in positive 2-rank}
\label{sec:finite-prime-avoidance}

\paragraph{Published input (Ozman--Pries).}
Let $\ell\ne p$ be prime, let $g\ge2$, and let $0\le f\le g$, with
$f\ne0$ when $g=2$.  For each irreducible component $Z$ of
$\mathcal M_g^f$, the Raynaud theta divisor of the geometric generic
curve of $Z$ contains no point of order $\ell$.  This is
Ozman--Pries, \emph{Ordinary and almost ordinary Prym varieties},
Asian J. Math. \textbf{23} (2019), 455--478,
Theorem~1.1(2), p.~456, equivalently Theorem~8.3, p.~473 \cite{OzmanPries2019}.
In particular its hypotheses allow $p=2$, $g=2$, and either $f=1$ or
$f=2$.  The assertion concerns every geometric point of order $\ell$,
not just a selected cyclic cover.

\paragraph{Finding A08 (finite prime avoidance).}
For either $f\in\{1,2\}$ and every finite set $S$ of odd primes, there
exist an integer $m\ge1$ and a smooth projective geometrically connected
curve $C/\mathbf F_{2^m}$ of genus two and $2$-rank $f$ such that
\[
 \Theta_C(\overline{\mathbf F}_2)\cap
 \bigl(J_{C^{(1)}}[\ell](\overline{\mathbf F}_2)\setminus\{0\}\bigr)
 =\varnothing\qquad(\ell\in S).
 \tag{FP.1}
\]
Here $C^{(1)}$ is the relative Frobenius twist, and $\Theta_C$ is the
theta support of the sheaf of locally exact differentials on $C^{(1)}$.
This is the convention of Ozman--Pries, Section~8.1, p.~472.
One can equivalently use the usual Frobenius-semilinear relabelling of
this support; the existence and order of its geometric torsion points
are unchanged.  For $f=2$ the curve is ordinary.

\paragraph{Recovered source scope.}
The ordinary finite-list obstruction already appears in recovered
Successor finding A.16 \cite[A.16]{ASrun}, which allows arbitrary
characteristic and genus.  Its full recovered proof uses generic
ordinary Pryms and finite cover moduli.  The argument below states the
geometric theta-support conclusion explicitly, gives a scheme-family
spreading proof without coarse-moduli descent, and also treats the
fixed $2$-rank-one stratum.

\begin{proof}
\emph{1. A curve at a geometric generic point.}
The strata $\mathcal M_2^1$ and $\mathcal M_2^2$ are nonempty.  For
example, the smooth projective models of
\[
 y^2+y=x^3+x^{-1},\qquad
 y^2+y=x+x^{-1}+(x-1)^{-1}
\]
have genus two and $2$-ranks one and two, respectively.  For an
Artin--Schreier equation with $r+1$ poles of odd orders $d_j$, the
genus is $\bigl(\sum_j(d_j+1)-2\bigr)/2$ and the $2$-rank is $r$;
see Pries--Zhu, \emph{The $p$-rank stratification of Artin--Schreier
curves}, Ann. Inst. Fourier \textbf{62} (2012), 707--726,
Lemma~2.6, p.~712 \cite{PriesZhu2012}.

Fix an irreducible component of $\mathcal M_2^f$ over
$\overline{\mathbf F}_2$ and a curve $C_\Omega$ at its geometric generic
point, where $\Omega$ is an algebraically closed field of characteristic
two.  Apply the published input to each $\ell\in S$, on this same
component and this same geometric generic curve.  Thus $C_\Omega$
satisfies all the required avoidance conditions simultaneously.

\emph{2. Spreading the curve while preserving its $2$-rank.}
Choose a point $P\in C_\Omega(\Omega)$ and a projective embedding.
The equations and the point involve finitely many coefficients.
They therefore descend to an integral finitely generated
$\mathbf F_2$-subalgebra $A\subset\Omega$; the precise limit reference
for the scheme and the marked-point morphism is EGA~IV$_3$,
Theorem~8.8.2(ii) and (i), printed p.~28, respectively;
see also the Stacks Project, Lemma~32.10.1(1)--(3),
Tag~\texttt{01ZM} \cite{EGAIV3,Stacks}.
After replacing
$T=\operatorname{Spec}A$ by a nonempty affine open, they give a smooth
projective family $\pi:\mathcal C\to T$ of geometrically connected
genus-two curves with a section, whose base change to $\Omega$ is
$(C_\Omega,P)$.  This uses the usual finite-presentation spreading
argument: the projective model is proper, the nonsmooth locus has
closed image missing the generic point, and we remove that image.
The geometric-connectedness locus is constructible by EGA~IV$_3$,
Theorem~9.7.7(ii), p.~79; since it contains the generic point of
the integral base, it contains a nonempty open.  The genus is
preserved by constancy of the Euler characteristic in this smooth
proper family, EGA~III$_2$, Theorem~7.9.4, pp.~76--77.
There is no need to choose a universal curve over a coarse moduli
space or to descend a generic curve to that coarse space's function
field.

We may also ensure that every geometric fibre has $2$-rank $f$.
For completeness, locally trivialize the rank-two bundle
$E=R^1\pi_*\mathcal O_{\mathcal C}$ and let $H$ be the matrix of its
Frobenius-semilinear endomorphism.  On a geometric fibre, the stable
rank of Frobenius is the $2$-rank, and in dimension two the image has
stabilized by the second iterate.  Hence the $2$-rank is
$\operatorname{rank}(H H^{(2)})$, where $H^{(2)}$ denotes entrywise
squaring.  All minors of size $f+1$ vanish at the generic point and
thus vanish identically on the integral base; some minor of size $f$
is nonzero at the generic point.  Removing its zero locus gives the
assertion.  Equivalently one can use the locally closed $p$-rank
strata recalled in Ozman--Pries, Section~2.5, p.~458.

\emph{3. The relative theta incidence is closed.}
Put $\mathcal C^{(1)}=\mathcal C\times_{T,F_T}T$, and let
$F:\mathcal C\to\mathcal C^{(1)}$ be relative Frobenius.
For a smooth relative curve in characteristic two it is finite locally
free of degree two, and
\[
 0\longrightarrow\mathcal O_{\mathcal C^{(1)}}
 \longrightarrow F_*\mathcal O_{\mathcal C}
 \longrightarrow\mathcal B\longrightarrow0
 \tag{FP.2}
\]
has a line-bundle quotient $\mathcal B$, compatible with base change.
These claims can be checked in a relative etale coordinate, where
$F_*\mathcal O$ has basis $1,t$ over the Frobenius-twisted coordinate
ring.  On each fibre, $\mathcal B$ is the sheaf used in the definition
of the Raynaud theta divisor.

Let $\mathcal J=\operatorname{Pic}^0(\mathcal C^{(1)}/T)$, an abelian
scheme.  The chosen section gives a Poincare line bundle $\mathcal P$
on $\mathcal C^{(1)}\times_T\mathcal J$, normalized along that section.
Upper semicontinuity of cohomology for its proper projection to
$\mathcal J$ shows that
\[
 Z_\Theta=
 \{z\in\mathcal J:
 h^0((\mathcal C^{(1)}\times_T\mathcal J)_{\overline z},
       (\operatorname{pr}_1^*\mathcal B\otimes\mathcal P)_{\overline z})>0\}
 \tag{FP.3}
\]
is a closed subset, where $\overline z$ denotes a geometric point
over $z$.  Its geometric fibres over $T$ are
exactly the supports of the relevant theta divisors.  Thus the
argument requires only this closed incidence locus, and does not
require a separate global assertion that a chosen relative theta
divisor is Cartier.  The precise semicontinuity theorem is
Grothendieck, EGA~III$_2$, Theorem~7.7.5(I), printed p.~67;
equivalently, the Stacks Project, Lemma~36.32.1, Tag~\texttt{0BDN} \cite{EGAIII2,Stacks}.
Its hypotheses hold because the base is noetherian, the projection
is proper and smooth, and
$\operatorname{pr}_1^*\mathcal B\otimes\mathcal P$ is a line bundle,
hence coherent and flat over $\mathcal J$.

\emph{4. All nonzero torsion points are removed at once.}
For each odd prime $\ell$, the group scheme $\mathcal J[\ell]$ is
finite etale over $T$.  Its zero section is both open and closed, so
\[
 \mathcal J[\ell]^\times:=\mathcal J[\ell]\setminus0_T
\]
is still finite etale over $T$.  The intersection
$Z_\ell=Z_\Theta\cap\mathcal J[\ell]^\times$ is closed in this finite
scheme.  Consequently its image $D_\ell\subset T$ is closed, because
finite morphisms are closed.  A geometric fibre over $t$ contains a
nonzero $\ell$-torsion point on its theta support if and only if
$t\in D_\ell$.

By Step~1, $D_\ell$ misses the generic point of $T$ for every
$\ell\in S$: otherwise its nonempty generic fibre would remain
nonempty after the field extension to $\Omega$.  Therefore
\[
 U=T\setminus\bigcup_{\ell\in S}D_\ell
\]
is a nonempty open subscheme, containing the generic point.
It is of finite type over $\mathbf F_2$, so it has a closed point $u$
with finite residue field $\kappa(u)=\mathbf F_{2^m}$, by the
Nullstellensatz.  The curve $C=\mathcal C_u$ has genus two and
$2$-rank $f$, and satisfies (FP.1) over an algebraic closure of its
residue field.  This proves the claim.
\end{proof}

\paragraph{Ordinary strengthening of A08 (any finite list of odd orders).}
For $f=2$, the finite set $S$ may be replaced by any finite set
$\mathcal N$ of odd integers $n\ge3$: one can require
\[
 \Theta_C(\overline{\mathbf F}_2)\cap
 \bigl(J_{C^{(1)}}[n](\overline{\mathbf F}_2)\setminus\{0\}\bigr)
 =\varnothing\qquad(n\in\mathcal N).
 \tag{FP.4}
\]
In particular there is no uniform finite list of odd prime-power
orders for the sufficient incidence condition in A11.

\begin{proof}
Yang, \emph{Raynaud--Tamagawa theta divisors and new-ordinariness of
ramified coverings of curves}, J. Algebra \textbf{587} (2021),
263--294, Theorem~3.5 \cite{Yang2021}, applies to arbitrary cyclic degree $m$
prime to the characteristic over a generic pointed stable curve.
For an unpointed smooth curve every ramification divisor $D(j)$
is zero, which is Frobenius stable by Definition~3.3 and
Lemma~3.2.  Hence every connected cyclic etale cover of the
generic genus-two curve, of any odd degree, has ordinary new
factor.  The generic curve itself is ordinary, as also recorded
in Yang's Theorem~1.1.  These statements concern this same
geometric generic curve.

To verify that this excludes odd-order theta points, suppose $L$
has exact odd order $d>1$ on its Raynaud theta support, viewed on
the Frobenius twist.  Choose $e$ with $2e\equiv1\pmod d$ and set
$M=(F^*L)^e$ on the original curve.  Then $M$ has order $d$,
$M^2=F^*L$, and $M^{(1)}=L$; the last equality follows by applying
$F^*$, which is injective on odd-order torsion.
The twisted Frobenius sequence
\[
 0\longrightarrow L\longrightarrow F_*F^*L
 \longrightarrow B\otimes L\longrightarrow0
\]
therefore identifies $H^0(B\otimes L)$ with the kernel of the
semilinear Frobenius arrow
$H^1(C,M)\to H^1(C,M^2)$.
The two spaces are lines, since $C$ has genus two and the line
bundles are nontrivial of degree zero.  This arrow is zero.
In the connected degree-$d$ Kummer cover defined by $M$, it is
one arrow in the nontrivial-character sum of $H^1(\mathcal O)$.
That sum is Frobenius stable, and its Frobenius is not injective,
so its stable rank is smaller than its dimension.  Its new
factor is therefore not ordinary, contradicting Yang's theorem.
Thus the geometric generic theta support contains no nonzero
point of any odd order.

Now use this generic ordinary curve in Steps~2--4 above.  For each
$n\in\mathcal N$, the scheme $\mathcal J[n]\setminus0_T$ is finite
etale.  Its theta incidence has closed image missing the generic
point, because every nonzero element has odd order dividing $n$.
Remove the finitely many images and choose a closed point of the
remaining open.  This gives (FP.4).
\end{proof}

This stronger odd-order assertion is stated only for the ordinary
case $f=2$.  Yang's Section~3.6.2 and Corollary~3.6 retain a prime
degree hypothesis on a generic fixed lower-$p$-rank stratum; no
prime-power extension for $f=1$ is inferred here.

\paragraph{The precise nonuniformity conclusion.}
Taking $S$ to be all odd primes at most an arbitrary real number $B$
shows that no bound depending only on characteristic two, genus two
and the fixed positive $2$-rank can guarantee a nonzero prime-order
theta point below that bound for every curve over finite extensions
of $\mathbf F_2$.  More generally, no fixed finite list of candidate
odd primes works for all those curves.  The extension degree $m$ is
allowed to depend on the list; this proof gives no numerical bound
for $m$ or for the defining equations.

This statement has quantifiers
\[
 \text{for every finite }S\quad\text{there exists }C_S
 \quad\text{avoiding every }\ell\in S.
\]
It does not assert that one finite-field curve avoids every odd prime,
and it does not obstruct a curve-dependent prime of unbounded size.
An infinite intersection of the open avoidance loci need not contain
a closed point.  No analytic hypothesis is used in this deduction.


\paragraph{Finding A09 (an explicit rank-one theta point of order thirteen).}
Let
\[
 C/\mathbb F_2:\quad Y^2+X^2Y=X^5+X,\qquad P=\infty.
\]
This curve has genus two and $2$-rank one. Put
\[
 k=\mathbb F_2[u,v]/(u^2+u+1,\ v^2+u^2v+1)\simeq\mathbb F_{16},
 \qquad R=(u,v)\in C(k).
\]
Then $L=\mathcal O_C(R-P)$ has exact order $13$ and lies on the
Raynaud theta divisor. Since $2$ has order $12$ modulo $13$,
this is a positive-$2$-rank instance of $T(C)$.

\begin{proof}
The substitution $x=X^{-1}$, $y=Y/X^2$ gives
$y^2+y=x^3+x^{-1}$. Its two Artin--Schreier poles have odd orders
$3$ and $1$. Riemann--Hurwitz gives $2g-2=-4+4+2=2$, and
Deuring--Shafarevich gives $f_C-1=2(-1)+2=0$, so $g=2$ and $f_C=1$.
There is one point $P$ above $X=\infty$, with
$\operatorname{ord}_P(X)=-2$ and $\operatorname{ord}_P(Y)=-5$.
The differential $dX/X^2=d(1/X)$ has divisor $2P$.
It is locally exact, and hence its corresponding section of the
rank-one sheaf $B^1$ has divisor $P$. Thus
$B^1\simeq\mathcal O_C(P)$ and the theta support is $C-P$.
Here the Frobenius twist is identified with the same equation over
$\mathbb F_2$; no choice of a twist over a nonperfect field is used.

The field presentation is explicit. Over $\mathbb F_4=\mathbb F_2(u)$,
putting $w=v/u^2$ gives $w^2+w=u^2$. Its trace to $\mathbb F_2$ is
$u^2+u=1$, so the second quadratic is irreducible. The relation
$u^3=1$ gives $u^5+u=1$, proving $R\in C(k)$.

Define
\begin{align*}
 A(X)={}&u(v+1)(X^6+X^3)+u^2vX^5+uX^4
                 +(uv+1)X^2+uX+u^2,\\
 B(X)={}&X^4+uX^3+X+1,\qquad F=A(X)+B(X)Y.
\end{align*}
Direct polynomial reduction in the displayed field gives
\[
 A(X)^2+X^2A(X)B(X)+(X^5+X)B(X)^2=(X+u)^{13},
 \tag{PR.1}
\]
and
\[
 F(R)=0,\qquad F(\iota R)=u\ne0,
 \qquad\iota(X,Y)=(X,Y+X^2).
 \tag{PR.2}
\]
The left side of (PR.1) is the norm $F\,\iota(F)$.
As $u\ne0$, the two points over $X=u$ are unramified and distinct.
Equations (PR.1)--(PR.2) imply that every finite zero of $F$ is $R$
and that its multiplicity is exactly $13$. There are no finite poles.
At $P$, the term $X^4Y$ has pole order $13$, while every other term
has smaller pole order. Therefore
\[
 \operatorname{div}(F)=13R-13P.\tag{PR.3}
\]
The class $R-P$ is nonzero: otherwise a rational function with that
divisor would give a degree-one map from a genus-two curve to
$\mathbb P^1$. Since $13$ is prime, (PR.3) proves its exact order.
Also $B^1\otimes L\simeq\mathcal O_C(R)$, so $L\in\Theta_{B^1}$.
Finally, the successive powers of $2$ modulo $13$ are
\[
 1,2,4,8,3,6,12,11,9,5,10,7,
\]
which are all the nonzero residues.
\end{proof}

\paragraph{An explicit cover and the arithmetic check.}
The normalization of the Kummer equation
\[
 W^{13}=F
\]
is a geometrically connected etale $\mu_{13}$-torsor over $C_k$.
Indeed, (PR.3) makes every valuation divisible by $13$, so it is
etale; it is geometrically connected because $F$ cannot be a
thirteenth power, again by nontriviality of $R-P$. With the local
parameter $t=X^2/Y$ at $P$, the unit $t^{13}F$ has residue $1$.
The root $tW=1$ therefore gives a $k$-rational point above $P$.
After extending constants to $\mathbb F_{4096}$, which contains
$k$ and $\mu_{13}$, its deck group is the constant group
$\mathbb Z/13$. Its genus is $14$ and its new Jacobian factor has
$2$-rank zero by A11 below (or by the prime case A03).
For this assertion about the displayed Kummer model, note that
$C$ and its theta support are defined over $\mathbb F_2$, so
Frobenius relabelling preserves the theta incidence; the explicit
twist normalization in A11 applies to its defining torsion class.
The direct torsor construction here is not asserted to be the same
arithmetic model as the Lang cover over $\mathbb F_{64}$ in the
recovered A1 row B.39, working identifier B.LIT7.

The file \texttt{theta13\_direct.py} constructs the identity using
Hasse--Taylor linear algebra over
$\mathbb F_2[z]/(z^4+z+1)$, with $u=z^2+z$, $v=z^3$.
Its certificate lists every coefficient. The separate verifier
\texttt{verify\_theta13\_norm.py} imports no producer code and
checks (PR.1)--(PR.2) by polynomial reduction over $\mathbb F_2$.
These are exact finite computations; they do not extrapolate to a
classification of positive-rank curves.

\paragraph{Nonabelian content of the example.}
The same curve has
\[
 \#C(\mathbb F_2)=4,\qquad \#C(\mathbb F_4)=4,\qquad
 H_2(U)=U^4+U^3+2U+4,\qquad \#J_C(\mathbb F_2)=8.
 \tag{PR.10}
\]
Indeed, the Artin--Schreier model has its two rational branch
points, and $x=1$ contributes two points. For
$x\in\mathbb F_4\setminus\mathbb F_2$, its right side is
$x^3+x^{-1}=1+x^2=x$,
whose absolute trace is $1$, so no additional points occur.
The two counts and the genus-two functional equation give the
displayed characteristic polynomial; evaluating it at $1$ gives
$8$. Thus all four rational points have the same prime-to-$2$
abelian Kummer class, because their divisor differences lie in
the $2$-group $J_C(\mathbb F_2)$. Nevertheless A09 and A11 show
that their full geometrically prime-to-$2$ point-section classes
are pairwise distinct. This is a concrete positive-rank instance
where the nonabelian covering argument supplies information
absent from the abelian quotient.

\paragraph{Finding A10 (three is excluded for every nonordinary genus-two base).}
Let $C$ be a smooth proper genus-two curve over an algebraically
closed field of characteristic two, with $f_C<2$. Then its Raynaud
theta divisor contains no nonzero point of order $3$.
Consequently every connected cyclic etale degree-three cover has
ordinary new factor of dimension two. In particular, this conclusion
also holds when $f_C=1$.

\begin{proof}
Work on the Frobenius twist, temporarily denoting that curve by
$C$ as well, and write $B=B^1$ there. In
characteristic two this is a line bundle with $B^{\otimes2}\simeq
\omega_C$ and degree one. The square relation follows in an etale
coordinate by comparing transition functions of locally exact
differentials with those of $\omega_C$.
From the relative Frobenius exact sequence,
$H^0(B)$ is the kernel of Frobenius on $H^1(\mathcal O_C)$.
Nonordinarity implies that Frobenius is not bijective, so $H^0(B)\ne0$.
Thus $B\simeq\mathcal O(W)$ for a point $W$, and
$2W\sim K_C$. Its theta support is $C-W$.

If $L\ne0$ had order $3$ on this support, write
$L=\mathcal O(R-W)$ with $R\ne W$. Then $3(R-W)$ would be
principal, so there would be a function with unique pole of exact
order $3$ at $W$. But Riemann--Roch and $2W\sim K_C$ give
\[
 h^0(2W)=2=h^0(3W).
\]
Every function with pole divisor at most $3W$ therefore already has
pole divisor at most $2W$, a contradiction.
For a cyclic degree-three cover, every nontrivial eigenspace
$H^1(C,L^i)$ is one-dimensional. The Frobenius maps between the
two eigenspaces have zero kernel precisely when the associated
theta incidences are absent, by the same exact sequence used in
A11. They are both bijective, so the new factor is ordinary.
\end{proof}

\subsection{A prime-power bridge and the remaining order obstruction}
\label{sec:prime-power-bridge}

\paragraph{Finding A11 (a sufficient prime-power incidence condition).}
Let $C$ be a smooth proper genus-two curve in characteristic two.
Let $n=\ell^r$ with $\ell$ an odd prime and $r\ge1$. Suppose that
\[
 \langle 2,-1\rangle=(\mathbb Z/n\mathbb Z)^\times
 \tag{PR.4}
\]
and that the Raynaud theta divisor contains a class of exact order
$n$. Then over an algebraic closure there is a connected cyclic etale
degree-$n$ cover $D_n\to C$ whose primitive conductor-$n$ factor has $2$-rank zero.
Writing $D_{n/\ell}$ for its intermediate degree-$n/\ell$ cover of
$C$, this factor is precisely
\[
 \operatorname{Prym}(D_n/D_{n/\ell})
\]
up to isogeny. In the finite-field covering formulation of A1, this
single cover separates every proposed pair of distinct colliding
points. Thus prime-power incidence with (PR.4) is sufficient; prime
order and the stronger equality $\langle2\rangle=(\mathbb Z/n)^\times$
are not required.

\begin{proof}
\emph{1. Symmetry and a zero Frobenius map.}
Let $B$ be the locally exact differential bundle on $C^{(1)}$ and
$L\in\operatorname{Pic}^0(C^{(1)})$ the given class. Since
$B^{\otimes2}\simeq\omega_{C^{(1)}}$ and $\chi(B\otimes L)=0$,
Serre duality gives
\[
 h^0(B\otimes L)=h^0(B\otimes L^{-1}).
 \tag{PR.5}
\]
Thus both $L$ and $L^{-1}$ lie on theta. Tensoring the relative
Frobenius sequence by $L$ gives
\[
 0\longrightarrow L\longrightarrow F_*F^*L
       \longrightarrow B\otimes L\longrightarrow0.
 \tag{PR.6}
\]
The two degree-zero line bundles $L$ and $F^*L$ are nontrivial:
$n$ is odd and Frobenius pullback preserves its exact order.
They have no sections, and each has one-dimensional $H^1$ by
Riemann--Roch. Therefore $H^0(B\otimes L)\ne0$ says that the
induced map from $H^1(L)$ to $H^1(F^*L)$ has nonzero kernel;
being a map of lines, it is zero. The same applies to $L^{-1}$.
To make the character labels precise, choose $e$ with $2e\equiv1
\pmod n$ and put $M=(F^*L)^e$ on $C$. This class still has exact
order $n$, and $M^2=F^*L$. If $M^{(1)}$ denotes its Frobenius
twist on $C^{(1)}$, then $F^*(M^{(1)})=M^2=F^*L$.
The pullback $F^*$ is injective on odd-order torsion, so
$M^{(1)}=L$. Consequently (PR.6) gives literally the
Frobenius-semilinear arrow from $H^1(C,M)$ to $H^1(C,M^2)$;
the inverse class gives the arrow from $H^1(C,M^{-1})$ to
$H^1(C,M^{-2})$.

\emph{2. Every primitive orbit is nilpotent.}
Over an algebraic closure, the Kummer correspondence gives the
connected cyclic etale cover associated to $M$, choosing its
character labels so that the eigensheaves are $M^i$.
When $C$ is over a finite field, all these data descend after a
finite constant extension. Its cohomology decomposes as
\[
 H^1(D_n,\mathcal O_{D_n})=
   H^1(C,\mathcal O_C)\oplus
   \bigoplus_{i=1}^{n-1} V_i,\qquad
   V_i=H^1(C,M^i),\quad \dim V_i=1.
 \tag{PR.7}
\]
Absolute Frobenius sends $V_i$ to $V_{2i}$. The primitive part is
the sum over $i\in(\mathbb Z/n)^\times$.
Condition (PR.4) means that the action of multiplication by $2$
on this unit group has at most two orbits, represented by $1$ and
$-1$. By Step~1, each orbit contains a zero arrow. A sufficiently
large iterate of Frobenius is therefore zero on every $V_i$ in
the primitive sum: traverse the finite orbit until the zero arrow
is reached. The stable Frobenius part of that sum is zero, which
is the assertion that the primitive factor has $2$-rank zero.
The one-dimensional spaces are those on the original genus-two
base; no inference from a single singular higher-dimensional
matrix is being made.

\emph{3. Prime powers identify the entire relative factor.}
Choose a generator $\sigma$ of the order-$n$ deck group and put
$a=\sigma^{n/\ell}$, of order $\ell$.
The quotient by $\langle a\rangle$ is $D_{n/\ell}$.
On the character with index $i$, the element $a$ acts trivially
if and only if $\ell\mid i$. For $n=\ell^r$, its nontrivial
characters are exactly the primitive indices $\gcd(i,n)=1$.
Thus the primitive factor is exactly the new factor for this
degree-$\ell$ quotient, proving the stated Prym identification.
The dimension is $(n-n/\ell)(g(C)-1)=\varphi(n)$, as a check.

\emph{4. The point-separation consequence.}
Now work in the finite-field situation of the covering formulation
of A1. Enlarge constants to define the cover, make its deck group
constant, and choose a rational point $P_n$ above the given base
point $P$. If a distinct point $Q$ collided with $P$, the covering
condition would provide a rational $Q_n$ above $Q$ with
\[
 \delta=[Q_n-P_n]\in J_{D_n}[2^\infty].
\]
The point $(a-1)\delta$ is still $2$-primary and belongs to
$\operatorname{im}(a-1)=\ker(\operatorname{Nm})^0
=\operatorname{Prym}(D_n/D_{n/\ell})$.
An abelian variety of $2$-rank zero has no nonzero geometric
$2$-power torsion points. Thus $(a-1)\delta=0$.
Equivalently, the divisor
\[
 aQ_n+P_n-Q_n-aP_n
 \tag{PR.8}
\]
is principal. Its positive and negative parts both have degree two
and are disjoint: the deck action is free, and the orbits above
the distinct base points $P$ and $Q$ are disjoint. It therefore
defines a degree-two map $D_n\to\mathbb P^1$.

Such a map is separable, since a purely inseparable map of degree
two from a smooth curve over an algebraically closed field to
$\mathbb P^1$ would make the curve rational. Hence $D_n$ would be
hyperelliptic. Its genus is $n+1\ge4$, so its hyperelliptic pencil
is unique. The order-$\ell$ automorphism $a$ acts on its base
$\mathbb P^1$; this action is nontrivial, because the kernel of
the action on a separable degree-two function field has order
at most two. A nontrivial odd-order automorphism of $\mathbb P^1$
has a fixed point. The fibre over that point contains at most two
geometric points, and an odd-prime cyclic group acts trivially on
a set of size at most two. This supplies a fixed point of $a$ on
$D_n$, contradicting the freeness of the etale deck action.
Thus the proposed collision is impossible.
\end{proof}

Condition (PR.4) is a finite group computation for the actual integer
$n$ in question. No non-Wieferich assumption is needed for the
statement or proof. Hooley's theorem supplies prime-level
primitive-root information under its analytic hypotheses; it does
not supply a theta point of any of these prime or prime-power
orders, nor does it imply (PR.4) at an arbitrary higher power.

\paragraph{Attempted extension using primary projections.}
For every genus-two curve over $\overline{\mathbb F}_2$, the
Raynaud theta support is a translate of its Abel curve, since
$\deg B=1$. It therefore generates its Jacobian: a proper translate
of a subgroup containing theta would contain all differences of
points of the Abel curve, whereas those differences generate the
Jacobian.
Poonen's Theorem~1.8 consequently applies to this theta curve:
given a finite set $S_0$ of primes, there is a density-one set
$S\supseteq S_0$ for which projection of theta points onto the
$S$-primary torsion of $J$ is onto \cite[Theorem~1.8]{Poonen2005}.
In particular, take $S_0$ to contain $2$ and a selected odd prime
$\ell$, prescribe an exact-order-$\ell^r$ point in the $\ell$-part,
and prescribe zero at every other prime of $S$. A resulting theta
point has order
\[
 n=\ell^r m,\qquad (\ell,m)=1,\qquad
 \text{every prime dividing }m\text{ lies outside }S.
 \tag{PR.9}
\]
This gives odd-order points, and iterating with all previously
occurring primes in $S_0$ gives infinitely many of pairwise coprime
orders. It applies also to ordinary curves with split Jacobians.
It does not require or prove that $m=1$.
Indeed, for a nonordinary genus-two curve take $\ell=3$ and $r=1$.
Finding A10 proves that the resulting point cannot have order $3$,
so its extra factor $m$ is necessarily greater than one. The
uncontrolled factors cannot simply be omitted from this projection
argument.

Two attempted ways to remove $m$ fail at definite algebraic steps.
First, multiplying the point by $m$ removes the extra primary
components but places it on $[m](\Theta)$; this operation does not
preserve $\Theta$. Second, pass from a cyclic degree-$\ell m$
cover to its intermediate degree-$m$ base. Its relative
degree-$\ell$ factor contains the characters of every conductor
$\ell d$ with $d\mid m$, including conductor $\ell$ itself.
Theta incidence at exact order $\ell m$ controls only the
primitive conductor-$\ell m$ block. On the intermediate base the
character spaces have dimension $m$, so one zero constituent does
not make those whole spaces nilpotent.

Even if $\langle2,-1\rangle=(\mathbb Z/\ell m)^\times$, symmetry
and that group equality kill only the primitive block. They leave
the proper-conductor blocks untouched. More intrinsically, when
$n$ has at least two distinct prime factors, the least common
multiple of its proper divisors is $n$. A deck element that acts
trivially on all proper-conductor characters must therefore be
the identity; there is no nontrivial element to substitute for
$a$ in Step~4 using only the primitive-zero conclusion. For a prime
power, that least common multiple is $n/\ell$, precisely the
reason Step~3 works.

These observations limit the single-deck-element and intermediate-cover
strategies. They do not rule out a different primitive-factor argument:
A15 below uses a product of deck operators and a gonality bound to bypass
the proper-conductor factors on every rank-one genus-two base, and more
generally on a base with no separable elliptic double quotient. A16
exhibits this bypass with conductor $15$ and positive $2$-rank in the
proper-conductor factors. The remaining universal question is incidence
or orbit coverage on a single cyclic torsion subgroup; it is not a
requirement to kill every proper-conductor block by the same method.


\subsection{An explicit Lang--Weil test for pure primary theta torsion}
\label{sec:lang-weil-theta}

\paragraph{Source scope.}
The Lang quotient construction used here is the one in Poonen,
\emph{Multiples of subvarieties in algebraic groups over finite fields},
IMRN 2005, no.~24, 1487--1498, Lemma~5.1 (author PDF, p.~6).
The recovered \texttt{ANABELIAN\_SUCCESSOR}, A.10--A.11, also uses
finite-field character sums on Abel--Jacobi images.  The following
statement gives explicit constants for the genus-two theta translate
and a sufficient test for a point whose \emph{entire order} is an odd
prime power.  It is unconditional.

\paragraph{Finding A12 (a curve-dependent sufficient test).}
Let $C/\mathbf F_Q$ be a smooth projective geometrically connected
genus-two curve in characteristic two, and assume that $C$ has an
$\mathbf F_Q$-rational point.  Write $X=C^{(1)}$ for its one-fold
relative Frobenius twist,
$J=\operatorname{Pic}^0(X)$, and
\[
 0\longrightarrow\mathcal O_X\longrightarrow F_*\mathcal O_C
 \longrightarrow B\longrightarrow0
\]
for relative Frobenius and its rank-one quotient.  Let
$\Theta=\{L\in J:h^0(X,B\otimes L)>0\}$ denote the Raynaud theta
support.  For any subgroup $H\subseteq G:=J(\mathbf F_Q)$, put
$a=[G:H]$ and
$N_H=\#(\Theta(\mathbf F_Q)\cap H)$.  Then
\[
 N_H\ \ge\
 \max\left\{0,\
 \left\lceil
 \frac{Q+1-\left\lfloor2(a+1)\sqrt Q\right\rfloor}{a}
 \right\rceil\right\}.
 \tag{LW.1}
\]
In particular,
\[
       Q+1-2(a+1)\sqrt Q>a
 \tag{LW.2}
\]
ensures that $\Theta(\mathbf F_Q)\cap H$ contains a nonzero point.
The simpler condition
\[
                       Q>(2a+3)^2
 \tag{LW.3}
\]
implies (LW.2).  If $C$ is ordinary, the origin is not in $\Theta$,
and positivity of the numerator in (LW.1) already suffices.

For an odd prime $\ell$ and $r\ge1$, take
$H=J(\mathbf F_Q)[\ell^r]$.  If (LW.2) holds and
\[
             \langle 2,-1\rangle=(\mathbf Z/\ell^r)^\times,
 \tag{LW.4}
\]
then the prime-power bridge A11 applies: after a finite extension of
constants there is a connected cyclic cover of degree $\ell^s$,
for some $1\le s\le r$, whose final degree-$\ell$ Prym has $2$-rank zero.
Consequently the point-separation conclusion of A11 holds for $C$.
This last conclusion uses A11, whereas (LW.1)--(LW.3) are purely
finite-field counting statements.

\begin{proof}
\emph{1. The theta translate and the Lang quotient.}
Choose $P_0\in X(\mathbf F_Q)$.  The line bundle $B$ has degree one.
A degree-one line bundle on a genus-two curve has a nonzero section
exactly when it is $\mathcal O_X(R)$ for a unique geometric point $R$.
Thus the theta support is the image of the closed embedding
\[
 i_B:X\longrightarrow J,\qquad
 R\longmapsto\mathcal O_X(R)\otimes B^{-1}.
\]
It is the usual Abel--Jacobi embedding
$i_{P_0}(R)=\mathcal O_X(R-P_0)$ translated by the
$\mathbf F_Q$-rational point
$\tau=\mathcal O_X(P_0)\otimes B^{-1}$.
No rational effective divisor representing $B$ is required.

Let $\Phi$ denote the $Q$-power Frobenius endomorphism of $J$.
The homomorphism $\Phi-1$ has differential $-1$, so is an
\'etale isogeny.  Its kernel is the constant group scheme
$G=J(\mathbf F_Q)$.  Factoring through the constant subgroup $H$
gives \'etale isogenies
\[
  J\xrightarrow{\pi_H} J/H\xrightarrow{\psi}J,
  \qquad \psi\pi_H=\Phi-1,
  \qquad \ker\psi=G/H.
 \tag{LW.5}
\]
In particular $\deg\psi=a$.  Form the pullback
\[
                  Y=X\times_{i_B,J,\psi}(J/H).
 \tag{LW.6}
\]
It is an \'etale cover of $X$ of degree $a$.

\emph{2. Geometric connectedness, including the characteristic part.}
We give the argument to make clear that $a$ need not be odd.
Over an algebraic closure, if (LW.6) were disconnected, its
monodromy would be a proper subgroup $A_0$ of the constant abelian
deck group $A=G/H$.  Quotienting by $A_0$ would give a nontrivial
intermediate \'etale isogeny
$\psi':J'\to J$ such that $i_B$ lifts to a morphism $v:X\to J'$.
Indeed the resulting $A/A_0$-torsor on $X$ has trivial geometric
monodromy and hence has a section.
Replace $v$ by $v-v(P_0)$.  It now sends $P_0$ to zero and lifts
$i_{P_0}$.  The universal property of the Jacobian supplies a
homomorphism $u:J\to J'$ with
$u i_{P_0}=v-v(P_0)$.  It follows, again by that universal property,
that $\psi'u=\operatorname{id}_J$.  The two abelian varieties have
the same dimension, so $u$ is an isogeny and
$1=\deg(\psi'u)=\deg\psi'\deg u$, a contradiction.

This is also the full abelian \'etale fundamental-group property
of the Abel--Jacobi map; it applies to characteristic-primary
\'etale covers as well as prime-to-characteristic covers.
See Milne, \emph{Jacobian Varieties}, author version of
12 June 2021, Section~9, Proposition~9.1, p.~30;
its characteristic-primary argument is completed on pp.~31--32 \cite{MilneJV}.
Thus $Y$ is smooth, projective and geometrically connected.
Unramified Riemann--Hurwitz gives
\[
                         g(Y)=a(g(X)-1)+1=a+1.
 \tag{LW.7}
\]

\emph{3. Rational fibres and the exact count.}
For $R\in X(\mathbf F_Q)$, put $x=i_B(R)\in G$.
Choose $y\in J(\overline{\mathbf F}_Q)$ with
$(\Phi-1)y=x$, and put $z=\pi_H(y)\in\psi^{-1}(x)$.
Since $\pi_H$ is defined over $\mathbf F_Q$,
\[
                 \Phi(z)-z=\pi_H(x).
 \tag{LW.8}
\]
Thus $z$ is rational if and only if $x\in H$.
The deck group $G/H$ is constant, so in this case all $a$ points
in the fibre are rational; otherwise none is rational.
Consequently
\[
                           \#Y(\mathbf F_Q)=aN_H.
 \tag{LW.9}
\]
The Weil bound for the curve $Y$ and (LW.7) give
$\#Y(\mathbf F_Q)\ge Q+1-2(a+1)\sqrt Q$.
For the precise cohomological input, see Deligne,
\emph{La conjecture de Weil. I}, formula~(1.5.1), p.~275, and
Theorem~1.6, p.~276 \cite{WeilI}.  Since $\#Y(\mathbf F_Q)$ is an integer,
(LW.9) yields (LW.1).

There is at most one point of $X$ mapping to the origin of $J$.
Condition (LW.2) gives $N_H>1$ and hence a nonzero point of $H$ on
$\Theta$.  If $\sqrt Q>2a+3$, then
$Q-2(a+1)\sqrt Q>\sqrt Q>a-1$, proving (LW.3).
If $C$ is ordinary, Frobenius on $H^1(\mathcal O_C)$ is invertible.
The cohomology sequence defining $B$ therefore gives
$H^0(X,B)=0$, so the origin is not on $\Theta$.

\emph{4. The prime-power consequence.}
A nonzero point in $H=G[\ell^r]$ has exact order $\ell^s$ for some
$1\le s\le r$.  Reduction
$(\mathbf Z/\ell^r)^\times\to(\mathbf Z/\ell^s)^\times$ is surjective.
Thus (LW.4) implies
$\langle2,-1\rangle=(\mathbf Z/\ell^s)^\times$.
Apply A11 to this actual order, allowing the finite constant
extension needed to define the deck action and a chosen point
above the prescribed base point.
\end{proof}

\paragraph{Finding A13 (an ordinary example with order $4159$).}
Let $C/\mathbf F_2$ be the smooth projective model of
\[
                  y^2+y=x+x^{-1}+(x-1)^{-1}.
 \tag{LW.10}
\]
It is ordinary of genus two.  Its Raynaud theta support contains
at least four $\mathbf F_{128}$-rational points of exact order
$4159$.  Moreover
$\langle2,-1\rangle=(\mathbf Z/4159)^\times$.
Thus A11 applies to this ordinary curve with a cyclic cover of
prime degree $4159$.

\begin{proof}
There are exactly three simple poles in (LW.10), at $0,1,\infty$.
The Artin--Schreier genus and $p$-rank formulas give genus two and
$2$-rank two; see Pries--Zhu,
\emph{The $p$-rank stratification of Artin--Schreier curves},
Lemma~2.6, p.~712 \cite{PriesZhu2012}.  The three totally ramified points are rational.
For a polynomial model one may put $h=x(x+1)$ and $Y=hy$, obtaining
\[
               Y^2+(x^2+x)Y=x^5+x^3+x^2+x.
 \tag{LW.11}
\]

There are three points over $\mathbf F_2$, precisely the three
ramified points.  Over
$\mathbf F_4=\mathbf F_2(u)$ with $u^2+u+1=0$, the only nonpole
$x$-values are $u,u^2$.  The right side of (LW.10) at these
values is respectively $u^2,u$, both of trace one to $\mathbf F_2$.
Neither has an Artin--Schreier preimage, so
\[
                      \#C(\mathbf F_2)=\#C(\mathbf F_4)=3.
 \tag{LW.12}
\]
If $\alpha_1,\ldots,\alpha_4$ are the Frobenius eigenvalues, then
$\sum\alpha_i=0$ and $\sum\alpha_i^2=2$.
The functional equation of the genus-two Weil polynomial
therefore gives
\[
                          P(T)=T^4-T^2+4.
 \tag{LW.13}
\]
Write $\beta_1,\beta_2$ for the roots of $T^2-T+4$.
The roots of (LW.13) occur in opposite pairs with squares
$\beta_1,\beta_2$.  For odd $m$ the Jacobian point-count formula
\cite[Theorem~19.1(b)]{MilneAV} gives
\[
 \#J(\mathbf F_{2^m})=(1-\beta_1^m)(1-\beta_2^m)
                   =1-t_m+4^m,
 \quad t_m=\beta_1^m+\beta_2^m.
 \tag{LW.14}
\]
The exact recurrence
$t_0=2,\ t_1=1,\ t_m=t_{m-1}-4t_{m-2}$ yields
\[
                  (t_0,\ldots,t_7)=(2,1,-7,-11,17,61,-7,-251).
\]
Consequently
\[
                    \#J(\mathbf F_{128})=16636=4\cdot4159.
 \tag{LW.15}
\]
Here the relative Frobenius twist has the same model, since the
equation is defined over $\mathbf F_2$.

Trial division by the integers from $2$ through
$\lfloor\sqrt{4159}\rfloor=64$ proves that $4159$ is prime.
The following exact modular calculations certify its order condition:
\[
\begin{gathered}
 2079=3^3\cdot7\cdot11,\qquad 4158=2\cdot2079,\\
 2^{2079}\equiv1,\quad
 2^{693}\equiv2554,\quad
 2^{297}\equiv1245,\quad
 2^{189}\equiv102\pmod{4159}.
\end{gathered}
 \tag{LW.16}
\]
They imply $\operatorname{ord}_{4159}(2)=2079$.
This order is odd, so $-1\notin\langle2\rangle$ and
$\langle2,-1\rangle$ has order $4158$.

In (LW.1), now take $Q=128$ and
$H=J(\mathbf F_{128})[4159]$.  By (LW.15), $|H|=4159$ and $a=4$.
The coarse test already holds because $128>11^2$.
For the integer count,
$113^2<100\cdot128<114^2$ gives
\[
       N_H\ge\left\lceil\frac{129-113}{4}\right\rceil=4.
 \tag{LW.17}
\]
The curve is ordinary, so none of these theta points is the origin;
all have exact order $4159$.  This proves the assertion and
verifies both hypotheses of A11.
\end{proof}

\paragraph{Arithmetic evidence for A13.}
The standard-library certificate
\path{lang_weil_4159_certificate.py} checks the finite-field
counts, recurrence, primality, modular orders and integer bound.
An independent script
\path{verify_ordinary_4159_independent.py} obtains (LW.15) instead
as the determinant of the $11$-by-$11$ Sylvester matrix for
$P(T)$ and $T^7-1$.  These exact arithmetic checks complement
the proof above; they do not constitute formal proof-assistant
verification.

\paragraph{Finding A14 (a fixed-prime barrier to the sufficient test).}
Fix a genus-two curve over $\mathbf F_q$ and an odd prime
$\ell\ne\operatorname{char}\mathbf F_q$.  Let $\Phi$ be Frobenius on
the rank-four module $T_\ell J$, and choose $d\ge1$ with
$\Phi^d\equiv1\pmod\ell$.  One may take its order modulo $\ell$,
or the explicit integer
$d=|\operatorname{GL}_4(\mathbf F_\ell)|
 =\prod_{j=0}^3(\ell^4-\ell^j)$.
Set
\[
                 C_\ell=\ell^{v_\ell(\#J(\mathbf F_{q^d}))}.
 \tag{LW.18}
\]
For every $m\ge1$,
\[
       |J(\mathbf F_{q^m})[\ell^\infty]|
                    \le C_\ell\,\ell^{4v_\ell(m)}
                    \le C_\ell m^4.
 \tag{LW.19}
\]
Let $M$ be the first integer $m\ge4$ such that
\[
                         q^{3m}\ge4C_\ell^2m^8.
 \tag{LW.20}
\]
This is a finite, computable integer.  For every $m\ge M$ and
every $r\ge1$, condition (LW.2) fails for
$Q=q^m$ and $H=J(\mathbf F_{q^m})[\ell^r]$.
The assertion concerns the sufficient numerical test, and does
not assert absence of theta torsion.

\begin{proof}
Put $U=\Phi^d-1\in\ell M_4(\mathbf Z_\ell)$.
For every positive integer $t$, the binomial identity gives
\[
 \Phi^{dt}-1=tU V_t,\qquad
 V_t=1+\sum_{j=2}^t\frac{\binom tj}{t}U^{j-1}.
 \tag{LW.21}
\]
Since
$\binom tj/t=\binom{t-1}{j-1}/j$ and
$j-1-v_\ell(j)\ge1$ for odd $\ell$ and $j\ge2$,
every summand after $1$ lies in $\ell M_4(\mathbf Z_\ell)$.
Thus $V_t\equiv1\pmod\ell$ and its determinant is a unit.
No eigenvalue of $\Phi^d$ is one, by the Weil absolute values.
Taking determinants and valuations in (LW.21) proves the exact formula
\[
 v_\ell\#J(\mathbf F_{q^{dt}})
   =v_\ell\#J(\mathbf F_{q^d})+4v_\ell(t).
 \tag{LW.22}
\]
The inclusion
$J(\mathbf F_{q^m})\subseteq J(\mathbf F_{q^{dm}})$
and (LW.22) with $t=m$ prove (LW.19).

Write $s=q^{m/2}$ and
$a=[J(\mathbf F_{q^m}):J(\mathbf F_{q^m})[\ell^r]]$.
The Weil lower bound for the abelian surface and (LW.19) imply
\[
                     a\ge\frac{(s-1)^4}{C_\ell m^4}.
 \tag{LW.23}
\]
Condition (LW.2), after rearrangement, would require
\[
                        a<\frac{(s-1)^2}{2s+1}.
 \tag{LW.24}
\]
It is therefore impossible whenever
$(s-1)^2(2s+1)\ge C_\ell m^4$.
For $m\ge4$ we have $s\ge2$, and
$(s-1)^2(2s+1)\ge s^3/2$.
Condition (LW.20) gives $s^3/2\ge C_\ell m^4$,
proving failure of (LW.2) at that $m$.
Finally $q^{3m}/m^8$ tends to infinity and is strictly increasing
for integral $m\ge4$: its consecutive ratio is at least
$8(4/5)^8>1$.  Hence $M$ exists and the same conclusion holds
for every $m\ge M$.
\end{proof}

\paragraph{What an analytic prime theorem still does not supply.}
A12 allows a curve-dependent successful test and A13 proves that
it can succeed on an ordinary curve.  A14 rules out making this
particular test succeed automatically by enlarging constants at one
fixed odd prime.  Allowing the prime to vary would require additional
information: a sufficiently large pure prime-power subgroup of
$J(\mathbf F_{q^m})$ together with the unit-group condition (LW.4).
For example, (LW.2) requires its order to dominate the cofactor on
the scale specified by (LW.24).
The Jacobian orders are the integers
$\det(1-\Phi^m)$; a density theorem for primes at which $2$ has a
specified multiplicative order does not by itself control large
prime-power divisors of this sequence.
This is not asserted to be the strongest numerical criterion for
point separation.  The recovered \texttt{ANABELIAN\_SUCCESSOR},
A.10, gives a different rational criterion: for a quotient
$\chi:J_C(\mathbf F_Q)\twoheadrightarrow\mathbf Z/\ell$ with
kernel $K$, writing $N=\#C(\mathbf F_Q)$, the conditions
$N\ge2|K|$ and $\ell>N$, together with $2$ primitive modulo $\ell$,
give a suitable rational prime-degree cover.  That criterion can
permit larger cofactors and does not require pure-primary theta
incidence.  Its order hypothesis is stronger than (LW.4):
the prime $4159$ of A13 has $2$ of index two, not a primitive root.
The factorization and order conditions must therefore be stated
for the particular criterion being used.
Once the finite Lang quotient is specified, its rational points
are already controlled by the unconditional finite-field Weil bound.
Thus no Dedekind-zeta GRH, Artin holomorphy assumption, or
number-field least-prime theorem is a premise of A12--A14.
An inference that such an analytic theorem forces the missing
geometric incidence or the needed factorization would require a
separate argument.


% Finding A15. Prepared independently of final/.  One-fold Frobenius twist
% notation C^{(1)} agrees with the root paper's revised convention.
\subsection{Rank one permits a primitive factor of arbitrary odd conductor}
\label{sec:rank-one-composite-bridge}

\paragraph{Finding A15 (the composite primitive-factor criterion).}
Let $C/k$ be a smooth projective geometrically connected curve of genus
two over a finite field of characteristic two, with geometric $2$-rank
$f_C=1$. Let $n>1$ be odd. The following conditions suffice for the
point-section separation conclusion of A1.
First, there is a geometrically connected cyclic etale cover
$D\to C$ of degree $n$ whose primitive conductor-$n$ Jacobian factor
has $2$-rank zero. More generally, for a specified $P\in C(k)$ it
suffices to have a finite extension $L/k$, such a cover over $L$ with
constant deck group $G=\langle\sigma\rangle$, and a point
$P'\in D(L)$ over $P$, for which
\[
 A=\prod_{\ell\mid n}(1-\sigma^{n/\ell}),\qquad
 P_n=\operatorname{im}(A:J_D\longrightarrow J_D),\qquad
 P_n(L)[2^\infty]=0.                                      \tag{CB.1}
\]
The product is over distinct prime divisors of $n$.
In particular, the first condition holds if the Raynaud theta divisor
on $\operatorname{Pic}^0(C^{(1)})$ contains a point of exact odd order
$n$ and
\[
                    \langle 2,-1\rangle=(\mathbf Z/n)^\times.
                                                               \tag{CB.2}
\]
Here $C^{(1)}$ is the one-fold relative Frobenius twist. Thus for
$2$-rank one the order $n$ need not be a prime power. No assertion of
universal theta incidence is included.

\paragraph{Exact source scope.}
The geometric criterion underlying (CB.1), with the extra hypothesis
that the base is not geometrically a separable double cover of a
genus-one curve, is the recovered \texttt{ANABELIAN\_SUCCESSOR}
\emph{final row A.14}, working identifier AC.11, in
\path{agent_criteria/criteria_section.tex}. The refinement here
removes that extra hypothesis for the rank-one genus-two stratum and
combines the criterion with theta incidence. Final row A.35 is a
different, finite Fourier counting criterion; it is not the source of
the composite projector argument. We give the full argument below.

\begin{proof}
\emph{1. Rank one excludes the exceptional double quotient.}
Work first over an algebraic closure. Suppose that a genus-two curve
$C$ admits a separable degree-two map to a genus-one curve $E$.
The function-field extension is quadratic and Galois, with group
$\mathbf Z/2$. Riemann--Hurwitz says that its different has degree
$2g_C-2=2$. Every ramified point in this characteristic-two quadratic
extension has different exponent at least two, so there is exactly
one ramified point. The Deuring--Shafarevich formula consequently gives
\[
 f_C-1=2(f_E-1)+1,\qquad f_C=2f_E\in\{0,2\}.              \tag{CB.3}
\]
This contradicts $f_C=1$. We use the geometric $p$-rank form of that
formula, stated for $p$-group covers, for example, in
Anbar--Beelen--Nguyen, Theorem~1, equation~(1)
\cite{AnbarBeelenNguyenP}; no assumption of ordinarity is needed.

\emph{2. The primitive projector produces a small separable function.}
Suppose that distinct points $P,Q$ collided in the covering formulation
of A1. The pointed cover in (CB.1), after restricting the collision
to $L$, supplies $Q'\in D(L)$ above $Q$ with
$\delta=[Q'-P']\in J_D(L)[2^\infty]$. Hence $A\delta=0$.
Put $r=\omega(n)$ and $d=2^r$. Each element
$\sigma^{n/\ell}$ has order $\ell$, and these elements belong to
distinct prime-primary factors of the cyclic group. Their $2^r$
subset products are therefore all distinct. Expanding $A(Q'-P')$
gives a nonzero principal divisor with exactly $d$ simple zeros and
$d$ simple poles. They lie in the two disjoint free deck orbits of
$P'$ and $Q'$. Let $f$ be the resulting rational function on
$D_{\overline L}$. It has degree $d$ and is separable: its simple
zeros (or poles) rule out a square in the function field.

The relevant degree bound is
\[
                         n>(2^{\omega(n)}-1)^2=(d-1)^2.
                                                               \tag{CB.4}
\]
It holds for every odd $n>1$, with no additional size assumption.
For $r=1$ use $n\ge3>1$; for $r=2$ use $n\ge15>9$.
For $r\ge3$, the distinct odd prime factors give
$n\ge3\cdot5^{r-1}>4^r>(2^r-1)^2$; the middle inequality holds
at $r=3$ and its ratio increases by $5/4$ thereafter.

\emph{3. The field generated by the conjugate functions.}
Write $K=\overline L(D)$,
$M=\overline L(\gamma f:\gamma\in G)$, and $e=[K:M]$.
Poonen's Castelnuovo--Severi consequence gives
\[
                  e\mid d,\qquad g(M)\le(d/e-1)^2.
                                                               \tag{CB.5}
\]
Here is the entire induction needed from that consequence. Start with
one degree-$d$ function, whose rational field has index $d$ and genus
zero. If an intermediate field has index $d/s$ and genus at most
$(s-1)^2$, adjoining another degree-$d$ function with relative degree
$a$ gives index $d/(as)$. Castelnuovo--Severi bounds its genus by
\[
 a(s-1)^2+(a-1)(as-1)
  =(as-1)^2-as(a-1)(s-1)\le(as-1)^2.
\]
Induction over the finitely many conjugates proves (CB.5).
This is exactly the use of Proposition~2.1, Lemma~2.2, and
Corollary~2.3 of \cite{Poonen2007Gonality}, over the perfect field
$\overline L$.

Since $D\to C$ is etale, $g(D)=n+1$. Equations (CB.4)--(CB.5)
exclude $e=1$. Thus $e>1$ is a power of two. Let $Z$ be the smooth
projective curve with function field $M$, and let $v:D\to Z$ have
degree $e$. The group $G$ preserves $M$. Its induced action on $Z$
is faithful, since a kernel $H$ would satisfy
$|H|\mid[K:M]=e$ as well as $|H|\mid n$.
The action on $Z$ is also free. Indeed, a stabilizer of a geometric
point of $Z$ would contain an element of odd prime order $\ell$.
It would partition the fibre divisor of $v$ into orbits of length
$\ell$, since the deck action on $D$ is free; fibre multiplicities
are constant on these orbits. This forces $\ell\mid e$, impossible.

Riemann--Hurwitz for the free action now yields
\[
                    g(Z)-1=n\bigl(g(Z/G)-1\bigr).
\]
The bound $0\le g(Z)\le(d/e-1)^2<n$ forces
$g(Z)=g(Z/G)=1$: quotient genus zero would give $g(Z)=1-n<0$,
and quotient genus at least two would give $g(Z)\ge n+1$.

\emph{4. A reduced fibre forces the forbidden degree-two map.}
Write $f=h\circ v$, with $h$ a nonconstant function on $Z$.
Choose a zero $z$ of $h$. The fibre divisor $v^*z$ is reduced,
because the zero divisor of $f$ is reduced. It has $e$ distinct
geometric points, all contained in the two original deck orbits.
Each such orbit maps injectively to $Z$, by equivariance and the
free action on $Z$. Hence the fibre contains at most one point from
each orbit, giving $e\le2$. Thus $e=2$.

The extension $K/\overline L(f)$ is separable, so $K/M$ is separable.
Its quadratic involution $\tau$ is normalized by $G$, because $G$
preserves its fixed field $M$; it is centralized by $G$ since
$\operatorname{Aut}(\mathbf Z/2)=1$. Therefore taking the quotient
by $G$ yields a separable degree-two map
\[
                  C=D/G\longrightarrow Z/G
\]
to a genus-one curve. Its degree is two explicitly because
$[K:M^G]=2n$ and $[K:K^G]=n$.
This contradicts Step~1 and proves the
criterion (CB.1). A primitive factor of geometric $2$-rank zero
has no nonzero geometric $2$-power torsion points. Its rational
$2$-primary group is therefore zero after every finite constant
extension, so this geometric condition implies the required
pointed version of (CB.1) for every base point.

\emph{5. Theta incidence annihilates every primitive orbit.}
Let $L$ have exact order $n$ on the Raynaud theta divisor of
$C^{(1)}$, and let $B$ denote its locally exact differential bundle.
In characteristic two $B$ is a line bundle with $B^2\simeq\omega$.
Serre duality therefore gives
$h^0(B\otimes L)=h^0(B\otimes L^{-1})$, so both signs lie on theta.
For relative Frobenius $F:C\to C^{(1)}$, tensor its exact sequence
by $L$. The induced map
\[
 H^1(C^{(1)},L)\longrightarrow H^1(C,F^*L)
\]
has a nonzero kernel, since $H^0(B\otimes L)\ne0$.
Both spaces are lines: the two line bundles are nontrivial of
degree zero on genus-two curves. The map is consequently zero.
The same holds for $L^{-1}$.

Choose $b$ with $2b\equiv1\pmod n$ and put $M_0=(F^*L)^b$ on $C$.
Then $M_0$ has exact order $n$ and
$F^*(M_0^{(1)})=M_0^2=F^*L$. Since Frobenius pullback is
injective on odd-order torsion, $M_0^{(1)}=L$. On the connected
cyclic cover associated to $M_0$, label the eigensheaves by $M_0^i$.
For $i\ne0\bmod n$ the cohomology spaces
$V_i=H^1(C,M_0^i)$ are one-dimensional, and Frobenius maps
$V_i$ to $V_{2i}$. The arrows $1\to2$ and $-1\to-2$ just shown
are zero. Condition (CB.2) makes every doubling orbit in
$(\mathbf Z/n)^\times$ contain $1$ or $-1$. Thus every primitive
orbit has nilpotent Frobenius, and its stable part is zero.

Finally $A$ acts on character $i$ by
$\prod_{\ell\mid n}(1-\zeta_\ell^i)$, which is nonzero precisely
when $\gcd(i,n)=1$. Since $n$ is odd, the character projectors are
defined after an extension of $\mathbf Z_2$ and split the
$2$-primary representation. The image $P_n$ is exactly the
primitive factor up to isogeny; its stable Frobenius part is the
primitive sum just considered. It therefore has $2$-rank zero.
All the geometric cover data descend after a finite constant
extension, proving the stated theta consequence.
\end{proof}

A18 below removes this geometric exclusion for the geometric
primitive-rank-zero criterion in characteristic two, in every genus
at least two. The rational-only criterion retains the distinction.

Steps~2--5 apply to any genus-two base known not to admit a geometric
separable double cover of a genus-one curve; the rank-one hypothesis
is what makes this geometric exclusion automatic in Step~1.

\paragraph{How many prime factors can the single-point condition allow?}
For odd $n$, the Chinese remainder theorem and the cyclicity of units
modulo an odd prime power give
\[
 (\mathbf Z/n)^\times/\bigl((\mathbf Z/n)^\times\bigr)^2
                  \simeq (\mathbf Z/2)^{\omega(n)}.
\]
If (CB.2) holds, this quotient is generated by the images of $2$ and
$-1$, so necessarily $\omega(n)\le2$. This is an arithmetic
restriction on the one-point sufficient condition, not on the
geometric criterion (CB.1). Exact finite examples are
\[
\begin{array}{c|rrrrrr}
 n&15&21&35&45&33&105\\ \hline
 \varphi(n)&8&12&24&24&20&48\\
 \operatorname{ord}_n(2)&4&6&12&12&10&12\\
 |\langle2,-1\rangle|&8&12&24&24&10&24
\end{array}
\]
Thus the first four integers satisfy (CB.2); having at most two
distinct prime factors is not by itself sufficient, as $n=33$ shows.
These computations certify only the unit groups. Theta incidence at
these orders is a separate condition, illustrated by A16 for $n=15$.

\paragraph{Several theta points: the exact coverage condition.}
For any fixed cyclic torsion class $M_0$ of odd order $n$, define
\[
 Z(M_0)=\{i\in(\mathbf Z/n)^\times:
                    M_0^{(1)i}\in\Theta_C\}.
\]
The primitive factor of its cyclic cover has $2$-rank zero if and
only if $Z(M_0)$ meets every orbit of multiplication by $2$ on
$(\mathbf Z/n)^\times$. Indeed, the arrow $V_i\to V_{2i}$ is
zero exactly at the displayed theta incidence. A cyclic sequence
of maps of lines has zero stable part exactly when at least one
arrow is zero; if every arrow is nonzero, their composition is a
bijective semilinear map of a line. This proves the equivalence.
Consequently, on the rank-one base, sufficiently many theta points
\emph{in this one cyclic subgroup} suffice, even when $\omega(n)>2$.
Condition (CB.2) and symmetry give this coverage from a single point.

\paragraph{What the primary-projection theorem still leaves open.}
Applying Poonen's theorem with $S_0\supset\{2,\ell\}$ can prescribe
an $\ell^a$ component, but gives order $n=\ell^a m$, where all primes
of $m$ lie outside a subsequently chosen density-one set $S$.
In the proof, Lemma~6.4 first finds a point $h$ in a bounded-index
subgroup supported on a density-zero prime set and then takes
$U$ to be \emph{all} prime divisors of its order. It bounds the
density quantity $\eta(U)$, not $|U|$ or $\omega(\operatorname{ord}h)$.
Enlarging finite $S_0$ therefore does not give the extra restriction
$\omega(m)\le1$ needed even to qualify for (CB.2), or the unit-group
equality itself \cite[Theorem~1.8 and Lemmas~6.4--6.6]{Poonen2005}.
For $\ell=3$, A10 forces $m>1$ on the rank-one stratum; A15 now
permits such a cofactor, but does not manufacture a suitable one.
Theta points of pairwise coprime orders likewise do not establish
the displayed orbit-coverage condition: on their fibre-product
cover they control proper-conductor characters, while the primitive
characters are tensor products whose theta incidence is not supplied.
This identifies the additional incidence statement that would be
needed to complete this proposed route.


\subsection{A checked composite-degree witness}

\paragraph{Finding A16 (exact order fifteen, with ordinary proper factors).}
Let $k=\mathbb F_2[z]/(z^3+z+1)\simeq\mathbb F_8$, and let
\[
 C/k:\quad Y^2+X^2Y=X^5+(z+1)X^3+X,\qquad
 P=\infty,\qquad R=(z^2,z).
 \tag{CE.1}
\]
Then $C$ has genus two and $2$-rank one. The class
$M=\mathcal O_C(R-P)$ has exact order $15$, and its Frobenius twist
$M^{(1)}$ lies on the Raynaud theta divisor of $C$.
There is a pointed connected cyclic etale cover of degree $15$
after a finite constant extension, of genus $16$ and $2$-rank $7$.
Its primitive conductor-$15$ factor has dimension $8$ and $2$-rank
zero, whereas its conductor-$3$ and conductor-$5$ factors are
ordinary of dimensions $2$ and $4$. Thus the whole new factor has
positive $2$-rank $6$. Finding A15 nevertheless separates every
proposed pair of distinct colliding point sections on this base.

\begin{proof}
The reciprocal coordinate $x=1/X$, $y=Y/X^2$ gives
$y^2+y=x^3+(z+1)x+x^{-1}$. Its two poles have orders $3$ and $1$,
so the genus is two and the $2$-rank is one, by the same
Artin--Schreier calculation used in A09
\cite[Lemma~2.6]{PriesZhu2012}. The point $P$ is the unique point
above $X=\infty$, with pole orders $2$ and $5$ for $X$ and $Y$.
As in A09, $\operatorname{div}(dX/X^2)=2P$; consequently
\[
 B^1_C\simeq\mathcal O_{C^{(1)}}(P^{(1)}),\qquad
 \Theta_C=C^{(1)}-P^{(1)}.
 \tag{CE.2}
\]
This twist is retained explicitly: the coefficients of (CE.1) are
in $\mathbb F_8$, not all in $\mathbb F_2$.

Define the following polynomials over $k$:
\begin{align*}
 A(X)={}&z+(z+1)X+(z^2+z)X^2+(z^2+z+1)X^3\\
        &+(z^2+z)X^4+(z^2+1)X^6+(z^2+z+1)X^7,\\
 B(X)={}&1+(z+1)X+z^2X^2+(z^2+z)X^3\\
        &+(z^2+z+1)X^4+X^5,\qquad F=A+BY.
\end{align*}
Reduction modulo $z^3+z+1$ gives the exact identities
\begin{align}
 A^2+X^2AB+(X^5+(z+1)X^3+X)B^2&=(X+z^2)^{15},
 \tag{CE.3}\\
 F(R)=0,\qquad F(\iota R)&=z\ne0,
 \quad\iota(X,Y)=(X,Y+X^2).
 \tag{CE.4}
\end{align}
The independent symbolic checker verifies every coefficient of
(CE.3) and both evaluations. Since $X=z^2$ is unramified, the norm
identity proves that the only finite zero of $F$ is $R$, with
multiplicity $15$. There are no finite poles. The leading term
$X^5Y$ has pole order $15$ at $P$, while $A$ has pole order at
most $14$. Thus
\[
                       \operatorname{div}(F)=15(R-P).
 \tag{CE.5}
\]

The class is nonzero, since a principal divisor $R-P$ would make
$C$ rational. Order $3$ is excluded by
$h^0(3P)=h^0(2P)=2$, because $2P$ is canonical. To exclude order
$5$, a function with divisor $5(R-P)$ would lie in
\[
 H^0(C,\mathcal O(5P))=\langle1,X,X^2,Y\rangle.
\]
Its $Y$ coefficient would be nonzero; after scaling it would be
$Y+a_2X^2+a_1X+a_0$. Its norm has coefficient $1$ at $X$, whereas
the monic polynomial $(X+r)^5$, $r=X(R)=z^2$, has coefficient
$r^4=z\ne1$ there. This is a contradiction. The only remaining
divisor of $15$ is $15$, proving the exact order. Equation (CE.2)
now shows that $M^{(1)}$ is on theta, as claimed. In the normalization
of A11 one obtains the original class $M$, since
$F_C^*(M^{(1)})=M^2$ and $2^{-1}\equiv8\pmod{15}$.

The normalization of $W^{15}=F$ is etale because every valuation
of $F$ is divisible by $15$. It is geometrically connected:
being a third or fifth power would make respectively
$5(R-P)$ or $3(R-P)$ principal, contrary to the exact order.
For the parameter $t=X^2/Y$ at $P$, the unit $t^{15}F$ has residue
$1$, so $tW=1$ gives a $k$-point above $P$. Over
$\mathbb F_{4096}$ the deck group is the constant group
$\mathbb Z/15$, and Riemann--Hurwitz gives genus $16$.

The units modulo $15$ consist of the two multiplication-by-$2$
orbits
\[
                \{1,2,4,8\},\qquad\{14,13,11,7\}.
\]
The theta point and its inverse give one zero Frobenius arrow on
each orbit, so the primitive factor has $2$-rank zero by the
character calculation of A11. It has dimension $\varphi(15)=8$.
The conductor-$3$ factor is ordinary of dimension $2$, by A10.

For the conductor-$5$ factor, the same norm-coefficient argument
above shows that every nonzero order-$5$ point on $C-P$ has
$X$ coordinate $1$. But at $X=1$ equation (CE.1) is
$Y^2+Y=z+1$, and $\operatorname{Tr}_{\mathbb F_8/\mathbb F_2}(z+1)=1$.
Thus there is no such $k$-point. All four nonzero powers of
$M^3$ are $k$-rational. The Abel embedding and Frobenius twisting
preserve this rationality, so none lies on the relevant theta
support. All four character arrows are therefore bijective, and
the conductor-$5$ factor is ordinary of dimension $4$.
The character decomposition, whose idempotents have odd
denominators, now gives
$f_D=1+2+4+0=7$. Finally A15 applies to the rank-one base and its
rank-zero primitive factor. This proves the point-separation claim
without asserting that the whole new factor has rank zero.
\end{proof}

\paragraph{Computation scope.}
The retained producer \texttt{construct\_theta15.py} obtains the
coefficients by Taylor linear algebra in $\mathbb F_8$.
The separate \texttt{verify\_theta15\_symbolic.py} imports no
producer arithmetic and reduces the polynomial identity using
SymPy over $\mathbb F_2$. Both pass. The example was found in a
finite search of $5488$ nonbranch rational points on $784$
displayed rank-one models over $\mathbb F_8$; this search is not
claimed to classify geometric curves or all torsion orders.
No claim is made that this curve lacks a prime-power witness:
the example verifies the composite-cover conclusion itself.

\paragraph{A smaller prime-degree comparison.}
On the rank-one curve
\[
 C_5/\mathbb F_2:\quad
 Y^2+X^2Y=X^5+X^4+X^3+X,
 \qquad R_5=(1,0),\quad P_5=\infty,
\]
the function $F_5=Y+X+1$ has norm $(X+1)^5$, vanishes at $R_5$,
has value $1$ at its conjugate, and has unique pole of order $5$
at $P_5$. Hence $\operatorname{div}(F_5)=5(R_5-P_5)$, and the
nonzero theta class has exact order $5$. Since $2$ is primitive
modulo $5$, A03 supplies a degree-$5$ separating cover. Meanwhile
$Q_5=(0,0)$ satisfies $\operatorname{div}(X)=2(Q_5-P_5)$, so
$P_5$ and $Q_5$ have the same prime-to-$2$ abelian Kummer class.
The pole calculation and theta-characteristic calculation are
identical to those above; this comparison is a direct proved
special case, not a claim of minimality among all covers.


\subsection{A positive-density arithmetic supply for composite conductors}
\label{sec:three-times-artin-prime}

\paragraph{Finding A17 (an explicit family of candidate orders).}
Assume the Dedekind-GRH premise M. Let
\[
 \mathcal Q=\{q\text{ prime}:q\equiv5\pmod8,
                         \operatorname{ord}_q(2)=q-1\}.
\]
Then $\mathcal Q$ has natural density $\mathcal A/2$ among the
primes, where
\[
 \mathcal A=\prod_{\ell\ {
m prime}}
                      \left(1-\frac1{\ell(\ell-1)}\right)>0
\]
is Artin's constant. For every $q\in\mathcal Q$,
\[
                  \langle2,-1\rangle=(\mathbf Z/3q)^\times.
                                                               \tag{AF.1}
\]
The analytic input used here is Lenstra's fixed finite Frobenius
condition version of Hooley's theorem. Its sufficient GRH family is
\[
 \zeta_{K_n}(s),\qquad
 K_n=\mathbf Q(\zeta_n,2^{1/n}),\qquad
 n\ge1\text{ squarefree}.                                    \tag{AF.2}
\]
In particular, this is a number-field Kummer Dedekind-zeta family,
which is contained in M. The density is an arithmetic candidate
supply; it is not a theta-incidence assertion.

\begin{proof}
\emph{1. Degrees and the finite congruence field.}
For odd squarefree $m>1$, the prime $2$ is unramified in
$\mathbf Q(\zeta_m)$. At any prime above $2$, the polynomial
$T^m-2$ is Eisenstein, because the valuation of $2$ is one.
It is irreducible there and hence over $\mathbf Q(\zeta_m)$.
Consequently
\[
                         [K_m:\mathbf Q]=m\varphi(m).
                                                               \tag{AF.3}
\]
The Galois group is the affine semidirect product
$C_m\rtimes(\mathbf Z/m)^\times$. For clarity, a generator
$\tau$ of $C_m$ sends $2^{1/m}$ to $\zeta_m2^{1/m}$, while the
cyclotomic automorphism indexed by $b$ may be chosen to fix
$2^{1/m}$ and sends $\zeta_m$ to $\zeta_m^b$.
Its conjugation sends $\tau$ to $\tau^b$. Since $m$ is odd,
$b=2$ is available, and its commutator with $\tau$ equals
$\tau$. The commutator subgroup therefore contains $C_m$,
and equals it because the quotient by $C_m$ is abelian.
Every abelian subextension of $K_m/\mathbf Q$ thus lies in
$\mathbf Q(\zeta_m)$.

Let $F=\mathbf Q(\zeta_8)$. Its intersection with $K_m$ is an
abelian extension of $\mathbf Q$, so lies in
$\mathbf Q(\zeta_m)\cap\mathbf Q(\zeta_8)=\mathbf Q$.
The latter equality follows from the coprime cyclotomic conductors.
Thus
\[
                       [FK_m:\mathbf Q]=4m\varphi(m).
                                                               \tag{AF.4}
\]
Equations (AF.3)--(AF.4) also hold for $m=1$ with $K_1=\mathbf Q$.

\emph{2. The density from finite inclusion-exclusion.}
The condition $q\equiv5\pmod8$ prescribes the automorphism
$\sigma_5:\zeta_8\mapsto\zeta_8^5$ on $F$. This automorphism
acts nontrivially on $\sqrt2=\zeta_8+\zeta_8^{-1}$, so $2$ is a
quadratic nonresidue at these primes. The prime-$2$ obstruction
to primitive order is already excluded.
For an odd prime $\ell$, the obstruction that
$\ell\mid(q-1)/\operatorname{ord}_q(2)$ is exactly complete
splitting in $K_\ell$, outside the finitely many ramified primes.
For a finite set $E$ of odd primes, the compositum of their
$K_\ell$ is $K_{\prod_{\ell\in E}\ell}$.
By (AF.4), the simultaneous condition of splitting in $K_m$ and
having Frobenius $\sigma_5$ on $F$ has density
$1/(4m\varphi(m))$ for odd squarefree $m$.
Finite inclusion-exclusion therefore gives the density
\[
 \sum_{m\mid\prod_{\ell\in E}\ell}
               \frac{\mu(m)}{4m\varphi(m)}.
                                                               \tag{AF.5}
\]
Lenstra's Theorem~3.1, with $K=\mathbf Q$, $W=\langle2\rangle$,
and the fixed Frobenius condition $(F,\{\sigma_5\})$, identifies
the limiting finite Chebotarev density with the natural density
under (AF.2) \cite[\S2 and Theorem~3.1]{Lenstra1977}.
Thus the analytic infinite-limit step is an application of that
theorem, not an inference from finite inclusion-exclusion alone.
The limit of (AF.5), an absolutely convergent Euler product, is
\[
 \frac14\prod_{\ell\ {
m odd}}
                \left(1-\frac1{\ell(\ell-1)}\right)
                         =\frac{\mathcal A}{2}>0.
\]

\emph{3. The composite unit group.}
For $q\in\mathcal Q$, the Chinese remainder theorem gives
\[
 \operatorname{ord}_{3q}(2)
       =\operatorname{lcm}(2,q-1)=q-1=\frac{\varphi(3q)}2.
\]
A power of $2$ equal to $-1$ modulo $q$ has exponent
$(q-1)/2$ modulo $q-1$, which is even because $q\equiv5\pmod8$.
A power equal to $-1$ modulo $3$ has odd exponent. Therefore
$-1\notin\langle2\rangle$ modulo $3q$. Adjoining it doubles
the subgroup of order $\varphi(3q)/2$, proving (AF.1).
\end{proof}

\paragraph{The exact remaining geometric gate.}
For a fixed genus-two curve over a finite field of
characteristic two, a theta point of exact order $3q$ for even one
$q\in\mathcal Q$ would satisfy A15 on its original locus, and
A18 on every genus-two base. Finding A16 exhibits this for
one particular curve and $q=5$. A17 supplies infinitely many
candidate primes for each fixed curve, but does not establish
such a point for every curve. Poonen's projection theorem with
$S_0\supseteq\{2,3\}$ supplies order $3m$ after prescribing an
exact order-three component; A10 forces $m>1$. Its proof neither
forces $m$ to be a single prime nor places that prime in
$\mathcal Q$. Those additional incidence requirements remain
unproved; the positive density in A17 does not remove them.


\subsection{A primitive rank-zero cover removes the elliptic exception}
\label{sec:geometric-primitive-all-genera}

\paragraph{Finding A18 (geometric primitive criterion in all genera).}
Let $C/k$ be a smooth projective geometrically connected curve of
genus $g\ge2$ over a finite field of characteristic two.
Suppose that, over a finite extension of constants, there is a
geometrically connected cyclic \'etale cover $D\to C$ of odd degree
$n>1$ whose primitive conductor-$n$ Jacobian factor has geometric
$2$-rank zero.  Then the finite-field point-section separation
conclusion in the covering formulation of A1 holds for $C$.
There is no hypothesis excluding a separable double map from $C$
to a genus-one curve, and no restriction that $n$ be a prime power.

More precisely, with constant deck group $G=\langle\sigma\rangle$,
write
\[
 A=\prod_{\ell\mid n}(1-\sigma^{n/\ell}),\qquad
 P_n=\operatorname{im}(A:J_D\longrightarrow J_D),
 \tag{GP.1}
\]
where the product runs over the distinct prime divisors of $n$.
Under a putative collision, if $P_n(L)[2^\infty]=0$ on a pointed
model over a finite field $L$, the exceptional outcome of the
projector argument necessarily satisfies
\[
                         f(P_n)=r\varphi(n)>0
 \tag{GP.2}
\]
for an integer $r\ge1$.  Here $r$ is the number of geometric
ramification points of a resulting separable degree-two map
$C_{\overline L}\to E$ with $g(E)=1$.
For genus two, $r=1$, and the entire nontrivial new factor of
$D/C$ in this exceptional case is ordinary.
The geometric rank-zero hypothesis therefore eliminates this
exception.  A rational torsion condition alone is not being
replaced by a geometric one.

\begin{proof}
\emph{1. The projector reduction.}
Extend the constants so that the proposed distinct colliding
points $P,Q$, the cover, its cyclic deck group and a chosen lift
$P'\in D(L)$ are all defined over one finite field $L$.
Geometric $2$-rank zero implies
$P_n(L)[2^\infty]=0$ over every such extension.
The covering formulation of a collision gives a lift $Q'$ over
$Q$ with
\[
                       \delta=[Q'-P']\in J_D(L)[2^\infty].
\]
Hence $A\delta=0$.

We recall the complete geometric reduction from A15,
Steps~2--4; those steps do not use the rank or genus-two
hypothesis on the base except for the lower bound on $g(D)$.
Put $s=\omega(n)$ and $d=2^s$.
The $2^s$ subset products of the elements $\sigma^{n/\ell}$
are distinct, because they lie in independent prime-primary
factors of the cyclic group.  The two deck orbits of $P',Q'$
are free and disjoint.  Thus the expansion of the principal
divisor $A(Q'-P')$ has precisely $d$ distinct simple zeros and
$d$ distinct simple poles.  The resulting function $f$ has
separable degree $d$.

Over $\overline L$, put
\[
 K=\overline L(D),\qquad
 M=\overline L(\gamma f:\gamma\in G),\qquad e=[K:M].
\]
The Castelnuovo--Severi induction in A15 gives
\[
                 e\mid d,\qquad g(M)\le(d/e-1)^2.
 \tag{GP.3}
\]
The exact published input is Poonen,
\emph{Gonality of modular curves in characteristic $p$},
Proposition~2.1, Lemma~2.2 and Corollary~2.3,
author p.~3 \cite{Poonen2007Gonality}.  The induction adjoins functions of degree $d$;
if the previous index is $d/u$ and the adjoining degree is $a$,
the new index is $d/(au)$ and its genus is at most
$a(u-1)^2+(a-1)(au-1)\le(au-1)^2$.
For every odd $n>1$ one has $n>(d-1)^2$:
for $s=1,2$ use $n\ge3,15$, and for $s\ge3$ use
$n\ge3\cdot5^{s-1}>4^s$.
Since $g(D)=n(g-1)+1\ge n+1$, (GP.3) excludes $e=1$.
Thus $e>1$ is a power of two.

Let $Z$ be the smooth projective curve with function field $M$.
The induced $G$-action on $Z$ is faithful: a kernel would have
order dividing both $e$ and $n$.
It is free: a stabilizer containing an element of odd prime
order $\ell$ would partition a fibre divisor of $D\to Z$
of degree $e$ into orbits of length $\ell$, with constant
multiplicities on each orbit, forcing $\ell\mid e$.
Consequently
\[
                 g(Z)-1=n\bigl(g(Z/G)-1\bigr).
\]
Together with $g(Z)\le(d/e-1)^2<n$, this forces
$g(Z)=g(Z/G)=1$.

Write $f=h\circ v$ for $v:D\to Z$.
A fibre above a zero of $h$ is reduced, since the zero divisor
of $f$ is reduced.  It has $e$ points, all in the two original
deck orbits.  Each of these orbits maps injectively to $Z$,
because the action on $Z$ is free.  Hence $e\le2$, so $e=2$.
The extension $K/M$ is separable and quadratic, with involution
$\tau$.  It is separable because it is an intermediate extension
of the separable extension $K/\overline L(f)$.
Since $G$ preserves $M$, it normalizes $\langle\tau\rangle$;
the unique nontrivial element is fixed by conjugation, so
$\tau$ commutes with $G$.

\emph{2. The exceptional diagram is an actual pullback.}
Set
\[
 F=K^G=\overline L(C),\qquad
 E_0=M^G=K^{G\times\langle\tau\rangle}.
\]
Let $E$ be the smooth projective curve with function field $E_0$.
We already know that $g(E)=1$.
The group $G\times\langle\tau\rangle$ acts faithfully on $K$,
and has order $2n$, so
$K/E_0$ is Galois of that degree.  The fields fixed by
$G$ and by $\langle\tau\rangle$ have intersection $E_0$ and
compositum $K$.  Equivalently,
\[
       [F:E_0]=2,\qquad [M:E_0]=n,\qquad FM=K.
 \tag{GP.4}
\]
In particular $C\to E$ is separable of degree two and
$Z\to E$ is the cyclic \'etale cover of degree $n$ already
obtained from the free $G$-action.
The fibre product $C\times_E Z$ is connected by (GP.4).
It is smooth, since its projection to $C$ is the base change
of an \'etale morphism.  Its function field is $K$, so it is
the curve $D$ itself.

For each divisor $m\mid n$, let $H_m$ be the subgroup of $G$
of order $n/m$, and set
\[
       Z_m=Z/H_m,\qquad D_m=D/H_m.
\]
The same field calculation gives
\[
        D_m=C\times_E Z_m,\qquad
        \deg(Z_m/E)=m,\qquad \deg(D_m/C)=m.
 \tag{GP.5}
\]
Every $Z_m\to E$ is \'etale between genus-one curves.
After choosing compatible origins over $\overline L$ it is
an isogeny of odd degree.  Such an isogeny induces an
isomorphism on the geometric $2$-primary Tate modules:
its kernel has odd order, and multiplication by its degree
is invertible on every $2$-primary group.  Therefore
\[
                              f(Z_m)=f(E).
 \tag{GP.6}
\]

\emph{3. Ramification and the exact ranks of all intermediate covers.}
Let $r$ be the number of geometric ramification points of
$C\to E$.  Riemann--Hurwitz gives different of degree
$2g-2>0$, so $r\ge1$.  This is a quadratic Galois extension
in characteristic two.  Its base change by the \'etale map
$Z_m\to E$ has exactly $mr$ ramification points: above each
branch point are $m$ distinct points of $Z_m$, and \'etale
base change preserves the local ramification data.
The Deuring--Shafarevich formula for this degree-two cover,
together with (GP.6), gives
\[
             f(D_m)-1=2\bigl(f(Z_m)-1\bigr)+mr,
 \qquad
             f(D_m)=2f(E)-1+mr.
 \tag{GP.7}
\]
The formula counts the geometric ramification points, since
each has ramification index two.  It does not replace that
count by the degree of the different.
For the exact primary locator of the formula, see
Anbar--Beelen--Nguyen,
\emph{A new tower with good $p$-rank meeting Zink's bound},
Theorem~1, equation~(1), author p.~3 \cite{AnbarBeelenNguyenP}.
Its hypotheses are a Galois $p$-group extension of function
fields over an algebraically closed field, as here.

\emph{4. Conductor decomposition forces positive primitive rank.}
Consider
$V=T_2(J_D)\otimes_{\mathbf Z_2}\mathbf Q_2$,
whose dimension is the geometric $2$-rank of $D$.
The odd group $G$ acts semisimply on $V$.
Pullback and norm for $D\to D_m$ have composition
$[n/m]$ in one direction and the sum of the subgroup
$H_m$ in the other.  Since $n/m$ is odd, they identify
$T_2(J_{D_m})\otimes\mathbf Q_2$ with $V^{H_m}$.

After extending the coefficient field to split $G$, decompose
$V$ according to the exact orders of its characters.
Let $w_j$ be the total dimension of the characters of exact
order $j$, for $j\mid n$.  The invariants just identified give
\[
                           f(D_m)=\sum_{j\mid m}w_j.
 \tag{GP.8}
\]
On character $i$ modulo $n$, the endomorphism $A$ in (GP.1)
acts by
$\prod_{\ell\mid n}(1-\zeta_\ell^i)$.
This is nonzero exactly when $\gcd(i,n)=1$.
Thus its image has the primitive character part as its
$2$-adic Tate module up to isogeny, and
$f(P_n)=w_n$.
M\"obius inversion of (GP.8), using (GP.7) and $n>1$, yields
\[
\begin{aligned}
 f(P_n)
   &=\sum_{m\mid n}\mu(n/m)\bigl(2f(E)-1+mr\bigr)\\
   &=r\sum_{m\mid n}\mu(n/m)m
     =r\varphi(n)>0.
\end{aligned}
 \tag{GP.9}
\]
This contradicts the geometric rank-zero hypothesis.
Hence the proposed collision cannot occur.
\end{proof}

\paragraph{Genus-two consequence and the rational distinction.}
For $g=2$, the different of a separable quadratic map to a
genus-one curve has degree two.  Each ramified point in
characteristic two has different exponent at least two.
Thus $r=1$, with a single ramified point of different exponent
two.  Formula (GP.7) gives
$f(D)-f(C)=n-1$, equal to $g(D)-g(C)$.
The pullback/norm decomposition of $J_D$ into $J_C$ and its
new factor, up to an odd-degree isogeny, therefore shows that
the whole new factor is ordinary.  In particular its primitive
factor is ordinary of dimension $\varphi(n)$.

Consequently, even the rational condition
$P_n(L)[2^\infty]=0$ suffices on a genus-two base if one also
knows that $P_n$ is \emph{not ordinary}: a collision would force
it to be ordinary.  Without this extra geometric condition,
the rational-only statement of A15 retains its stated
elliptic-quotient qualification.  In all genera, (GP.9)
describes the positive primitive rank forced by the
exceptional case.

The theta calculation of A15 now gives a smaller geometric
gate on \emph{every} genus-two curve in characteristic two,
of any $2$-rank: an exact odd-order theta point with
$\langle2,-1\rangle=(\mathbf Z/n)^\times$ is sufficient,
and so is the exact same-subgroup orbit-coverage condition
there.  The arithmetic family $n=3q$ of A17 is therefore
available as a family of candidate orders on ordinary
bielliptic bases as well.
This does not assert incidence for every curve.  The
one-dimensional eigenspace argument establishing that
theta criterion remains specific to genus two; A18's
all-genus assertion assumes a primitive factor of geometric
rank zero directly.


\subsection{A finite exact search for each prescribed odd prime support}
\label{sec:finite-support-theta}

\paragraph{Finding A19 (an explicit field for all primary exponents).}
Let $C/\mathbf F_Q$ be a smooth projective geometrically connected
genus-two curve in characteristic two. Write $X=C^{(1)}$ for its
one-fold relative Frobenius twist, $J=\operatorname{Pic}^0(X)$,
and $\Theta\subset J$ for its Raynaud theta support. Thus
\[
 \Theta=i_B(X),\qquad
 i_B(R)=\mathcal O_X(R)\otimes B^{-1},
                                                               \tag{FS.1}
\]
where $B$ is the locally exact differential line bundle of degree
one on $X$. This is a genus-two Abel curve translated over
$\mathbf F_Q$; a rational effective divisor for $B$ is not required.
Let $S$ be a finite set of odd primes, and write $J(\overline{\mathbf
F}_Q)\{S\}$ for the points whose orders have every prime divisor in
$S$, with arbitrary exponents. Let $\Phi$ be $Q$-power Frobenius on
$J$, and choose $d\ge1$ such that
\[
            \Phi^d=1\quad\hbox{on }J[\ell]
                         \qquad(\ell\in S).                  \tag{FS.2}
\]
Then
\[
 \Theta(\overline{\mathbf F}_Q)\cap J(\overline{\mathbf F}_Q)\{S\}
   =\Theta(\mathbf F_{Q^d})\cap J(\mathbf F_{Q^d})\{S\}.
                                                               \tag{FS.3}
\]
In particular the left side is finite and is computable by a
terminating finite-field search. One may take
\[
 \begin{split}
 d&=\operatorname{lcm}_{\ell\in S}
                  \operatorname{ord}(\Phi\mid J[\ell]),
       \quad\hbox{or the explicit larger integer}\\
 d&=\operatorname{lcm}_{\ell\in S}
                    |\operatorname{GL}_4(\mathbf F_\ell)|,
 \qquad |\operatorname{GL}_4(\mathbf F_\ell)|
                   =\prod_{j=0}^3(\ell^4-\ell^j).
 \end{split}                                                   \tag{FS.4}
\]
For $S=\varnothing$ take $d=1$ and the torsion group to be $\{0\}$.
No GRH or conjecture about prime distribution is a premise.

\paragraph{Literature scope.}
This is a specialized effective finite-support torsion statement
of Boxall type, proved in full below. The bibliographic locator
for Boxall's general work is
\emph{Sous-vari\'et\'es alg\'ebriques de vari\'et\'es
semi-ab\'eliennes sur un corps fini}, in \emph{Number Theory
(Paris, 1992--1993)}, LMS Lecture Note Ser.~215, Cambridge, 1995,
69--80. The primary paper of Scanlon--Voloch discusses related torsion
finiteness results and lists this reference
\cite[Introduction and reference \textnormal{[Box]}]{ScanlonVoloch1999}.
The original Boxall chapter's theorem text was not accessible in
this check; no theorem number, exact bound, or originality claim
is inferred from its bibliographic record. The proof and the
specific field (FS.4) do not depend on that unavailable text.

\begin{proof}
\emph{1. Two points is the exact geometric intersection bound.}
For every nonzero $T\in J(\overline{\mathbf F}_Q)$,
\[
                  \#\bigl(\Theta\cap(\Theta+T)\bigr)\le2
                                                               \tag{FS.5}
\]
for geometric points. Here is a proof using only genus two.
Suppose three distinct points of this intersection correspond to
pairs $(R_i,S_i)$ on $X$ with
$[R_i-S_i]=T$, for $i=1,2,3$. The $R_i$ are distinct, and so
are the $S_i$, because the Abel map is an embedding.
For $i=2,3$ we have
$R_1+S_i\sim R_i+S_1$. These two effective degree-two divisors
are distinct: equality, with $R_1\ne R_i$, would force
$R_1=S_1$ and hence $T=0$.
Thus the corresponding degree-two linear system moves.
Riemann--Roch in genus two forces its divisor class to be
canonical: for a degree-two divisor $D$, $h^0(D)\ge2$ implies
$h^0(K_X-D)\ge1$, and a degree-zero line bundle with a section
is trivial. Therefore both $S_2$ and $S_3$ are the unique
effective degree-one divisor in the class $K_X-R_1$.
This contradicts $S_2\ne S_3$ and proves (FS.5).

\emph{2. An elementary powering lemma.}
Let $\ell$ be odd, let $V=(\mathbf Q_\ell/\mathbf Z_\ell)^4$,
and let $U\in\operatorname{GL}_4(\mathbf Z_\ell)$ satisfy
$U\equiv1\pmod\ell$. If $v\in V$ is a finite-order point and
$w=(U-1)v$ has exact order $\ell^s$, $s\ge1$, then
\[
        (U^{\ell^{s-1}}-1)v=\ell^{s-1}w\ne0,
                                                               \tag{FS.6}
\]
and this point has exact order $\ell$.
To check the equality put $N=\ell^{s-1}$ and $M=U-1$.
The binomial expansion is
\[
 (U^N-1)v=Nw+\sum_{j=2}^{N}\binom Nj M^{j-1}w.
\]
For $j\ge2$,
\[
 v_\ell\!\left(\binom Nj\right)\ge s-1-v_\ell(j),
 \qquad j-2-v_\ell(j)\ge0.
\]
The first follows from
$\binom Nj=(N/j)\binom{N-1}{j-1}$, and the second uses that
$\ell$ is odd. Since $M\in\ell M_4(\mathbf Z_\ell)$, every
summand after $Nw$ is divisible by $\ell^s$ as an operator on
$w$, and is zero. For $s=1$ the sum is empty. This proves
(FS.6). In addition $U^{\ell^{s-1}}$ fixes every point of
$V[\ell]$, again because $U\equiv1\pmod\ell$.

\emph{3. Isolate one prime component of a nonrational point.}
Let $z\in\Theta(\overline{\mathbf F}_Q)\cap J\{S\}$ and
write $z=\sum_{\ell\in S}z_\ell$ for its primary decomposition.
Put $A=\Phi^d$. Its action on every $T_\ell J$ is congruent
to the identity modulo $\ell$. For a component of order
$\ell^a$, the action of $A$ belongs to the kernel of
\[
 \operatorname{GL}_4(\mathbf Z/\ell^a)
                         \longrightarrow\operatorname{GL}_4(\mathbf F_\ell),
\]
which is an $\ell$-group. Hence the orbit length of $z_\ell$
under $A$ is a power of $\ell$.

Suppose $Az\ne z$ and choose $\ell$ for which
$Az_\ell\ne z_\ell$. For every other prime $t\in S$, choose
a power of $t$ equal to the orbit length of $z_t$ under $A$.
Let $c$ be the product of these orbit lengths. Then
$\gcd(c,\ell)=1$, and $U=A^c$ fixes every other component,
but still moves $z_\ell$: raising a permutation of
$\ell$-power order to the prime-to-$\ell$ power $c$ cannot
turn a nontrivial orbit into a fixed point.

Use the identification
$J[\ell^\infty](\overline{\mathbf F}_Q)
 \simeq T_\ell J\otimes(\mathbf Q_\ell/\mathbf Z_\ell)$.
Let $s\ge1$ be the order exponent of $(U-1)z_\ell$, and put
$\sigma=U^{\ell^{s-1}}$. By (FS.6),
\[
                      \sigma z=z+T,
 \qquad T\in J[\ell]\setminus\{0\},\qquad \sigma T=T.
                                                               \tag{FS.7}
\]
The other primary components remain fixed. The map $\sigma$
on geometric points is a positive integral power of $Q$-power
Frobenius. As $\Theta$ is defined over $\mathbf F_Q$, all
\[
                \sigma^jz=z+jT\qquad(0\le j<\ell)
\]
belong to $\Theta$. They are $\ell$ distinct points. Each also
belongs to $\Theta+T$, by using the preceding point in the
cyclic translation orbit. This contradicts (FS.5), since
$\ell\ge3$. Consequently $Az=z$, which is exactly
$z\in J(\mathbf F_{Q^d})$. The reverse inclusion in (FS.3)
is immediate.

\emph{4. An explicit order bound and terminating computation.}
Let
\[
 h_d=\#J(\mathbf F_{Q^d}),\qquad
 N_S=\prod_{\ell\in S}\ell^{v_\ell(h_d)}.
                                                               \tag{FS.8}
\]
Every point in (FS.3) has order dividing $N_S$, and
\[
                 N_S\le h_d\le(Q^{d/2}+1)^4.                 \tag{FS.9}
\]
The last inequality is the genus-two Jacobian Weil bound
\cite[Theorem~19.1(c)]{MilneAV}.
If $P_J(T)=\det(T-\Phi)$ is its integer Weil polynomial, then
\[
 h_d=\prod_{j=1}^4(1-\alpha_j^d)
       =\operatorname{Res}_T(P_J(T),T^d-1),                 \tag{FS.10}
\]
where the $\alpha_j$ are its roots; the sign is positive
because there are four roots
\cite[Theorem~19.1(b)]{MilneAV}.
Thus $h_d$, its valuations at the finitely many primes in
$S$, and $N_S$ are exact finite integer computations.
Using the explicit second choice in (FS.4) does not require
computing the torsion representation first.

Enumerate $X(\mathbf F_{Q^d})$, form the line-bundle classes
$i_B(R)$ in the finite Jacobian, and keep precisely those killed
by $N_S$. This enumerates the entire geometric intersection
in (FS.3). Their exact orders, and any desired doubling-orbit
coverage within their cyclic subgroups, can then be tested by
finite group operations. A smaller bound may be obtained by
replacing $h_d$ with the exponent of its $S$-primary subgroup,
but this is unnecessary for termination. No assertion of
practical running time at the potentially large degree $d$ is made.
\end{proof}

\paragraph{Consequence for the remaining composite-order search.}
For each fixed prime $q$ in A17, taking $S=\{3,q\}$ now makes
the search over \emph{all} orders $3^a q^b$ finite, not only
the search at order $3q$. The actual order must still satisfy
the unit-group test or the same-subgroup orbit-coverage condition.
The primes $q$ range over an unbounded set, so this result does
not prove that one of the finite searches succeeds for every
curve.

There is also an exact reformulation of the squarefree gate.
For $q\ne3$, the curve $[3](\Theta)$ contains a nonzero point
of order $q$ if and only if $\Theta$ contains a point of order
$q$ or $3q$. Indeed a lift $z\in\Theta$ of such a point
satisfies $3qz=0$ but $3z\ne0$, so has one of those two orders;
the converse follows by multiplying by three.
For $q$ in A17 either order meets the genus-two sufficient
condition of A18. Applying Poonen's multiples theorem to the
generating curve $[3](\Theta)$ still permits a varying additional
integer multiplier. It does not put a prime-order point on that
curve itself. Thus neither this reformulation nor (FS.3)
replaces the remaining incidence assertion by prime density.


\subsection{A rank-one curve with no nonzero three-primary theta point}
\label{sec:no-three-primary-theta}

\paragraph{Finding A20 (all powers of three fail on the degree-thirteen example).}
For the genus-two rank-one curve of A09,
\[
 C/\mathbf F_2:\quad Y^2+X^2Y=X^5+X,
\]
its Raynaud theta divisor satisfies
\[
       \Theta_C(\overline{\mathbf F}_2)
          \cap J_C[3^\infty](\overline{\mathbf F}_2)=\{0\}.
 \tag{TP.1}
\]
Thus allowing arbitrary powers of the single prime $3$ does not
repair the universal positive-rank incidence assertion. This does
not obstruct point separation on this curve: A09 gives its exact
order-$13$ theta point and separating cover.

\begin{proof}
The curve and its Frobenius twist have the same equation over
$\mathbf F_2$. Finding A09 gives its Frobenius characteristic
polynomial on every odd-adic Tate module:
\[
                         P(T)=T^4+T^3+2T+4.
 \tag{TP.2}
\]
Let $\overline P$ be its reduction modulo $3$. Exact polynomial
division in $\mathbf F_3[T]$ gives
\[
                       T^{20}-1=\overline P(T)Q(T),
 \tag{TP.3}
\]
where
\[
\begin{aligned}
 Q(T)={}&T^{16}-T^{15}+T^{14}+T^{12}+T^{11}+T^{10}\\
        &-T^6+T^5-T^4-T^2-T-1.
\end{aligned}
\]
Cayley--Hamilton applied to the action of the $2$-power Frobenius
$\Phi$ on $J_C[3]$ therefore proves $\Phi^{20}=1$ there.
All $3^4$ geometric points of $J_C[3]$ are consequently rational
over $\mathbf F_{2^{20}}$.

The determinant formula for the rational points of a Jacobian
gives
\[
\begin{aligned}
 \#J_C(\mathbf F_{2^{20}})
   &=\det(1-\Phi^{20})
     =\operatorname{Res}(P(T),T^{20}-1)\\
   &=1099812538656
     =3^4\cdot13577932576,
\end{aligned}
 \tag{TP.4}
\]
and $13577932576\equiv1\pmod3$.
The resultant has the displayed sign because $P$ has even degree.
One exact computation uses the integral companion matrix
\[
 C_P=\begin{pmatrix}
 0&0&0&-4\\
 1&0&0&-2\\
 0&1&0&0\\
 0&0&1&-1
 \end{pmatrix}
\]
and evaluates $\det(I-C_P^{20})$; an independent polynomial
resultant calculation gives the same integer.
The determinant formula and Frobenius convention are those of
\cite[Theorem~19.1(b), pp.~40--41]{MilneAV}.

The rational $3$-primary subgroup in (TP.4) has exactly $3^4$
points and already contains all of $J_C[3]$, so it equals
$J_C[3]$ and is killed by $3$.
Apply A19 with the prime set $S=\{3\}$ and the Frobenius exponent
$d=20$ established in (TP.3). It puts every geometric
$3$-primary theta point in $J_C(\mathbf F_{2^{20}})$.
Every nonzero such point would thus have exact order $3$.
Finding A10 excludes order-three theta points on every nonordinary
genus-two curve, and A09 gives $f_C=1$. This is a contradiction.
Finally the origin is on theta, since in A09 the locally exact
differential line bundle is $\mathcal O_C(P)$ for the point at
infinity. This proves equality in (TP.1).
\end{proof}

The standard-library certificate \texttt{check\_three\_primary.py}
performs the companion-matrix and modular polynomial computations.
The independently written SymPy certificate
\texttt{verify\_three\_primary.py} uses a polynomial resultant
and imports no producer code. These verify the finite arithmetic
in (TP.3)--(TP.4); the geometric exclusion also uses the separately
proved A19 and A10.


\subsection{Generic avoidance of every odd torsion order}
\label{sec:generic-all-odd}

\paragraph{Finding A21 (both positive-rank genus-two strata).}
Let $k$ be an algebraically closed field of characteristic two,
let $f\in\{1,2\}$, and let $Z$ be an irreducible component of
$\mathcal M_2^f$.  Let $C$ denote its geometric generic curve,
put $X=C^{(1)}$, and let $B$ be the locally exact differential
line bundle on $X$.  Its Raynaud theta support
\[
 \Theta=\{L\in\operatorname{Pic}^0(X):
                       H^0(X,B\otimes L)\ne0\}
\]
contains no nonzero point of odd order:
\[
 \Theta\cap
 \bigl(\operatorname{Pic}^0(X)[n]\setminus\{0\}\bigr)
             =\varnothing\qquad(n\ge1\text{ odd}).             \tag{GO.1}
\]
Consequently, for every odd $n>1$, every connected cyclic
degree-$n$ etale cover of $C$ has ordinary whole new factor.
This includes prime powers and composite degrees on the
$2$-rank-one stratum.  All assertions are unconditional.

The statement about one geometric generic curve does not
assert that a single nonempty open subset avoids all odd
orders simultaneously.  For any prescribed finite list of
odd integers the avoidance conditions do hold on a common
nonempty open, and the finite-field specialization argument
of A08 then applies for either $f=1$ or $f=2$.

\paragraph{Exact published inputs.}
Achter--Pries, \emph{Monodromy of the $p$-rank strata of the
moduli space of curves}, IMRN 2008, Art.~rnn053,
Theorem~4.5 (author version p.~14), gives full
$\operatorname{Sp}_{2g}(\mathbf F_\ell)$ and
$\operatorname{Sp}_{2g}(\mathbf Z_\ell)$ monodromy for
$\ell\ne p$, including $g=2$ and $f=1,2$
\cite{AchterPries2008}.
For the generic-point interpretation used here, the proof
of Ozman--Pries, Proposition~4.3, printed p.~465, explicitly
states that the full symplectic level-$\ell$ torsor $Z[\ell]$
is irreducible \cite[Proposition~4.3, proof]{OzmanPries2019}.
Thus this input is stronger than mere
connectedness of a cover over a possibly nonnormal base.
The tricanonical framing construction in Achter--Pries,
Section~5.1 and Lemma~5.2, pp.~17--18, passes to a scheme
without reducing monodromy.  The lifting assertion for a
closed subgroup of $\operatorname{Sp}_4(\mathbf Z_\ell)$
surjecting modulo $\ell$ is also
Landesman--Swaminathan--Tao--Xu,
\emph{Lifting subgroups of symplectic groups over
$\mathbf Z/\ell\mathbf Z$}, Research in Number Theory \textbf{3}
(2017), Art.~14, Theorem~1; their Proposition~1(b) says that
$\operatorname{Sp}_4(\mathbf Z/\ell^a)$ is perfect for every
odd $\ell$ and $a\ge1$ \cite{LSTX2017}.  Finally, Hall,
\emph{Classical Groups and Geometry}, version 6 May 2015,
Theorem~12.26, printed p.~173, gives the center and simplicity
of $\operatorname{PSp}_4(\mathbf F_\ell)$ for all odd $\ell$,
including $\ell=3$ \cite{Hall2015}.

\begin{proof}
\emph{1. The monodromy statement at the geometric generic
point.}
Use tricanonical frames to replace $Z$ by a scheme carrying
its universal curve.  The frame morphism has geometrically
integral $\operatorname{GL}_5$ fibres.  In particular its
pullback to the irreducible full level-$\ell$ torsor
$Z[\ell]$ is irreducible.  Restrict to a nonempty smooth
affine open $T$ of the resulting integral scheme.  The
level-$\ell$ torsor remains irreducible on $T$, so its
geometric generic Galois image is the full
$\operatorname{Sp}_4(\mathbf F_\ell)$.
The closed $\ell$-adic image over this same generic field
is consequently $\operatorname{Sp}_4(\mathbf Z_\ell)$ by
the lifting theorem just cited.  This reasoning works for
every odd $\ell$ on the same generic curve.  It does not
assume that deleting a singular locus from an arbitrary
irreducible base preserves its fundamental group.

These statements apply equally to $J=\operatorname{Pic}^0(X)$.
Indeed relative Frobenius induces an isogeny between the
Jacobians of $C$ and $X$ of degree a power of two.  It is an
isomorphism on every odd-order Tate module, and is a
similitude of the corresponding alternating pairings.
Conjugating a full symplectic group by a similitude leaves
that symplectic group unchanged.

\emph{2. Independence at distinct odd primes.}
Write $n=\prod_{\ell\mid n}\ell^{a_\ell}$ and
$G_\ell=\operatorname{Sp}_4(\mathbf Z/\ell^{a_\ell})$.
Each $G_\ell$ is perfect.  Its reduction kernel modulo
$\ell$ is an $\ell$-group, as is seen by filtering the
congruence kernel by successive powers of $\ell$.
If $G_\ell$ has a nontrivial simple quotient, that quotient
is nonabelian, since $G_\ell$ is perfect.  The image of the
congruence kernel is a normal $\ell$-subgroup of this simple
quotient and must be trivial.  The quotient therefore
factors through $\operatorname{Sp}_4(\mathbf F_\ell)$.
Its center has trivial image in a nonabelian simple group,
and Hall's theorem now implies that the simple quotient is
exactly $\operatorname{PSp}_4(\mathbf F_\ell)$.
The orders
\[
 \bigl|\operatorname{PSp}_4(\mathbf F_\ell)\bigr|
       =\frac{\ell^4(\ell^2-1)(\ell^4-1)}2                  \tag{GO.2}
\]
strictly increase with the positive integer $\ell$.
Thus factors belonging to distinct odd primes have no
common nontrivial simple quotient.

For clarity, this implies the required product assertion
without an independence assumption.  If a subgroup
$H\subset A\times B$ surjects to $A$ and $B$, Goursat's
argument identifies it with the fibre product over
isomorphic quotients $A/N_A\simeq B/N_B$.  If $H$ is
proper, this common quotient is nontrivial and has a
nontrivial simple quotient.  This would be a common simple
quotient of $A$ and $B$.  Induct on the number of prime
factors.  A nonabelian simple quotient of a product of
perfect groups factors through one factor: the images of
the factors are commuting normal subgroups, and at most
one can be the whole nonabelian simple group.  There are
no abelian simple quotients.  The obstruction just
described is therefore absent at each induction step.

The generic Galois image on $J[n]$ surjects to every
$G_\ell$ by Step~1 and is contained in their product.
Consequently it is
\[
            \prod_{\ell\mid n}G_\ell
                  =\operatorname{Sp}_4(\mathbf Z/n).         \tag{GO.3}
\]
The equality on the right uses the Chinese remainder
decomposition of matrices and their defining symplectic
equations.

\emph{3. Transitivity and the genus-two intersection bound.}
Over the local ring $R=\mathbf Z/\ell^a$, a vector of exact
order $\ell^a$ is primitive.  If $v$ is primitive, the
perfect alternating form admits $w$ with
$\langle v,w\rangle=1$.  The span of $v,w$ is a
nondegenerate free direct summand; its orthogonal
complement is a free rank-two symplectic module.
Completing $v,w$ to a symplectic basis proves that
$\operatorname{Sp}_4(R)$ is transitive on primitive
vectors.  Applying this in each Chinese remainder factor
proves that the group in (GO.3) is transitive on the
points of exact order $n$.

Because $B$ has degree one, $\Theta$ is the translated
Abel curve
\[
                 R\longmapsto\mathcal O_X(R)\otimes B^{-1}.
                                                               \tag{GO.4}
\]
For every nonzero $T\in J$, it satisfies
\[
                   \#\bigl(\Theta\cap(\Theta+T)\bigr)\le2.
                                                               \tag{GO.5}
\]
This is the bound proved in A19.  A short proof is included
to specify its exact scope.  Three distinct intersection
points give pairs $(R_i,S_i)$ with $[R_i-S_i]=T$.
The $R_i$ and the $S_i$ are separately distinct.  For
$i=2,3$, the distinct effective divisors
$R_1+S_i\sim R_i+S_1$ move in a degree-two linear system.
Riemann--Roch in genus two makes this class canonical.
Both $S_2$ and $S_3$ would then be the unique effective
degree-one divisor in the class $K_X-R_1$, a contradiction.
If the two displayed degree-two divisors were equal as
divisors, their equality and $R_1\ne R_i$ would instead
force $T=0$.  This also is excluded.

Suppose now that $\Theta$ contains a point of exact odd
order $n>1$.  The line bundle $B$, and hence its theta
support, is defined over the generic field of the family.
Galois invariance and (GO.3) force $\Theta$ to contain
every point of exact order $n$.  Choose coordinates
$J[n]\simeq(\mathbf Z/n)^4$, let $e_1,e_2$ be two
distinct coordinate vectors, and choose an odd prime
$\ell\mid n$.  Set
\[
                   P=e_1,\qquad T=(n/\ell)e_2.
\]
The point $T$ has exact order $\ell$ and all
$P+jT$, $0\le j<\ell$, have exact order $n$, since their
first coordinate is $1$.  They therefore lie on $\Theta$.
Each also lies on $\Theta+T$, using the preceding point
in this cyclic list.  There are $\ell\ge3$ such distinct
points, contradicting (GO.5).  This excludes every exact
odd order and proves (GO.1).

\emph{4. The whole new factor of an odd cyclic cover.}
Work over the algebraic closure of the generic field.
Let $\pi:D\to C$ be a connected cyclic etale cover of
odd degree $n$.  Choose a line bundle $M$ of exact order
$n$ defining the cover, so that
\[
 \pi_*\mathcal O_D=\bigoplus_{i=0}^{n-1}M^i,\qquad
 H^1(D,\mathcal O_D)=\bigoplus_{i=0}^{n-1}H^1(C,M^i).
                                                               \tag{GO.6}
\]
For $i\ne0$, the latter summand has dimension one.
Put $L_i=(M^i)^{(1)}$ on $X=C^{(1)}$.  It is nontrivial
of odd order and $F_{C}^*L_i=(M^i)^2$.  Tensoring
$0\to\mathcal O_X\to F_{C*}\mathcal O_C\to B\to0$
by $L_i$ and using the projection formula gives
\[
 0\longrightarrow L_i\longrightarrow F_{C*}(M^i)^2
                 \longrightarrow B\otimes L_i
                 \longrightarrow0.                          \tag{GO.7}
\]
By (GO.1), $H^0(X,B\otimes L_i)=0$.  Thus
$H^1(X,L_i)\to H^1(C,M^{2i})$ is injective, hence
bijective between its one-dimensional source and target.
Under Frobenius-semilinear identification this is precisely
the arrow from character $i$ to character $2i$ in (GO.6).
Multiplication by two permutes all nonzero residues
modulo $n$, so Frobenius is bijective on the entire
nontrivial-character sum.

Its stable rank is therefore $n-1$; the trivial-character
sum has stable rank $f$.  Consequently
$f(D)=f+n-1$.  The norm--pullback isogeny decomposition
$J_D\sim J_C\times\operatorname{Prym}(D/C)$ and
Riemann--Hurwitz give
\[
 f\bigl(\operatorname{Prym}(D/C)\bigr)=n-1
       =\dim\operatorname{Prym}(D/C),
\]
which proves the final assertion.
\end{proof}

\paragraph{Scope relative to the earlier results.}
For prime $n$, the conclusion agrees with Ozman--Pries,
Theorem~1.1; for the ordinary stratum, Yang's
new-ordinarity theorem already supplies the arbitrary
odd-degree conclusion used in A08.
The argument here is a direct genus-two consequence of
published full monodromy and the intersection bound.  It
does not require an extension of the compact-type
degeneration argument from prime to composite degrees.
The full joint monodromy in (GO.3) also supplies the finite
level needed in A22 below.


\subsection{No fixed finite odd prime support works uniformly}
\label{sec:no-fixed-prime-support}

\paragraph{Finding A22 (all exponents and all products).}
Let $S$ be any finite set of odd primes and let
$f\in\{1,2\}$.  There is a finite field $\mathbf F_Q$ of
characteristic two and a smooth projective geometrically
connected genus-two curve $C/\mathbf F_Q$ of $2$-rank $f$
such that, for $X=C^{(1)}$, $J=\operatorname{Pic}^0(X)$ and
its Raynaud theta support $\Theta$,
\[
       \Theta(\overline{\mathbf F}_Q)\cap
       \bigl(J(\overline{\mathbf F}_Q)\{S\}\setminus\{0\}\bigr)
                         =\varnothing.                       \tag{SA.1}
\]
Here $J\{S\}$ allows arbitrary powers of the primes in
$S$ and arbitrary products of those powers.
The curve can be chosen in any prescribed nonempty open
of any geometric irreducible component of
$\mathcal M_2^f$.  In particular these examples are
Zariski dense in each such component.

This strengthens finite-list avoidance: no fixed finite
prime support, even with all exponents and mixed orders
allowed, supplies a universal pure odd-torsion theta
witness on either positive-rank stratum.  The assertion
is unconditional.  The field and curve may depend on $S$;
no uniform degree bound for that field is asserted.

\begin{proof}
The case $S=\varnothing$ is immediate, so assume that $S$
is nonempty and put $N=\prod_{\ell\in S}\ell$.

\emph{1. A family with a finite avoidance condition.}
Start with a component of $\mathcal M_2^f$ over
$\overline{\mathbf F}_2$.  Use the tricanonical frame scheme
and choose a nonempty smooth affine open $T$ as in A21.
There is a universal smooth proper genus-two curve
$\mathcal C\to T$, every fibre has $2$-rank $f$, and the
geometric generic monodromy on every odd level is full.
One can also restrict to the inverse image of any
prescribed nonempty moduli open.

Write $\mathcal X=\mathcal C^{(1)}$ for relative Frobenius
twist over $T$ and
$\mathcal J=\operatorname{Pic}^0(\mathcal X/T)$.
Relative Frobenius is finite locally free of degree two.
In the sequence
\[
 0\longrightarrow\mathcal O_{\mathcal X}
  \longrightarrow F_*\mathcal O_{\mathcal C}
  \longrightarrow\mathcal B\longrightarrow0,                 \tag{SA.2}
\]
the quotient $\mathcal B$ is locally free of rank one.
This can be checked in a relative etale coordinate, where
$1,t$ is the usual basis for relative Frobenius.
It has degree one on every fibre.  Consequently
\[
 \mathcal X\longrightarrow\mathcal J,\qquad
 R\longmapsto\mathcal O_{\mathcal X_t}(R)
                         \otimes\mathcal B_t^{-1}             \tag{SA.3}
\]
is a relative translated Abel map.  Its construction uses
the relative diagonal and the Picard functor; a rational
base point on the fibres is unnecessary.  Its image is
closed, because $\mathcal X$ is proper over $T$ and
$\mathcal J$ is separated over $T$, and its geometric
fibre image is precisely the Raynaud theta support.

The scheme
$\mathcal J[N]^\times=\mathcal J[N]\setminus\{0\}$
is finite etale over $T$: the zero section is open and
closed in this finite etale scheme.  The inverse image
of the closed image of (SA.3) in
$\mathcal J[N]^\times$ is closed, and its projection to
$T$ is closed because the projection is finite.
It misses the generic point by A21.  Remove this proper
closed subset and replace what remains by a nonempty
affine open $U$.  Then every geometric fibre over $U$
satisfies
\[
                  \Theta_t\cap
                  (\mathcal J_t[N]\setminus\{0\})
                           =\varnothing.                     \tag{SA.4}
\]
This is a finite incidence condition, so no passage from
infinitely many open subsets to one open has occurred.

All data descend to a sufficiently large finite field
$k_0=\mathbf F_{q_0}$.  Enlarge that constant field so
that $U$ is geometrically integral and
$q_0\equiv1\pmod{N^2}$.  The latter is possible because
$N$ is odd.  The family, its fixed-rank condition,
and (SA.4) are preserved by this constant extension.

\emph{2. A connected symplectic level torsor.}
Choose a primitive $N^2$th root of unity in $k_0$ to
identify the values of the Weil pairing with the standard
constant pairing.  Let
\[
 V=\operatorname{Isom}_{\mathrm{sympl}}
       \bigl((\mathbf Z/N^2)^4_U,\mathcal J[N^2]\bigr)
                      \longrightarrow U.                    \tag{SA.5}
\]
This is a finite etale torsor for the constant group
\[
       G=\operatorname{Sp}_4(\mathbf Z/N^2)
        =\prod_{\ell\in S}
                       \operatorname{Sp}_4(\mathbf Z/\ell^2).
\]
Its geometric generic monodromy is all of $G$ by A21.
The restriction of $V_{\overline{k}_0}$ to the generic
point of $U_{\overline{k}_0}$ is therefore connected.
Since $U$ is smooth and
geometrically integral, $V$ is geometrically integral.
Equivalently, its geometric connectedness follows from
the full monodromy, and smoothness makes distinct
irreducible components disjoint.
The total arithmetic image is also $G$: it contains
the full geometric image, and it is contained in $G$
because $\mu_{N^2}$ is constant over $k_0$.

\emph{3. Realizing a precise Frobenius class by a twist.}
For $\ell\in S$ put $u_\ell=1+\ell$ modulo $\ell^2$ and
choose
\[
 \gamma_\ell=
 \operatorname{diag}(u_\ell,u_\ell,
                         u_\ell^{-1},u_\ell^{-1})
     \in\operatorname{Sp}_4(\mathbf Z/\ell^2).
                                                               \tag{SA.6}
\]
Use the standard symplectic basis ordered as
$e_1,e_2,f_1,f_2$.  Then
\[
 \gamma_\ell\equiv I\pmod\ell,\qquad
 (\gamma_\ell-I)/\ell
                   \equiv\operatorname{diag}(1,1,-1,-1)
                             \pmod\ell.                       \tag{SA.7}
\]
The second matrix is invertible.
Let $\gamma=(\gamma_\ell)_{\ell\in S}\in G$.
Each $\gamma_\ell$ has order $\ell$, so $\gamma$ has
order $N$.

Here is a complete existence argument for a specialization
with this Frobenius class.  It is the finite-field
twisting proof of Chebotarev, with only qualitative
Lang--Weil needed.  Let $\mathcal F$ denote arithmetic
$q_0$-power Frobenius acting on the geometric points of
$V$.  For a symplectic frame $v$ define a new descent
action of the arithmetic generator by
\[
                       v\longmapsto
                               \mathcal F(v)\gamma^{-1}.
                                                               \tag{SA.8}
\]
Right multiplication by $\gamma^{-1}$ is a deck
automorphism defined over $k_0$ and commutes with
$\mathcal F$.  Since $\gamma^N=1$, this descent action
agrees with the original action on
$\operatorname{Gal}(\overline{k}_0/\mathbf F_{q_0^N})$.
It therefore defines a finite Galois descent datum on
$V_{\mathbf F_{q_0^N}}$.  The datum is effective because
$V$ is affine; the exact general reference is the
Stacks Project, Lemma~35.37.1, Tag~\texttt{0245}
\cite{Stacks0245}.
Write $V^\gamma$ for the descended twist.
The projection to $U$ descends because deck
automorphisms act trivially there.

The twist $V^\gamma$ is geometrically integral, being
geometrically isomorphic to $V$.  Lang--Weil gives
$V^\gamma(\mathbf F_{q_0^m})\ne\varnothing$ for every
sufficiently large $m$.  One precise reference for this
form, including affine varieties, is Poonen,
\emph{Rational points on varieties}, Theorem~7.7.1(ii)--(iii),
printed p.~222 \cite{Poonen2017RPV}.  Choose such an $m$
with $m\equiv1\pmod N$ and a rational point of
$V^\gamma$ over $\mathbf F_{q_0^m}$.  Under the geometric
identification with $V$, it gives a frame $v$ over a
point $t\in U(\mathbf F_{q_0^m})$ such that
\[
                \mathcal F^m(v)\gamma^{-m}=v,
       \qquad\text{hence}\qquad
                \mathcal F^m(v)=v\gamma.                      \tag{SA.9}
\]
Thus the $Q$-power Frobenius of the fibre, where
$Q=q_0^m$, has matrix $\gamma$ on its $N^2$-torsion in
the frame $v$.  Changing the frame conjugates this
matrix, which leaves the properties in (SA.7) unchanged.
This argument does not require $t$ to have minimal
residue-field degree $m$; the fibre is considered over
the stated field $\mathbf F_Q$.

\emph{4. All rational $S$-primary torsion has squarefree
exponent.}
Let $C/\mathbf F_Q$ be this fibre and
$\Phi$ its $Q$-power Frobenius on $J$.
For each $\ell\in S$, lift the chosen torsion frame to
a basis of $T_\ell J$.  Equations (SA.7)--(SA.9) imply
\[
             \Phi-I=\ell A_\ell,\qquad
                         A_\ell\in
                         \operatorname{GL}_4(\mathbf Z_\ell).
                                                               \tag{SA.10}
\]
On
$J[\ell^\infty](\overline{\mathbf F}_Q)
 \simeq(\mathbf Q_\ell/\mathbf Z_\ell)^4$,
the matrix $A_\ell$ is an automorphism.  Consequently
\[
 \begin{split}
 J(\mathbf F_Q)[\ell^\infty]
     &=\ker\bigl(\Phi-I\mid
                  J[\ell^\infty](\overline{\mathbf F}_Q)\bigr)\\
     &=J[\ell](\overline{\mathbf F}_Q).
 \end{split}                                                   \tag{SA.11}
\]
Indeed $\ell A_\ell v=0$ if and only if
$\ell v=0$.  In particular all the $\ell$-torsion
points are rational, and there are no rational points
of order $\ell^2$ or higher.  Taking the direct sum
over $S$ gives
\[
                  J(\mathbf F_Q)\{S\}
                          =J[N](\overline{\mathbf F}_Q).
                                                               \tag{SA.12}
\]
This equation includes mixed orders dividing $N$.

\emph{5. Excluding arbitrary exponents geometrically.}
Since $\Phi$ is the identity on $J[\ell]$ for every
$\ell\in S$, Finding~A19 applies with $d=1$.
It puts every geometric $S$-primary theta point in
$J(\mathbf F_Q)\{S\}$.  By (SA.12) this is $J[N]$,
and (SA.4) excludes every nonzero point of that group
from theta.  This proves (SA.1).

The construction started inside an arbitrary prescribed
nonempty open of a component.  Repeating it for each
such open proves the asserted Zariski density.
\end{proof}

\paragraph{What the result leaves open.}
A22 is a nonexistence result for a \emph{fixed finite}
choice of prime factors across all curves.  It does not
exclude a witness whose prime factors depend on the
curve, and it does not contradict A19's terminating
search for a specified finite support on one specified
curve.  It also does not establish the variable-prime
theta-incidence condition in A18.  An Artin-prime
density theorem continues to supply arithmetic
candidates, while an actual theta incidence at an
appropriate order remains a separate geometric
condition.


\subsection{An ordinary obstruction to the large-prime quotient test}
\label{sec:b3-ordinary-obstruction}

\paragraph{Finding A23 (the large-prime quotient test can fail over every extension).}
Let $C/\mathbf F_4$ be the smooth projective model of
\[
                   y^2+y=x+x^{-1}+(x-1)^{-1}.
 \tag{B3O.1}
\]
It is an ordinary genus-two curve.  Put
$Q=4^m$, $N_m=\#C(\mathbf F_Q)$ and
$J_m=\#J_C(\mathbf F_Q)$ for $m\ge1$.
For every prime $\ell\mid J_m$ one has
\[
                              \ell N_m<2J_m.
 \tag{B3O.2}
\]
Consequently the sufficient conditions in recovered A1
B.3, formula (B.CL.2), can never hold for a prime quotient
of this actual Jacobian over any constant extension of
$\mathbf F_4$, even if the primitive-root requirement is omitted.
Nevertheless the Raynaud theta support of this same curve has
nonzero points of exact order $4159$ over $\mathbf F_{4^7}$,
and $\langle2,-1\rangle=(\mathbf Z/4159)^\times$.
Thus (B3O.2) is a failure of a sufficient numerical test, not
an absence of geometric theta witnesses.

\begin{proof}
The equation is defined over $\mathbf F_2$.  Its three simple
poles at $0,1,\infty$ give genus two and $2$-rank two by the
Artin--Schreier formulas \cite[Lemma~2.6]{PriesZhu2012}.
There are three rational points over each of $\mathbf F_2$
and $\mathbf F_4$: the three ramified points account for all
$\mathbf F_2$-points; if $u^2+u+1=0$, the remaining
$\mathbf F_4$-values of $x$ are $u,u^2$, and the right side
of (B3O.1) has respective values $u^2,u$, both of absolute
trace one.  These have no Artin--Schreier preimages.
These are also the exact inputs (LW.12) of A13.

Let $\alpha_1,\ldots,\alpha_4$ be the Frobenius eigenvalues
over $\mathbf F_2$.  The two point counts imply
$\sum_i\alpha_i=0$ and $\sum_i\alpha_i^2=2$.
The genus-two functional equation therefore yields the
characteristic polynomial
\[
                         P_2(T)=T^4-T^2+4,
 \tag{B3O.3}
\]
as in (LW.13).  If $\beta_1,\beta_2$ are the roots of
$T^2-T+4$, the four squares $\alpha_i^2$ are
$\beta_1,\beta_1,\beta_2,\beta_2$.
Thus the Frobenius polynomial over $\mathbf F_4$ is
$(T^2-T+4)^2$; no isogeny classification is needed here.
Define
\[
 t_m=\beta_1^m+\beta_2^m,\qquad e_m=Q+1-t_m.
 \tag{B3O.4}
\]
The recurrence $t_0=2$, $t_1=1$ and
$t_m=t_{m-1}-4t_{m-2}$ shows that $t_m,e_m$ are integers.
The curve and Jacobian Frobenius point-count formulas give
\[
 \begin{aligned}
 N_m&=Q+1-2t_m=2e_m-(Q+1),\\
 J_m&=(1-\beta_1^m)^2(1-\beta_2^m)^2=e_m^2.
 \end{aligned}
 \tag{B3O.5}
\]
For the Jacobian formula see
\cite[Theorem~19.1(b)]{MilneAV}.
Since $|\beta_i|=2$, the Weil estimate gives
$e_m\ge(\sqrt Q-1)^2>0$.

If a prime $\ell$ divides $J_m=e_m^2$, then
$\ell\mid e_m$ and $\ell\le e_m$.
Also $N_m\ge0$, since it is a point count.  Hence
\[
 \ell N_m\le e_mN_m
       =2e_m^2-(Q+1)e_m
       <2e_m^2=2J_m,
 \tag{B3O.6}
\]
which proves (B3O.2) for every $m$ and every such prime.

The recovered source \cite[B.3, formula~(B.CL.2)]{A1run}
requires, for a surjection
$J_C(\mathbf F_Q)\twoheadrightarrow\mathbf Z/\ell$
with kernel $K$, both $N_m\ge2|K|$ and $\ell>N_m$,
in addition to $2$ being primitive modulo $\ell$.
Since $|K|=J_m/\ell$, its first inequality is exactly
$\ell N_m\ge2J_m$, contradicted by (B3O.2).
No primitive-divisor theorem or large-prime-factor theorem
can universally force these sufficient inequalities:
they are false for this ordinary curve over every extension
of its stated base field.

Finally A13 proves that the same model over $\mathbf F_2$
has at least four Raynaud-theta points of exact order
$4159$ over $\mathbf F_{128}=\mathbf F_{2^7}$, and that
$\langle2,-1\rangle=(\mathbf Z/4159)^\times$.
The inclusion
$\mathbf F_{2^7}\subset\mathbf F_{2^{14}}=\mathbf F_{4^7}$
preserves their order and theta incidence.  They therefore
remain valid geometric witnesses on the curve (B3O.1)
over its specified base $\mathbf F_4$.
The primitive-root condition in B.3 is stronger than this
signed-doubling condition: A13 has
$\operatorname{ord}_{4159}(2)=2079$, not $4158$.
The two arguments are consistent because B.3 gives only
one sufficient test and asserts no converse.
\end{proof}

\paragraph{Exact source locator for A23.}
The recovered B.3 is the complete claim and proof titled
\emph{The corrected pair estimate and its precise limitation}
in \path{ANABELIAN_A1_paper.tex}, lines 2488--2546 of the
retained checkpoint-03 text.  A13 supplies formulas
(LW.12)--(LW.17) and its complete theta-incidence proof.
The all-extension inequality (B3O.6) is symbolic; the
accompanying finite recurrence checks are arithmetic
cross-checks, not the proof of its universal quantifier.


\subsection{A Frobenius-orbit obstruction at inert prime powers}
\label{sec:frobenius-sidon}

\paragraph{Finding A24 (a necessary congruence for the A13 curve).}
Let $C/\mathbf F_2$ be the ordinary genus-two curve of A13,
\[
             y^2+y=x+x^{-1}+(x-1)^{-1},
\]
and let $X=C^{(1)}$, $J=\operatorname{Pic}^0(X)$ and $\Theta$
be its Raynaud theta support.  Let $n=\ell^r$, where $\ell$
is an odd prime and $r\ge1$, and assume the actual finite-level
condition
\[
                   \langle2,-1\rangle=(\mathbf Z/n)^\times.
                                                        \tag{SO.1}
\]
If $\Theta$ contains a point of exact order $n$, then $\ell$
is not inert in $\mathbf Q(\sqrt{-15})$.  Equivalently,
\[
        \ell\in\{3,5\}\quad\hbox{or}\quad
                    \ell\equiv1,2,4,8\pmod {15}.          \tag{SO.2}
\]
This is an unconditional geometric restriction.  In particular,
allowing higher powers satisfying (SO.1) does not restore an
inert underlying prime.  No assertion that every prime lifts to
all powers, and no general non-Wieferich hypothesis, is used.

The surviving arithmetic candidates are not exhausted by this
restriction.  Under premise M, the set of split primes satisfying
$\langle2,-1\rangle=(\mathbf Z/\ell)^\times$ has natural density
\[
                         \frac{72}{95}\mathcal A>0,        \tag{SO.3}
\]
where $\mathcal A$ is Artin's constant.  The smaller set of split
primes for which $2$ itself is primitive has density
$48\mathcal A/95$.  These conditional density statements supply
arithmetic candidates and do not assert theta incidence.

\begin{proof}[Proof of the unconditional restriction]
\emph{1. The finite-subgroup bound.}
For every nonzero $T\in J(\overline{\mathbf F}_2)$, A19 proves
\[
                   \#\bigl(\Theta\cap(\Theta+T)\bigr)\le2.
                                                        \tag{SO.4}
\]
Here and below these are numbers of geometric points.  If $H$ is
any finite subgroup of $J(\overline{\mathbf F}_2)$ and
$A\subseteq\Theta\cap H$, count the ordered pairs of distinct
points of $A$ by their nonzero difference.  Each difference has
at most two representations by (SO.4).  Thus
\[
                         |A|(|A|-1)\le2(|H|-1).           \tag{SO.5}
\]
In particular, for a prime-order point $z\in\Theta$ and its
Frobenius-generated $\mathbf F_\ell$-space $W$ of dimension $d$,
its Frobenius orbit $O$ satisfies
$|O|(|O|-1)\le2(\ell^d-1)$.  The same counting proof applies
to finite modules of prime-power exponent; it is not restricted
to vector spaces.

We will also use symmetry.  The characteristic-two locally exact
differential line bundle satisfies $B^{\otimes2}\simeq\omega_X$.
Riemann--Roch in degree one therefore gives
\[
 h^0(X,B\otimes L)=h^0(X,B\otimes L^{-1})
                        \qquad(L\in J),
\]
so $\Theta=-\Theta$.  This uses the same Raynaud line bundle as
in A03 and A11, not an arbitrary translate of the Abel curve.

\emph{2. The exact quadratic Frobenius relation.}
The equation and its Frobenius twist have the same model over
$\mathbf F_2$.  Let $\psi$ be $2$-power Frobenius and put
$\Phi=\psi^2$, the $4$-power Frobenius.  The checked A13 point
counts give the characteristic polynomial
\[
                    P_2(T)=T^4-T^2+4.
\]
Cayley--Hamilton on each odd Tate module gives the integral
identity
\[
                      \Phi^2-\Phi+4=0.                    \tag{SO.6}
\]
Thus the quadratic annihilator is established directly; it is
not inferred from the squared characteristic polynomial over
$\mathbf F_4$ without a semisimplicity argument.

Suppose henceforth that $\ell$ is inert.  Then
$Q(T)=T^2-T+4$ is irreducible modulo $\ell$, since its
discriminant is $-15$.  In particular $\ell\notin\{3,5\}$.
Let $\beta$ denote the class of $T$ in
$\mathbf F_{\ell^2}=\mathbf F_\ell[T]/(Q)$.  Its norm is
\[
                         \beta^{\ell+1}=4.                \tag{SO.7}
\]
On every nonzero vector of $J[\ell]$, the $\mathbf F_\ell[\Phi]$
span is a one-dimensional $\mathbf F_{\ell^2}$-space, hence a
two-dimensional $\mathbf F_\ell$-space, and the $\Phi$-orbit has
length
$s_1=\operatorname{ord}_{\mathbf F_{\ell^2}^{\times}}(\beta)$.

\emph{3. A uniform signed-orbit bound at the residue level.}
Write $h=(\ell-1)/2$.  Reduction of (SO.1) modulo $\ell$ gives
the signed-unit condition there.  In a cyclic group of even
order, adjoining $-1$ can enlarge a subgroup by at most two.
Consequently the order of $2$ is either $\ell-1$, or $h$ with
$h$ odd.  In either case
\[
                       \operatorname{ord}_\ell(4)=h.       \tag{SO.8}
\]
Let $H_1=\langle\beta,-1\rangle\subset\mathbf F_{\ell^2}^\times$.
Its norm image is $\langle4\rangle$, of order $h$, while the
norm kernel contains $-1$ and therefore has even order.
If this kernel had order two, then $|H_1|=2h=\ell-1$.
The unique subgroup of that order in the cyclic group
$\mathbf F_{\ell^2}^\times$ is $\mathbf F_\ell^\times$.
This would put $\beta$ in $\mathbf F_\ell$, contradicting
irreducibility.  The kernel consequently has order at least
four, and
\[
                          |H_1|\ge4h=2(\ell-1).            \tag{SO.9}
\]
This argument includes every inert eligible prime; it does not
require a large-prime cutoff or a list of exceptional cases.

\emph{4. Actual prime-power levels.}
Put $\mathcal O=\mathbf Z_\ell[T]/(Q)$, the ring of integers
of the unramified quadratic extension of $\mathbf Q_\ell$.
Identity (SO.6) makes $T_\ell J$ an $\mathcal O$-module.
The quadratic polynomial is irreducible over $\mathbf Q_\ell$,
so a nonzero element of $\mathcal O$ acts invertibly on
$T_\ell J\otimes\mathbf Q_\ell$.  The Tate module is therefore
torsion free over the discrete valuation ring $\mathcal O$,
and is free of rank two over it.  For
$R=\mathcal O/\ell^r\mathcal O$ this gives
\[
                              J[\ell^r]\simeq R^2.
\]
Let $z$ have exact additive order $\ell^r$.  In these coordinates
at least one component of $z$ is a unit of the local principal
ring $R$; otherwise its order would be at most $\ell^{r-1}$.
Thus $a\mapsto az$ is an injection $R\hookrightarrow R^2$.
The Frobenius-generated subgroup $H=Rz$ consequently has
exactly $\ell^{2r}$ points, and the $\Phi$-orbit of $z$ has
length $s_r=\operatorname{ord}_{R^\times}(\beta)$.

The reduction kernel $R^\times\to\mathbf F_{\ell^2}^\times$
is an $\ell$-group of exponent dividing $\ell^{r-1}$.
The latter assertion follows by repeated binomial expansion of
$1+\ell a$, with $\ell$ odd.  Therefore
\[
                           s_r=s_1\ell^b,
                         \qquad0\le b\le r-1.              \tag{SO.12}
\]
The norm $R^\times\to(\mathbf Z/\ell^r)^\times$ sends
$\beta$ to $4$.  The cyclic-unit argument used in (SO.8),
now applied to the actual condition (SO.1), gives
\[
                   \operatorname{ord}_{\ell^r}(4)
                         =h\ell^{r-1}.                    \tag{SO.13}
\]
This norm order divides $s_r$.  As $s_1$ is prime to $\ell$,
(SO.12)--(SO.13) force $b=r-1$.

Put $H_r=\langle\beta,-1\rangle\subset R^\times$.
If $s_1$ is odd, then $s_r=s_1\ell^{r-1}$ is odd, so adjoining
$-1$ doubles both orders.  If $s_1$ is even, then $-1$ already
belongs to both cyclic groups: the unique element of order two
in the units of this odd-residue-characteristic local ring is
$-1$.  To see uniqueness, $a^2=1$ implies
$(a-1)(a+1)=0$, and one of the factors is a unit because
$2$ is a unit; hence $a=1$ or $a=-1$.
It follows in both cases, with $k=\ell^{r-1}$, that
\[
                      |H_r|=k|H_1|\ge2(\ell-1)k.           \tag{SO.14}
\]
If $z\in\Theta$, Frobenius invariance and symmetry put its
entire signed orbit $H_rz$ in $\Theta\cap Rz$.  The injection
$R\hookrightarrow Rz$ makes its size exactly $|H_r|$.
Every inert prime here is at least seven, since $3,5$ ramify.
For $\ell\ge5$,
\[
\begin{split}
 |H_r|(|H_r|-1)
  &\ge k^2\,2(\ell-1)(2\ell-3)\\
  &>2\ell^2 k^2
   >2(\ell^{2r}-1).
\end{split}                                                \tag{SO.15}
\]
The middle inequality is
$2(\ell-1)(2\ell-3)-2\ell^2=2\ell^2-10\ell+6>0$.
The first follows from $k\ge1$ and
$(ak)(ak-1)\ge k^2a(a-1)$ for $a\ge1$.
This contradicts (SO.5), since $|Rz|=\ell^{2r}$.
This proves the stated
prime-power exclusion.  Quadratic reciprocity, or the two
quadratic characters of conductors $3$ and $5$, gives the split
classes $1,2,4,8$ modulo $15$, proving (SO.2).
\end{proof}

\paragraph{A finite consistency check, not a proof premise.}
Exact orders of $\beta$ in the three smallest inert residue fields
are
\[
\begin{array}{c|r|r|r}
 \ell&\operatorname{ord}(\beta)&s_1(s_1-1)&2(\ell^2-1)\\\hline
 7&24&552&96\\
 11&60&3540&240\\
 13&84&6972&336
\end{array}
\]
The accompanying \texttt{check\_A24\_orbits.py} verifies these
orders by quadratic-ring multiplication and an independent
$2$-by-$2$ multiplication-matrix implementation.  The uniform
norm-kernel proof above does not depend on this finite check.

\begin{proof}[Proof of the conditional residual densities]
Apply Lenstra's fixed finite Frobenius-condition theorem to
$W=\langle2,-1\rangle\subset\mathbf Q^\times$ and the
splitting condition in $F=\mathbf Q(\sqrt{-15})$.
The sufficient analytic family is
\[
 \zeta_{L_n}(s),\qquad
 L_n=\mathbf Q\bigl(\zeta_n,2^{1/n},(-1)^{1/n}\bigr),
                   \qquad n\ge1\text{ squarefree}.          \tag{SO.16}
\]
These Dedekind zeta functions are contained in premise M.
Lenstra's Section~2 permits an infinite finitely generated
subgroup with torsion, so it applies to this $W$ directly.
Theorem~3.1 and its Section~3.2 argument allow the fixed field
$F$ without requiring a separately enlarged GRH family
\cite[\S\S2--3, Theorem~3.1]{Lenstra1977}.

For odd squarefree $m$, write
$K_m=\mathbf Q(\zeta_m,2^{1/m})$.  As proved in A17,
\[
 [K_m:\mathbf Q]=m\varphi(m),\qquad
 K_m\cap\mathbf Q(\zeta_8)=\mathbf Q,
\]
and the maximal abelian subfield of $K_m$ is
$\mathbf Q(\zeta_m)$.  For completeness, Eisenstein at $2$
gives the degree, while the affine Galois group
$C_m\rtimes(\mathbf Z/m)^\times$ has commutator $C_m$:
the cyclotomic element $2$ conjugates a radical generator
$\tau$ to $\tau^2$.  The assertion about the abelian subfield
and the coprime-conductor intersection follow.

Taking odd roots of $-1$ adds nothing.  Taking a $2m$th root
of $-1$ adds $i$, and a $2m$th root of $2$ is generated by
$\sqrt2$ and $2^{1/m}$, using Bezout for $2,m$.  Thus
\[
 L_m=K_m,\qquad L_{2m}=\mathbf Q(\zeta_8)K_m,\qquad
 [L_m:\mathbf Q]=m\varphi(m),\quad
 [L_{2m}:\mathbf Q]=4m\varphi(m).                           \tag{SO.17}
\]
The maximal abelian subfields are respectively
$\mathbf Q(\zeta_m)$ and $\mathbf Q(\zeta_{8m})$.
The quadratic field $F$ has discriminant $-15$, is contained
in $\mathbf Q(\zeta_{15})$, and ramifies at both $3$ and $5$.
It follows that, for either $n=m$ or $n=2m$,
\[
 F\cap L_n=
 \begin{cases}
 F,&15\mid m,\\
 \mathbf Q,&15\nmid m.
 \end{cases}                                               \tag{SO.18}
\]
Indeed the ramification at a missing one of $3,5$ excludes
containment, while $15\mid m$ gives the containment directly.

For each finite collection of prime obstructions, ordinary
Chebotarev and inclusion-exclusion give the corresponding
finite sum of $\mu(n)/[FL_n:\mathbf Q]$.  Lenstra's theorem,
under (SO.16), identifies its limit with natural density.
The absolutely convergent sums can now be evaluated.  First
the total eligible density, without the splitting condition, is
\[
 D=\sum_{n\ {\rm squarefree}}\frac{\mu(n)}{[L_n:\mathbf Q]}
   =\frac34\prod_{q\ {\rm odd\ prime}}
                     \left(1-\frac1{q(q-1)}\right)
   =\frac32\mathcal A.
\]
Forcing $15\mid n$ changes the Euler factors at $3,5$ by
\[
 \frac{-1/6}{1-1/6}\cdot\frac{-1/20}{1-1/20}
                            =\frac1{95}.
\]
By (SO.18), the split density is therefore
\[
 \sum_n\frac{\mu(n)}{[FL_n:\mathbf Q]}
       =\frac12D+\frac12\frac{D}{95}
       =\frac{72}{95}\mathcal A.
\]
For $W=\langle2\rangle$ the even field is instead
$K_{2m}=K_m(\sqrt2)$ of degree $2m\varphi(m)$; the same
intersection argument applies.  The total density is then
$\mathcal A$, and its split part is $48\mathcal A/95$.
A sufficient GRH family for this smaller-set assertion is
$\zeta_{K_n}$ for squarefree $n$, as in A17.  Premise M
contains both stated sufficient families.
\end{proof}

\paragraph{Why the remaining candidates are not eliminated.}
The actual A13 witness $4159$ is split: $4159\equiv4\pmod{15}$.
Thus this same curve already prevents interpreting the
restriction as the absence of all eligible witnesses.
More generally, at a split prime away from the discriminant,
the quadratic Frobenius algebra is
$\mathbf F_\ell\times\mathbf F_\ell$.  For a point with
both eigencomponents nonzero, its generated space has dimension
two but its Frobenius orbit has at most $\ell-1$ points.
Such an orbit always satisfies the numerical inequality
$(\ell-1)(\ell-2)\le2(\ell^2-1)$, so this counting test
alone gives no contradiction.  An eigenline point can have a
shorter orbit; the A13 points over $\mathbf F_{128}$ have
$4$-power Frobenius orbit length dividing $7$.
The signed-unit condition controls the norm order; it does not
make theta invariant under multiplication by $2$.
Consequently A24 gives a necessary curve-dependent congruence
condition for prime and prime-power searches, while the
variable-prime theta-incidence gate of A18 remains open.


\subsection{A pairing obstruction to enlarging the rational torsion subgroup}
\label{sec:lang-pairing-rank-bound}

\paragraph{Finding A25 (a prime-level Lang test forces torsion rank at most two).}
Let $C/\mathbf F_Q$ be a genus-two curve in characteristic two,
let $J$ be its Jacobian, and put $G=J(\mathbf F_Q)$.
Let $r$ be an odd prime and
$\kappa=\dim_{\mathbf F_r}G[r]$.
Suppose either the nonzero-point numerical test (LW.2), or
the weaker ordinary-case positivity test following (LW.1),
holds for the subgroup $H$ under consideration.
Then the following implications hold:
\[
\begin{array}{c|c|c}
\text{hypothesis on }2\pmod r&H&\text{range forcing }\kappa\le2\\
\hline
\operatorname{ord}_r(2)=r-1&G[r]&r\ge11\\
\operatorname{ord}_r(2)=r-1&G[3r]&r\ge17\\
\langle2,-1\rangle=(\mathbf Z/r)^\times&G[r]&r\ge29\\
\langle2,-1\rangle=(\mathbf Z/r)^\times&G[3r]&r\ge41.
\end{array}
\tag{LP.1}
\]
In each range, writing $s=\sqrt Q$ and $J_Q=|G|$, the full
test (LW.2) therefore requires respectively
\[
 r>\left(J_Q\frac{2s+1}{(s-1)^2}\right)^{1/2}
 \quad\text{or}\quad
 r>\frac19\left(J_Q\frac{2s+1}{(s-1)^2}\right)^{1/2}.
\tag{LP.2}
\]
These thresholds have asymptotic sizes
$\sqrt2 Q^{3/4}$ and $(\sqrt2/9)Q^{3/4}$.
The corresponding ordinary thresholds replace $2s+1$ by $2s$.
These are necessary conditions for the specified sufficient
numerical tests, not for the existence of theta points.

\begin{proof}
The canonical principal polarization of the Jacobian gives a
nondegenerate alternating pairing
\[
                e_r:J[r]\times J[r]\longrightarrow\mu_r
\tag{LP.3}
\]
on its four-dimensional geometric $r$-torsion space, equivariant
for Frobenius.  For the pairing and its equivariance see
\cite[\S16, p.~29, and Lemma~16.2(e), p.~30]{MilneAV};
the principal polarization identifies the dual torsion space
with the original one.
If $\kappa\ge3$, the subspace $G[r]$ is not totally isotropic:
a totally isotropic subspace $W$ satisfies
$W\subseteq W^\perp$ and hence
$\dim W\le4-\dim W$.
Choose $P,T\in G[r]$ with $e_r(P,T)\ne1$.
Since $r$ is prime, this pairing value has exact order $r$.
Both points are fixed by Frobenius, so equivariance gives
\[
                         e_r(P,T)^Q=e_r(P,T).
\]
It follows that $r\mid Q-1$.
Writing $Q=2^v$ and $d=\operatorname{ord}_r(2)$, we obtain
\[
                         d\mid v,\qquad Q\ge2^d.
\tag{LP.4}
\]

The weaker of the two Lang tests in the statement already
implies
\[
 |H|>\frac{2sJ_Q}{(s-1)^2}
       \ge 2s(s-1)^2,
\tag{LP.5}
\]
where the last inequality uses $J_Q\ge(s-1)^4$.
For $Q\ge16$, the last expression is at least
$(9/8)Q^{3/2}$, because
$2(1-1/s)^2\ge2(3/4)^2=9/8$.
In particular, either successful test requires
$|H|>Q^{3/2}$ in this range.

If $H=G[r]$, then $|H|=r^\kappa\le r^4$.
If $H=G[3r]$, the ranges in (LP.1) exclude $r=3$, and
primary decomposition gives
$|H|=|G[3]|\,|G[r]|\le81r^4$.
If $\kappa\ge3$, (LP.4) now forces failure of (LP.5)
whenever
\[
                         c r^4<2^{3d/2},
\tag{LP.6}
\]
where $c=1$ or $81$ respectively.

For a primitive root, $d=r-1$.
The two starting inequalities are
\[
 11^4=14641<2^{15},\qquad
 81\cdot17^4=6765201<2^{24}.
\]
They persist on increasing the odd integer $r$ by two:
the left side grows by a factor at most
$(13/11)^4<2$, whereas the right side grows by $8$.
Thus (LP.6) holds in the first two ranges of (LP.1).

Under the signed-doubling hypothesis,
$r-1=|\langle2,-1\rangle|\le2d$, so
$d\ge(r-1)/2$.
For the other two ranges it suffices to check
\[
 29^4=707281<2^{21},\qquad
 81\cdot41^4=228886641<2^{30}.
\]
On increasing $r$ by two, the lower bound
$2^{3(r-1)/4}$ on the right of (LP.6) grows by
$2^{3/2}>2$, whereas the left side again grows by a
factor less than $2$.
All four ranges also imply $Q\ge16$ through (LP.4).
Consequently $\kappa\ge3$ contradicts either successful
Lang test, proving (LP.1).

Finally, solving (LW.2) for the subgroup size gives
$|H|>J_Q(2s+1)/(s-1)^2$.
With $\kappa\le2$, one has $|G[r]|\le r^2$ and
$|G[3r]|\le81r^2$, which proves (LP.2).
The asymptotics follow from $J_Q=Q^2+O(Q^{3/2})$.
\end{proof}

\paragraph{Implication for the proposed arithmetic route.}
In particular, first passing to an extension where all of
$J[r]$ is rational cannot force the prime-level Lang test
for an Artin prime $r\ge11$: it creates $\kappa=4$, the
excluded situation in (LP.1).
The factor $|G[3]|\le81$ can improve an individual index,
but the analogous strategy with $G[3r]$ fails for Artin
primes $r\ge17$.
This argument concerns these rational subgroup estimates.
It does not remove geometric theta witnesses, and does not
assert the same rank bound for truncated prime-power
subgroups $G[r^b]$ without an additional argument.


\subsection{A parity improvement for prescribed-point rational covers}
\label{sec:rank-one-augmentation}

\paragraph{Finding A26 (the rank-one augmentation criterion).}
Let $C/\mathbf F_Q$ be a smooth proper geometrically connected
genus-two curve in characteristic two, let $P\in C(\mathbf F_Q)$,
and suppose $N=\#C(\mathbf F_Q)$ is even.
Let $r$ be an odd prime with $\operatorname{ord}_r(2)=r-1$, and set
\[
 G=J_C(\mathbf F_Q),\qquad V=G/rG.
\tag{RA.1}
\]
If $|V|>N$, some surjection $G\twoheadrightarrow\mathbf Z/r$
gives a constant, geometrically connected, $P$-split cyclic
\'etale Lang cover whose relative class number is odd.
If $r\mid|G|$ and $r>N$, every such surjection has this property.
In both statements the prescribed point $P$ is arbitrary.
No lower bound on $N$ in terms of the kernel size is required.

Every genus-two curve of $2$-rank one satisfies the stated
parity hypothesis over its finite field of definition and
over every finite extension.  Thus both assertions apply
to every such curve whenever their remaining numerical
hypotheses hold.

\begin{proof}
We first prove the parity assertion.  The hyperelliptic
degree-two map in characteristic two is an Artin--Schreier
cover.  If its reduced geometric poles have odd orders
$d_1,\ldots,d_t$, its genus and $2$-rank are
\[
       g=\frac{\sum_{i=1}^t(d_i+1)-2}{2},
       \qquad f_C=t-1;
\]
see \cite[Lemma~2.6, p.~712]{PriesZhu2012}.
For $g=2$ and $f_C=1$, there are two poles, with
orders $1$ and $3$.  Frobenius preserves the ramification
orders, so it fixes both branch points individually.
Each has one geometric point above it, which is therefore
rational.  These are exactly the two fixed points of the
hyperelliptic involution.  Every other rational point is
paired with a distinct rational point by that involution.
Consequently $N$ is even.  The same argument applies over
every extension; in particular the two branch points remain
rational.

Now assume only the evenness of $N$.
In the odd-order group algebra $\mathbf F_2[V]$, put
\[
 \Theta_V=[0]+\sum_{R\in C(\mathbf F_Q)}[[R-P]\bmod rG].
\tag{RA.2}
\]
Its augmentation is $N+1=1$ in $\mathbf F_2$.
Recovered A1 B.2 gives the following exact criterion:
for a surjection $a:G\to\mathbf Z/r$, the associated
$P$-split Lang cover has odd relative class number if and
only if every nontrivial character of $\mathbf Z/r$ has
nonzero evaluation on $a(\Theta_V)$
\cite[B.2]{A1run}.
The connected Lang-pullback argument of A12, applied
to the Abel map based at $P$, proves geometric connectedness.
The zero point above $[P-P]=0$ gives the rational point
above the prescribed $P$, and the constant deck group
makes that fibre split.

Suppose all surjections to $\mathbf Z/r$ fail the criterion.
The irreducibility of $\Phi_r$ over $\mathbf F_2$, which
follows from $\operatorname{ord}_r(2)=r-1$, means that
one vanishing nontrivial character evaluation for such a
quotient forces all its nontrivial character evaluations
to vanish.  Every nontrivial character of $V$ factors
through one of these quotients.  Thus all nontrivial
Fourier coefficients of $\Theta_V$ vanish.
Fourier inversion is valid because $|V|$ is odd, and gives
\[
        \Theta_V=c\sum_{b\in V}[b],\qquad
        c=\frac{N+1}{|V|}=1\quad\text{in }\mathbf F_2.
\tag{RA.3}
\]
But $|V|>N$ leaves at least one empty Abel-image coset.
The zero coset is occupied by $P$, so an empty coset is
nonzero.  Its coefficient in (RA.2) is zero, contradicting
(RA.3).  Therefore some quotient succeeds.

For the assertion about every surjection, fix one
$a:G\to\mathbf Z/r$ and suppose it fails.
Apply the same argument to
$\Theta_a=[0]+\sum_R[a([R-P])]$ in
$\mathbf F_2[\mathbf Z/r]$.
Irreducibility and the augmentation force
$\Theta_a=\sum_{b\in\mathbf Z/r}[b]$.
Since $r>N$, an empty nonzero coset again contradicts
that equality.  As $a$ was arbitrary, every quotient
succeeds.  A surjection exists when $r\mid|G|$ by the
structure theorem for finite abelian groups.
\end{proof}

\paragraph{Relation to the recovered row statements.}
The Fourier identity (RA.3) is already the displayed
identity in recovered successor A.35, and the empty-coset
argument also occurs in A.10 \cite[A.10 and A.35]{ASrun}.
Those statements retain an additional pair-count
inequality; the present even-$N$ specialization obtains
an immediate contradiction from augmentation and does
not require that inequality.
Recovered A1 B.1 then gives the prescribed-point
separation implication for the resulting odd cyclic cover.
This is a rational relative-class-number condition over
the displayed field, not geometric primitive $2$-rank
zero; it need not survive an arbitrary constant extension.
The existence of a prime satisfying (RA.1) and $|V|>N$
is an additional requirement, not a consequence of
Hooley's positive density statement alone.


\subsection{A rank-one curve that defeats every augmentation size test}
\label{sec:rank-one-augmentation-obstruction}

\paragraph{Finding A27 (the augmentation size test can fail over every extension).}
Let $u\in\mathbf F_4$ satisfy $u^2+u+1=0$, and let
$C/\mathbf F_4$ be the smooth projective model of
\[
                    y^2+y=ux^3+x+u/x.
\tag{RO.1}
\]
This is a genus-two curve of $2$-rank one.
For every $m\ge1$, put $Q=4^m$, $G_m=J_C(\mathbf F_Q)$
and $N_m=\#C(\mathbf F_Q)$.
For every odd prime $r$ with $\operatorname{ord}_r(2)=r-1$,
one has
\[
                         |G_m/rG_m|<N_m.
\tag{RO.2}
\]
In particular, no prime and constant extension of this
curve can satisfy either sufficient numerical criterion
of A26.  This statement does not exclude rational
class-number witnesses that fail those sufficient
inequalities, or geometric theta witnesses.

\begin{proof}
\emph{1. Exact Frobenius data.}
The two Artin--Schreier poles of (RO.1) have odd orders
$3$ and $1$, so the genus and $2$-rank formulas give
$g=2$ and $f_C=1$ \cite[Lemma~2.6]{PriesZhu2012}.
The two ramified points are rational.
Direct finite-field enumeration gives
\[
                 N_1=6,\qquad N_2=28.
\tag{RO.3}
\]
The accompanying independent checker enumerates the
solutions in $\mathbf F_4$ and $\mathbf F_{16}$ directly,
including the two ramified points.  It embeds $u$ as
$z^2+z$ in $\mathbf F_{16}=\mathbf F_2[z]/(z^4+z+1)$.
The trace-zero method in the retained 72-model search
gives the same counts by a separate enumeration.

If $P_4(T)$ is the Frobenius characteristic polynomial,
(RO.3) gives $a_1=N_1-5=1$ and
$a_2=(a_1^2-(17-N_2))/2=6$.
The genus-two functional equation therefore gives
\[
 \begin{aligned}
 P_4(T)&=T^4+T^3+6T^2+4T+16\\
       &=(T^2-T+4)(T^2+2T+4)=A(T)B(T).
 \end{aligned}
\tag{RO.4}
\]
Only this exact polynomial identity is used below;
no elliptic isogeny classification is needed.

Let $t_m$ and $w_m$ be the sums of the $m$th powers of
the roots of $A$ and $B$, respectively.  They are integers,
with recurrences
\[
\begin{gathered}
 t_0=2,\quad t_1=1,\quad t_m=t_{m-1}-4t_{m-2},\\
 w_0=2,\quad w_1=-2,\quad w_m=-2w_{m-1}-4w_{m-2}.
\end{gathered}
\tag{RO.5}
\]
Put $e_m=Q+1-t_m$ and $h_m=Q+1-w_m$.
The Frobenius formulas \cite[Theorem~19.1(b)]{MilneAV} yield
\[
             |G_m|=e_mh_m,\qquad N_m=Q+1-t_m-w_m.
\tag{RO.6}
\]
All roots of $A$ and $B$ have complex absolute value $2$,
so
\[
 e_m\le Q+1+2\sqrt Q,
 \qquad N_m\ge L_m:=4^m+1-4\cdot2^m.
\tag{RO.7}
\]
Moreover $t_m\equiv1\pmod4$ for $m\ge1$ by (RO.5),
and hence
\[
                               4\mid e_m.
\tag{RO.8}
\]

\emph{2. The rational torsion rank for an eligible prime.}
Write $\kappa_m(r)=\dim_{\mathbf F_r}G_m[r]$.
Finite abelian-group structure gives
$|G_m/rG_m|=r^{\kappa_m(r)}$, and
$0\le\kappa_m(r)\le4$.
Let $F$ be the $4$-power Frobenius acting on $J_C[r]$.
Then $G_m[r]=\ker(F^m-1)$.

For $r\ge5$, the polynomials $A,B$ are coprime modulo $r$,
because their resultant is $36$: their difference is $3T$
and $A(0)=4$.
The characteristic polynomial (RO.4), Cayley--Hamilton and
Bezout therefore split the four-dimensional torsion space
as two $F$-stable, two-dimensional spaces
$W_A\oplus W_B$ on which $A(F)=0$ and $B(F)=0$,
respectively.  This uses no semisimplicity assumption.
Since
\[
                         (T-2)B(T)=T^3-8,
\]
one has $F^3=8$ on $W_B$.
If $F^m$ fixes a nonzero vector in $W_B$, it follows that
$2^{3m}\equiv1\pmod r$.
For an eligible prime this means
\[
                              r-1\mid3m.
\tag{RO.9}
\]
If $\kappa_m(r)\ge2$ and $W_B$ contributes no fixed
vector, all of $W_A$ is fixed by $F^m$.
Taking its determinant gives $4^m\equiv1\pmod r$,
so $r-1\mid2m$.
Consequently
\[
 r\ge5,\quad\kappa_m(r)\ge2
                  \quad\Longrightarrow\quad r\le3m+1.
\tag{RO.10}
\]

\emph{3. Every extension degree $m\ge10$.}
At $m=10$,
\[
       (3m+1)^4=923521<1044481=L_{10}.
\]
Also
$L_{m+1}-4L_m=8\cdot2^m-3>0$, while for $m\ge10$
\[
 \frac{(3(m+1)+1)^4}{(3m+1)^4}
          \le(34/31)^4<2.
\]
Induction gives
\[
                            (3m+1)^4<L_m\le N_m
                            \qquad(m\ge10).
\tag{RO.11}
\]
For $r=3$, $|G_m/3G_m|\le81<(3m+1)^4<N_m$.
For $r\ge5$ and $\kappa_m(r)\ge2$, (RO.10)--(RO.11)
prove (RO.2).  If $\kappa_m(r)=0$, then
$|G_m/rG_m|=1<N_m$.

It remains to consider $\kappa_m(r)=1$.
If $r\mid e_m$, (RO.8) and the fact that $r$ is odd give
\[
 r\le e_m/4\le\frac{Q+1+2\sqrt Q}{4}
                    <Q+1-4\sqrt Q\le N_m.
\tag{RO.12}
\]
The strict middle inequality holds for $\sqrt Q\ge8$:
after multiplying by four it is
$3((\sqrt Q)^2-6\sqrt Q+1)>0$.
If $r\nmid e_m$, then $r\mid h_m$ by (RO.6), and
$F^m$ has a fixed vector on $W_B$.
Equation (RO.9) and (RO.11) again give $r<N_m$.
Since $|G_m/rG_m|=r$ in this case, (RO.2) follows.

\emph{4. The nine remaining extension degrees.}
The exact recurrences (RO.5) and trial division give the
following complete table.  The middle column lists
\emph{all} eligible odd primes dividing $|G_m|$.
The last column is the maximum of
$r^{\min\{4,v_r(|G_m|)\}}$ over those primes, or $1$
when there are none.  It is an upper bound on
$|G_m/rG_m|$ for every eligible prime, including those
not dividing $|G_m|$.
\[
\begin{array}{c|r|l|r}
m&N_m&\text{eligible prime divisors}&\text{upper bound}\\
\hline
1&6&\varnothing&1\\
2&28&3&9\\
3&60&19&19\\
4&256&3,5,13&13\\
5&996&\varnothing&1\\
6&3976&3,19&81\\
7&16764&\varnothing&1\\
8&66016&3,5,13&13\\
9&260340&19,181&361.
\end{array}
\tag{RO.13}
\]
Every entry in the last column is strictly smaller than
$N_m$.  The independent checker retains the full
factorizations, prime certificates by trial division,
and multiplicative orders of $2$, as well as the
point enumerations and both induction anchors.
Thus (RO.2) holds for the remaining degrees and hence
for all $m\ge1$.
\end{proof}

\paragraph{Scope of the obstruction.}
The result refutes universal success of the numerical
augmentation criterion, even on $2$-rank-one curves.
The conclusion of A26 itself is a sufficient conditional
statement and remains valid.
No class-number calculation for every actual cyclic
cover is asserted here, and the numerical failure does
not imply absence of a rational witness or of a geometric
primitive-zero cover.


\subsection{The recovered canonical-class character criterion}
\label{sec:recovered-canonical-character}

\paragraph{Cited and rechecked result A28 (recovered successor A.26).}
Let $C/k$ be a smooth proper geometrically connected genus-two
curve over a finite field of characteristic $p$, let $P\in C(k)$,
and put
\[
             G=J_C(k),\qquad\kappa_P=[K_C-2P]\in G.
\tag{RC.1}
\]
Suppose $r\ne p$ is an odd prime and
$-1\in\langle p\rangle\subset(\mathbf Z/r)^\times$.
If $\kappa_P\notin rG$, there is a constant, geometrically
connected, $P$-split cyclic \'etale cover of degree $r$
over $k$ whose relative class number is prime to $p$.
More precisely, every surjection $a:G\to\mathbf Z/r$
with $a(\kappa_P)\ne0$ gives such a witness.
No size condition on $r$, $G/rG$ or $C(k)$ is required.

This is recovered \texttt{ANABELIAN\_SUCCESSOR}, final row
A.26, \emph{canonical-class sufficient condition}, not a
new prime-level theorem.  Its complete proof is retained
in the source, with the prescribed-point and character
hypotheses explicit \cite[A.26]{ASrun}.

The complete retained source proof was read in this run, including
its middle-coefficient Riemann--Roch symmetry, Frobenius
permutation of reduced characters, and faithful-character
contradiction.  The Lang-pullback connectedness and the
prescribed split fibre are the construction checked in A12.
The existing prime theorem is cited here rather than reproved.


\paragraph{The rank-one and Weierstrass interpretation.}
In characteristic two, if $W\in C(k)$ is a rational
Weierstrass point, then $K_C\sim2W$ and
\[
                           \kappa_P=2[W-P].
\tag{RC.2}
\]
Multiplication by two is invertible on $G/rG$ for odd $r$.
Thus the condition in A28 is exactly
$[W-P]\notin rG$.
For a prime $r$, inversion is a power of two if and
only if $\operatorname{ord}_r(2)$ is even.
For an arbitrary odd composite modulus, even order
alone is insufficient: $2$ has order four modulo $15$,
but its powers never equal $-1$ modulo $15$.
This distinction is retained in using higher conductors.

Any two rational Weierstrass points $W,W'$ satisfy
$2[W-W']=0$.  Every odd-order quotient kills their
difference, so the condition is independent of the
choice of $W$.
Rank-one genus-two curves have the two rational
ramification points proved in A26, so this interpretation
always applies to that stratum.
When $P$ is one of those points, $\kappa_P=0$ and
this criterion supplies no witness.
More generally, failure of its detecting-class
hypothesis does not assert that all actual covers fail.
Hooley supplies no missing prime factor of the fixed
finite-order class $\kappa_P$; extending constants
preserves that class's order.


% A29: canonical-class detection on a cyclic prime-power packet.
% Authored in work/; never edits final/.
\subsection{A prime-power canonical-class criterion at a prescribed point}
\label{sec:prime-power-canonical}

\paragraph{Finding A29 (detecting a nonzero primary canonical class).}
Let $C/k$ be a smooth projective geometrically connected genus-two
curve over a finite field $k=\mathbf F_q$ of characteristic $p$, let
$P\in C(k)$, and put
\[
 G=J_C(k),\qquad \kappa_P=[K_C-2P]\in G.
 \tag{PC.1}
\]
Suppose that an odd prime $r\ne p$ satisfies
\[
 \operatorname{ord}_r(p)\ \text{is even},\qquad
 (\kappa_P)_r\ne0,
 \tag{PC.2}
\]
where the subscript denotes the $r$-primary component in the finite
abelian group $G$. Then there are an integer $a\ge1$ and a constant,
geometrically connected cyclic etale $P$-split cover
$D\to C$ of degree $n=r^a$ over $k$ itself such that, for its unique
intermediate quotient $E\to C$ of degree $r^{a-1}$,
\[
 \begin{aligned}
 \mathcal P&=\ker(\operatorname{Nm}:J_D\longrightarrow J_E)^0,\\
 \#\mathcal P(k)&=\frac{\#J_D(k)}{\#J_E(k)},
 \qquad p\nmid\#\mathcal P(k).
 \end{aligned}
 \tag{PC.3}
\]
Consequently the covering formulation of A1 separates the point
section of $P$ from the point section of every distinct $Q\in C(k)$.
The statement only needs the rational $p$-primary group of this
primitive factor to vanish; its geometric $p$-rank need not be zero.
It does not assert that $\#J_D(k)/\#J_C(k)$ is prime to $p$ when
$a>1$.

\paragraph{Recovered inputs and the refinement.}
The prime-degree canonical-class criterion is already proved in
recovered successor \emph{final row A.26}, in every characteristic:
if $-1\in\langle p\rangle\subset(\mathbf Z/r)^\times$ and
$\kappa_P\notin rG$, a prime-degree rational witness exists
\cite[A.26]{ASrun}. Its proof uses the middle-coefficient
Riemann--Roch identity used below. Recovered A1 B.13--B.14 already
identify the primitive factor in a cyclic prime-power tower as the
appropriate rational norm obstruction \cite[B.13--B.14]{A1run}.
The refinement here combines these inputs to replace
$\kappa_P\notin rG$ by the strictly weaker finite-group condition
$(\kappa_P)_r\ne0$. We prove the whole combination, including
its pointed-cover consequence. This attribution does not make a
claim of bibliographic priority beyond the recovered source audit.

\begin{proof}
\emph{1. A cyclic quotient detecting the primary component.}
Decompose the $r$-primary subgroup as
$G_r\simeq\bigoplus_i\mathbf Z/r^{e_i}$.
Since $(\kappa_P)_r\ne0$, at least one of its coordinates is
nonzero. Projection to that coordinate, preceded by the canonical
primary projection of $G$, gives a surjection
\[
 b:G\twoheadrightarrow A=\mathbf Z/r^a,
 \qquad b(\kappa_P)\ne0,
 \tag{PC.4}
\]
with $a=e_i$. Conversely, any cyclic $r$-power quotient detecting
$\kappa_P$ implies $(\kappa_P)_r\ne0$.
In particular (PC.4) does not require $\kappa_P$ to survive in
$G/rG$: an element can be nonzero in a cyclic $r^a$-factor while
being divisible by $r$ in that factor.

Let $d_a=\operatorname{ord}_{r^a}(p)$. Reduction modulo $r$
shows that $\operatorname{ord}_r(p)$ divides $d_a$, so $d_a$ is
even. The group $(\mathbf Z/r^a)^\times$ is cyclic and its
unique element of order two is $-1$. Hence
\[
                  p^{d_a/2}\equiv-1\pmod{r^a}.
 \tag{PC.5}
\]
This uses the actual order at the chosen power. No assertion of
maximal order growth, or non-Wieferich hypothesis, is needed.

\emph{2. The constant pointed Lang cover.}
Write $\Phi$ for $q$-power Frobenius on $J=J_C$ and
$H=\ker b$. The map $\Phi-1$ is an etale isogeny with constant
kernel $J(k)=G$. It factors as
\[
 J\longrightarrow J/H\xrightarrow{\psi}J,
 \qquad \ker\psi=G/H=A.
\]
Pull $\psi$ back by $i_P:C\to J$, $R\mapsto[R-P]$, to obtain
$D=C\times_{J}(J/H)$. This is an etale cover of degree $n$ with
constant deck group $A$. Its fibre over $P$ is the constant kernel
of $\psi$, so it has a chosen rational point $P'$ above $P$ and
the whole fibre is rational.

For completeness, this pullback is geometrically connected.
Otherwise its geometric monodromy would be a proper subgroup
$A_0\subset A$. Quotienting by $A_0$ gives a nontrivial intermediate
etale isogeny $\psi':J'\to J$ whose pullback to $C$ has trivial
geometric monodromy. Over an algebraic closure, $i_P$ therefore
lifts to $v:C\to J'$. Subtract $v(P)$, which lies in $\ker\psi'$.
The Jacobian universal property gives a homomorphism $u:J\to J'$
with $u i_P=v-v(P)$ and $\psi'u=\operatorname{id}_J$.
The homomorphism $u$ is an isogeny, since its composition with
$\psi'$ is an isomorphism and both varieties have dimension two.
Multiplicativity of degree now contradicts $\deg\psi'>1$.
This is the connected Lang-pullback argument of A12.

\emph{3. The genus-two character polynomial and its symmetry.}
Let $\chi:A\to\mu_n$ be a nontrivial characteristic-zero
character. The unramified character Euler product for this cover is
\[
 L_P(\chi,T)
 =\sum_{d\ge0}\left(
   \sum_{Z\in\operatorname{Div}^{+}_d(C)(k)}
      \chi\bigl(b([Z-dP])\bigr)\right)T^d.
\]
The convention for Frobenius translation is the Lang convention
of recovered A1 B.2; reversing that convention replaces every
character by its inverse and leaves the argument unchanged.
In genus two its exact polynomial is
\[
 \begin{aligned}
 L_P(\chi,T)&=1+c_1(\chi)T
                  +q\chi\bigl(b(\kappa_P)\bigr)T^2,\\
 c_1(\chi)&=\sum_{R\in C(k)}\chi\bigl(b([R-P])\bigr).
 \end{aligned}
 \tag{PC.6}
\]
Here is a direct verification of the truncation. For $d>2$,
Riemann--Roch makes the number of effective divisors in each
degree-$d$ line-bundle class independent of the class. Summing
the nontrivial character over $\operatorname{Pic}^d(C)(k)$ gives
zero. In degree two, every line-bundle class has one effective
divisor except the canonical class, which has $q+1$. The excess
is $q\chi(b(\kappa_P))$. The degree-zero coefficient is one,
and degree one gives the displayed point sum. The rational point
$P$ identifies every degree component with $J(k)$ in these sums.

Riemann--Roch in degree one also gives
\[
 c_1(\chi)=\chi\bigl(b(\kappa_P)\bigr)c_1(\chi^{-1}).
 \tag{PC.7}
\]
Indeed $R\mapsto\iota R$, where $R+\iota R\sim K_C$, is a
bijection on $C(k)$. This follows directly because $K_C-R$ has
a unique effective divisor of degree one, defined over $k$.
The identity
$[\iota R-P]=\kappa_P-[R-P]$ proves (PC.7).
Thus no rational Weierstrass point is required.

\emph{4. Every faithful character is a unit at $p$.}
Reduce the cyclotomic integers modulo a fixed prime above $p$,
and denote the resulting character by $\bar\chi$. Since
$p\nmid n$, reduction preserves its order and its faithfulness.
Put $S_\chi=\overline{c_1(\chi)}$. If $j=d_a/2$, (PC.5)
and the integer coefficients of the point sum give
\[
 S_{\chi^{-1}}=S_\chi^{p^j},\qquad
 S_\chi=\bar\chi\bigl(b(\kappa_P)\bigr)S_\chi^{p^j}.
 \tag{PC.8}
\]
Suppose that $L_P(\chi,1)$ vanishes modulo that prime.
Equation (PC.6) gives $S_\chi=-1$, since $q=0$ in the residue
field. Its $p^j$th power is also $-1$, so (PC.8) forces
$\bar\chi(b(\kappa_P))=1$.
For a faithful $\chi$ this contradicts (PC.4). Therefore
\[
 L_P(\chi,1)\ \text{is a unit at }p
       \quad\text{for every character of exact order }r^a.
 \tag{PC.9}
\]
Nonfaithful characters may have zero reduction. The argument
deliberately makes no assertion about those factors.

\emph{5. The top-step Prym has prime-to-$p$ rational order.}
Choose a generator $\sigma$ of $A$ and set
$\tau=\sigma^{n/r}$, of order $r$; let $E=D/\langle\tau\rangle$.
Pullback and norm for $D\to E$ have composition $[r]$ on $J_E$.
They yield a $k$-isogeny
$J_E\times\mathcal P\to J_D$, so finite-field point counts give
the first equality in (PC.3).

For the exact character factorization, the Euler products of
$D/C$ and $E/C$ are products of the corresponding character
Euler factors. The characters for $E/C$ are exactly those of $A$
trivial on $\langle\tau\rangle$. Since $n$ is a prime power,
the remaining characters are exactly the faithful ones. Hence
\[
 \frac{\#J_D(k)}{\#J_E(k)}
       =\prod_{\substack{\chi\in\widehat A\\
                          \operatorname{ord}(\chi)=r^a}}
                       L_P(\chi,1).
 \tag{PC.10}
\]
This also follows by cancelling the imprimitive character factors
in the zeta-function ratios of recovered A1 B.2. The right side
is an integer, by its identification with $\#\mathcal P(k)$.
Each factor is a unit at the chosen prime above $p$ by (PC.9),
so the integer is prime to $p$. Thus
$\mathcal P(k)[p^\infty]=0$. The dimension check is
$\dim\mathcal P=(n-n/r)(g(C)-1)=\varphi(n)$.

\emph{6. The prime-power deck argument separates the point.}
Suppose that a distinct $Q\in C(k)$ collided with $P$ in the
covering formulation of A1. Applied to the pointed cover just
constructed, that formulation supplies $Q'\in D(k)$ above $Q$
with
\[
                  \delta=[Q'-P']\in J_D(k)[p^\infty].
 \tag{PC.11}
\]
This is the same collision-to-cover implication used in recovered
A1 B.1 and in A11; no converse to that implication is used.
The point $(\tau-1)\delta$ is rational, is $p$-primary, and lies in
\[
             \operatorname{im}(\tau-1)=\mathcal P.
 \tag{PC.12}
\]
To verify the equality, the image is connected and lies in the
norm kernel. Conversely, on $\mathcal P$ the endomorphism
$[r]-\sum_{i=0}^{r-1}\tau^i$ is $[r]$ and is divisible by
$\tau-1$ as a polynomial in $\tau$. Multiplication by $r$ is
surjective on $\mathcal P$, proving the reverse inclusion.
Equation (PC.3) therefore gives $(\tau-1)\delta=0$.

It follows that
$\tau Q'+P'-Q'-\tau P'$ is principal. Its four points are
distinct: the deck action is free, and the fibres over $P$ and
$Q$ are disjoint. Its positive and negative parts each have
degree two with simple support. The resulting rational function
gives a separable degree-two map $D_{\bar k}\to\mathbf P^1$;
the simple poles rule out a $p$th power if $p=2$, and the degree
itself excludes inseparability if $p\ne2$.

Since $g(D)=n+1\ge4$, its hyperelliptic pencil is unique.
The order-$r$ automorphism $\tau$ descends to an automorphism
of that projective line of order $r$: the kernel of the descent
action has order at most two. Over $\bar k$ this odd-order
automorphism has a fixed point. Its fibre contains at most two
geometric points, so an order-$r$ cyclic group acts trivially on
that fibre. This gives a fixed point of $\tau$ on $D$, contrary
to its free deck action. The assumed collision is impossible.
\end{proof}

\paragraph{The characteristic-two Weierstrass form and its exact residue.}
If $p=2$ and $W\in C(k)$ is a Weierstrass point, then
$K_C\sim2W$, so $\kappa_P=2[W-P]$. Multiplication by two is
an automorphism on every odd-primary subgroup. Consequently A29
applies precisely when $[W-P]$ has a nonzero $r$-primary component
for some odd prime $r$ with even $\operatorname{ord}_r(2)$.
Any two rational Weierstrass points differ by a class killed by
two, so this condition is independent of their choice.
It supplies a finite-field rational primitive-factor witness with
no size inequality and no analytic hypothesis.

This detector leaves unchanged those points for which the odd part
of $[W-P]$ is supported entirely on primes with odd
$\operatorname{ord}_r(2)$, including the case where $[W-P]$ is
purely $2$-primary. In particular it says nothing at a Weierstrass
point $P$, where $\kappa_P=0$. A constant extension cannot remove
this obstruction to the detector: the map on rational points of
the Jacobian is injective, so the finite order and primary support
of the fixed class $\kappa_P$ do not change. This is a limitation
of the sufficient construction, not a nonexistence theorem for
other cover witnesses. The A09 example already has non-Weierstrass
rational points whose differences from a rational Weierstrass
point are purely $2$-primary, because its rational Jacobian has
order eight.

\paragraph{Strictness on an actual curve.}
The difference from the prime-level canonical detector occurs on
the explicit curve of A09, without an abstract group realization
assumption. Write its point at infinity as $W$, and retain the
point $R=(u,v)\in C(\mathbf F_{16})$ whose class
$L=[R-W]$ has exact order $13$. For the prescribed point $R$,
\[
                   \kappa_R=-2L
 \tag{PC.13}
\]
also has exact order $13$. The finite etale surjective isogeny
$[13]:J_C\to J_C$ has a nonempty fibre over $\kappa_R$.
A point of this finite fibre is rational over some finite
extension $k'/\mathbf F_{16}$; write it as $U$, so that
$13U=\kappa_R$. Thus
$\kappa_R\in13J_C(k')$ while its $13$-primary component remains
nonzero. For any odd prime $r\ne13$, choose an integer $s$ with
$rs\equiv1\pmod{13}$. Then
$\kappa_R=r(s\kappa_R)$ in $J_C(k')$. Consequently
\[
                  \kappa_R\in rJ_C(k')
                    \quad\text{for every odd prime }r.
 \tag{PC.14}
\]
Every prime-level canonical detector in A28 therefore fails over
this field. Nevertheless $\operatorname{ord}_{13}(2)=12$, so
A29 constructs a cyclic $13^a$ pointed witness with the top-step
rational Prym of odd order. Any quotient detecting $\kappa_R$
must have $a\ge2$, because its image in every quotient of order
$13$ is zero. This proves strictness of the canonical detection
criteria on an actual curve. It does not assert that all actual
prime-degree covers fail: the independently constructed geometric
order-$13$ witness of A09 remains available. No explicit degree
for $k'$ is needed or claimed.

Recovered successor A.44 concerns a different pointed-involution
constraint at a Weierstrass base point and its primitive Frobenius
multipliers \cite[A.44]{ASrun}. It is compatible with the present
residual case. In general composite degree an even multiplicative
order alone does not force $-1\in\langle p\rangle$; the use of a
single odd prime power in (PC.5) is essential to the stated
hypothesis and to the simple order-$r$ deck endgame.


\subsection{The unramified double-cover repair and a genus-three obstruction}
\label{sec:genus-three-pointing-counterexample}

\paragraph{Finding A30 (the genus-two character truncation is essential).}
Over $k=\mathbf F_4$, let $D$ be the smooth projective model of
\[
                   z^2+z=\frac{1}{(t^2+t)^3}.
\tag{DC.1}
\]
It has genus three and $2$-rank one.  Let $P$ be the point
above $t=\infty$ with $z=0$, and let $W$ be the unique point
above $t=0$.  Both are rational; $W$ is Weierstrass and $P$
is not.  There is a surjection
\[
            b:J_D(k)\longrightarrow\mathbf Z/3,
                 \qquad b([W-P])\ne0,
\tag{DC.2}
\]
whose associated constant, geometrically connected, $P$-split
cyclic \'etale cover $D'\to D$ has \emph{even} relative class
number.  Since $\operatorname{ord}_3(2)=2$, the genus-two
conclusion of A28 that every such detecting quotient gives
an odd relative class number does not extend to genus three.
This does not rule out witnesses from other quotients or
higher powers for the same pointed genus-three curve.
The canonical class is also
detected: $K_D\sim4W$, so $b([K_D-4P])\ne0$.

More precisely, if $H_Y(T)$ denotes the numerator of the zeta
function of a curve $Y/k$, and
$B=(\ker(J_{D'}\to J_D))^0$, then
\[
 \overline{H_D}(T)=1-T,\qquad
 \overline{H_{D'}}(T)=(1-T)^2(1-T^3),\qquad
 \overline{H_B}(T)=(1-T)(1-T^3)
                         \quad\text{in }\mathbf F_2[T].
\tag{DC.3}
\]
The cover has degree three and $g(D')=7$.

\begin{proof}
\emph{1. The curve and its elliptic quotient.}
The rational function in (DC.1) has poles of order three
at $t=0,1$, and no others.  The Artin--Schreier genus and
$2$-rank formulas give $g(D)=3$ and $f_D=1$
\cite[Lemma~2.6, p.~712]{PriesZhu2012}.
Its hyperelliptic involution is $z\mapsto z+1$.
The two points above infinity have $z=0,1$ and are
unramified; the points above $0,1$ are its two ramification
points.  This proves the assertions about $P$ and $W$.
Since the hyperelliptic degree-two divisor class is $2W$,
the canonical class in genus three is $4W$.

Put
\[
 E/k:\ y^2+y=x^3,\qquad O=\infty,\qquad Q=(0,0).
\]
There is a separable degree-two map
\[
 f:D\longrightarrow E,\qquad
                 x=(t^2+t)^{-1},\quad y=z.
\tag{DC.4}
\]
Indeed its function-field extension is
$t^2+t=1/x$.  The two simple poles of $1/x$ on $E$ are
$Q=(0,0)$ and $-Q=(0,1)$, so these and only these are its
branch points.  Furthermore $f(P)=Q$ and $f(W)=O$.

There are nine points on $E(k)$: each of the four values
of $x$ has two values of $y$, and there is $O$.
Every affine point $(a,c)$ has order dividing three.
Its tangent line is $y=c+a^2(x-a)$; substituting
$x=a+v$ and using $a^4=a$ leaves $v^3=0$.
Thus the tangent meets the cubic three times at that point.
Consequently
\[
              E(k)\simeq(\mathbf Z/3)^2,
              \qquad \operatorname{ord}(Q)=3.
\tag{DC.5}
\]
The point count also gives
$H_E(T)=1+4T+4T^2$, in particular $\overline{H_E}=1$.

\emph{2. A translated Lang isogeny with the required fibres.}
Choose $a:E(k)\twoheadrightarrow\mathbf Z/3$ with $a(Q)=1$,
and write $H=\ker a$.
Let $F$ be $4$-power Frobenius on $E$.  Factoring the Lang
isogeny through $H$ gives
\[
 E\longrightarrow Y:=E/H\xrightarrow{\psi}E,
             \qquad \psi\pi_H=F-1.
\tag{DC.6}
\]
The map $\psi$ is an \'etale isogeny of degree three with
constant kernel $E(k)/H$.  The usual Lang-fibre computation
of A12 shows that a fibre $\psi^{-1}(R)$ is rational exactly
when $R\in H$; if rational, all three points are rational.
Use the translated cover
\[
                  \eta:Y\longrightarrow E,
                           V\longmapsto\psi(V)+Q.
\tag{DC.7}
\]
Its fibre above $Q$ has three rational points.  Its fibre
above $-Q$ has Frobenius translation
$a(-Q-Q)=a(Q)=1$, because $3Q=0$.
It is therefore one closed point of degree three.
The fibre above $O$ likewise has nontrivial translation
$a(-Q)=-1$.

Define $D'=D\times_{E,\eta}Y$.  This is a degree-three
constant \'etale cover of $D$, and its fibre above $P$ is
split.  It is geometrically connected: over an algebraic
closure the function-field extensions of $E$ defining $D$
and $Y$ have coprime degrees two and three, and the latter
is Galois.  They are linearly disjoint.
Unramified Riemann--Hurwitz gives $g(D')=3(3-1)+1=7$.

The rational quotient defining this pointed cover is
$b=a\circ f_*:J_D(k)\to\mathbf Z/3$.
Indeed $f_*[R-P]=[f(R)-Q]$, which is exactly the Lang
translation used in (DC.7).  In particular
\[
                 b([W-P])=a(O-Q)=-1\ne0.
\tag{DC.8}
\]
This also proves that $b$ is surjective.

\emph{3. The exact zeta calculation modulo two.}
For a separable degree-two cover $X\to Z$ of curves in
characteristic two, let $\mathcal R$ be its set of closed
branch points on $Z$.  Comparing Euler factors gives
\[
 Z_X(T)\equiv Z_Z(T^2)
                   \prod_{R\in\mathcal R}(1-T^{\deg R})
                                \pmod2.
\tag{DC.9}
\]
At an unramified closed point of degree $d$, either there
are two degree-$d$ points above it or one degree-$2d$ point.
Both local factors reduce to $(1-T^{2d})^{-1}$.
At a ramified point the local factor is $(1-T^d)^{-1}$;
multiplication by $1-T^d$ changes the corresponding
$Z_Z(T^2)$ factor to this one.  This proves (DC.9) as an
identity of formal power series.

Since the constant field has even cardinality, passing to
zeta numerators in (DC.9) gives the polynomial identity
\[
 \overline{H_X}(T)=\overline{H_Z}(T^2)
             \frac{\prod_{R\in\mathcal R}(1-T^{\deg R})}{1-T}.
\tag{DC.10}
\]
For $D\to E$, the two branch points are rational and
$\overline{H_E}=1$, giving $\overline{H_D}=1-T$.
The elliptic curve $Y$ is isogenous to $E$, so
$H_Y=H_E$ and $\overline{H_Y}=1$.
For $D'\to Y$, the branch divisor is the \'etale pullback
of $Q+(-Q)$.  Its closed-point degrees are $1,1,1,3$ by
Step~2.  Formula (DC.10) therefore gives
\[
       \overline{H_{D'}}(T)
              =\frac{(1-T)^3(1-T^3)}{1-T}
              =(1-T)^2(1-T^3).
\]
Pullback and norm for the degree-three cover show that
$J_D\times B$ is isogenous to $J_{D'}$.  Dividing their
Frobenius polynomials proves the last identity in (DC.3).
Finally
$\#B(k)=H_B(1)\equiv0\pmod2$.
The point-count and isogeny identities are the usual
Frobenius identities, as in
\cite[Theorem~19.1(b) and its Frobenius proof, pp.~40--41]{MilneAV}.

\emph{4. The individual reduced character polynomials.}
The failure is visible before taking their product.
Let $\chi:\mathbf Z/3\to\overline{\mathbf F}_2^\times$
send $1$ to $\zeta$, a primitive cube root.
The nontrivial unramified character sheaf on $E$ defined
by (DC.7) has $L$-polynomial $1$.
For example, Riemann--Roch on the elliptic curve makes
the number of effective divisors in a rational divisor
class of each positive degree independent of the class,
so every positive-degree nontrivial character coefficient
vanishes.

The character Euler-factor argument of Step~3 applies
with a character value $\xi$ at a closed point of $E$:
both unramified splitting types give
$(1-\xi^2T^{2d})^{-1}$ modulo two, while a branch point
contributes the correction $1-\xi T^d$.
The character values at $Q,-Q$ are $1,\zeta$ respectively.
Thus the degree-four unramified character polynomial on
the genus-three curve $D$ has reduction
\[
       \overline{L_{\chi,D,P}}(T)=(1-T)(1-\zeta T),
       \qquad
       \overline{L_{\chi^{-1},D,P}}(T)
                                      =(1-T)(1-\zeta^2T).
\tag{DC.11}
\]
Both values at $T=1$ vanish.  The genus-two reduction is
linear, whereas (DC.11) has a genuine quadratic term.
This is precisely why its inversion symmetry does not
give the genus-two contradiction.
\end{proof}

\paragraph{What happens to the proposed unramified double-cover repair.}
The underlying geometric observation is correct.
For a rank-one genus-two Artin--Schreier curve written,
after choosing the two rational branch points as $0,\infty$,
as
\[
 C:\ y^2+y=f_1(x)+f_3(x),
\]
where $f_1$ has just the simple pole at $0$ and $f_3$
just the triple pole at $\infty$, form the connected
biquadratic fibre product
\[
 z_1^2+z_1=f_1(x),\qquad z_2^2+z_2=f_3(x),
                         \qquad y=z_1+z_2.
\tag{DC.12}
\]
Over an algebraic closure the two Artin--Schreier classes
are independent, since their ramification supports differ.
The inertia groups over $0,\infty$ are respectively the
$z_1$- and $z_2$-translation groups.
The diagonal translation is in neither inertia group,
so its quotient map $\widetilde C\to C$ is \'etale of
degree two.  It is the unique nontrivial geometric etale
double cover: $f_C=1$ means that
$H^1_{\mathrm{et}}(C_{\bar k},\mathbf F_2)$ has dimension one.
Its source has genus three by Riemann--Hurwitz.
The quotient by $z_2$-translation is
the genus-zero $z_1$-curve.  Hence $z_2$-translation is
the hyperelliptic involution of $\widetilde C$.

At the simple pole its two lifts are interchanged by
this involution, so neither is Weierstrass.  At the
triple pole both lifts are fixed and are Weierstrass.
Equivalently, the hyperelliptic presentation on the
genus-zero $z_1$-curve has two poles of order three,
so the same formula gives $2$-rank one.
Finite constant extensions, or the appropriate constant
twist, can make a prescribed fibre rational, but do not
change these geometric conclusions.

This construction nevertheless does not repair the A1
argument.  First, its nontrivial geometric degree-two
quotient does not factor through the geometric
prime-to-two fundamental group, whose finite quotients
have odd order.  A collision in the quotient used by A1
therefore does not supply lifts to this cover.
Second, if $P_s,P_t$ are the simple- and triple-pole
points on $C$, their difference has exact order two:
$\operatorname{div}(x)=2P_s-2P_t$, and a nonzero
degree-one principal divisor would make $C$ rational.
For lifts $\widetilde P_s,\widetilde P_t$ one has
\[
 \pi_*[\widetilde P_t-\widetilde P_s]=[P_t-P_s].
\tag{DC.13}
\]
Every odd cyclic character pulled back from $C$ thus
kills this difference.  A new odd character detecting
it upstairs cannot supply the A28 detector by descent
from such a character on $C$.

The explicit curve (DC.1) is exactly this construction
for $C:y^2+y=x^3+x^{-1}$, with $z_1=t$ and $z_2=z$.
Its example (DC.11) gives an additional obstruction:
even as an argument about the upstairs curve alone,
the proposed unrestricted higher-genus symmetry rule
is false.


\subsection{An exact theta witness on the numerical-obstruction curve}
\label{sec:a27-theta-seven}

\paragraph{Finding A31 (a degree-seven geometric witness on A27).}
Let $k=\mathbf F_4=\mathbf F_2(u)$, with $u^2+u+1=0$, and let
\[
 C/k:\qquad y^2+y=u t^3+t+u/t
 \tag{TS.1}
\]
be the smooth proper rank-one genus-two curve of A27.  Let $W$ be
its point over $t=0$ and let $R=(u,u^2)$.  Then
$M=\mathcal O_C(R-W)$ has exact order seven.  On the Frobenius
twist $C^{(1)}$, its twist $M^{(1)}$ lies on the Raynaud theta
divisor.  Consequently A11 gives a geometrically connected cyclic
degree-seven cover with geometric primitive $2$-rank zero.  After
a finite constant extension it gives a constant split witness at
any prescribed point of $C$.

For the prescribed point $R$, the canonical class
$\kappa_R=[K_C-2R]$ has exact order seven.  Thus A29's detector
does not apply to this class, even after extending constants,
because $\operatorname{ord}_7(2)=3$ is odd.  This is a concrete
instance of the remaining detector case that the geometric method
does handle.  It is not a universal incidence theorem.  It is also
compatible with every numerical failure proved in A27.

\begin{proof}
\emph{1. An explicit isomorphic model.}
The coordinate change
\[
 \begin{aligned}
 X&=u/t, &
 Y&=u^2t^{-2}(y+u)+u^2,\\
 t&=u/X, &
 y&=u+(Y+u^2)/X^2
 \end{aligned}
 \tag{TS.2}
\]
and its displayed inverse identify the function field of (TS.1)
with that of
\[
 \begin{aligned}
 C':\quad Y^2+X^2Y&=f(X),\\
 f(X)&=X^5+X^4+uX^3+u^2X^2+uX+u.
 \end{aligned}
 \tag{TS.3}
\]
Here is a coefficient check, retaining all constant terms.
Put $\alpha=u^2t^{-2}$ and $\beta=t^{-2}+u^2$, so
$Y=\alpha y+\beta$.  Since $X^2\alpha=\alpha^2$,
the equation on the left becomes
\[
 \begin{aligned}
 Y^2+X^2Y
 &=\alpha^2(ut^3+t+u/t)+\beta^2+X^2\beta\\
 &=u^2t^{-5}+ut^{-4}+ut^{-3}
       +ut^{-2}+u^2t^{-1}+u\\
 &=f(u/t).
 \end{aligned}
 \tag{TS.4}
\]
Both are models of their smooth proper function-field curves, so
this birational identification extends to an isomorphism.
The two odd Artin--Schreier pole orders in (TS.1) are three and
one; thus the genus and $2$-rank are those already checked in A27.
The point $W$ becomes the unique point at infinity in (TS.3),
with $\operatorname{ord}_W(X)=-2$ and
$\operatorname{ord}_W(Y)=-5$.  The point $R$ becomes $(1,u)$.
The affine model is smooth: the only possible singular
$X$-coordinate is zero, and $f'(0)=u\ne0$ there.

\emph{2. A principal divisor of degree seven.}
Set
\[
                F=(X+u)Y+X.
 \tag{TS.5}
\]
The hyperelliptic involution sends $Y$ to $Y+X^2$.
Taking its function-field norm gives the exact polynomial identity
\[
 \begin{aligned}
 F\,\iota(F)
 &=X^2+X^3(X+u)+(X+u)^2f(X)\\
 &=(X+1)^7.
 \end{aligned}
 \tag{TS.6}
\]
Indeed expansion yields coefficient one in every degree from
zero through seven.  At the two points above $X=1$,
\[
       F(1,u)=0,\qquad F(1,u^2)=u^2\ne0.
 \tag{TS.7}
\]
They are distinct unramified points since $X^2=1$ there.
Thus (TS.6) assigns all seven zeros of $F$ to $R$, with no
zero at its conjugate.  The function is regular on the affine
curve and has a pole of exact order seven at $W$:
the term $XY$ has pole order seven, whereas $uY$ and $X$
have pole orders five and two.  Consequently
\[
                       \operatorname{div}(F)=7(R-W).
 \tag{TS.8}
\]
The class $[R-W]$ is nonzero.  Otherwise a function with divisor
$R-W$ would give a degree-one map from a genus-two curve to
$\mathbf P^1$, which is impossible.  Since seven is prime,
(TS.8) proves that its exact order is seven.

\emph{3. The theta bundle and the Frobenius twist.}
The differential
\[
                    \omega=dX/X^2=d(1/X)
 \tag{TS.9}
\]
is exact and has divisor $2W$.  At a finite point with $X\ne0$,
$X$ is an etale local coordinate, so $\omega$ is nonvanishing.
At $X=0$, the differentiated equation gives
$\omega=dY/f'(X)$; since $f'(0)=u\ne0$, this is again a
nonvanishing regular differential.  Its divisor is therefore
supported at $W$ and has degree $2g-2=2$.

Let $B_C$ denote the line of locally exact differentials on
$C^{(1)}$.  The section supplied by (TS.9) has divisor
$W^{(1)}$ on that twist: a local zero of order two in an exact
differential corresponds to a zero of order one in the
$\mathcal O_{C^{(1)}}$-line.  Hence
\[
 B_C\simeq\mathcal O_{C^{(1)}}(W^{(1)}),\qquad
 B_C\otimes M^{(1)}
       \simeq\mathcal O_{C^{(1)}}(R^{(1)}).
 \tag{TS.10}
\]
The second line bundle has a nonzero section, proving the asserted
theta incidence.  This keeps the twist explicit over $\mathbf F_4$:
the coefficients of (TS.3) are squared, and
$R^{(1)}=(1,u^2)$ in the resulting coordinates.
Twisting (TS.8) preserves the exact order seven.

\emph{4. The cover and the exact limitation of the detectors.}
The powers of two modulo seven are $1,2,4$.  Adjoining their
negatives gives every nonzero residue, so
\[
                  \langle2,-1\rangle=(\mathbf Z/7)^\times.
 \tag{TS.11}
\]
Finding A11 applies to (TS.10) and gives a cyclic degree-seven
cover with primitive geometric $2$-rank zero.  Since the degree
is prime, this is the whole new factor.  Descending the cover and
its deck transformations and rationalizing a point above any
prescribed base point require only a finite constant extension.
The geometric rank-zero property survives these operations; the
relative rational class number is then odd, and A11's pointed
deck argument separates any putative collision.

Finally (TS.9) also gives $K_C\sim2W$, so
$\kappa_R=-2[R-W]$ has exact order seven.  Its only nonzero odd
primary part is the one for seven, whose order of two is odd.
Injectivity of rational Jacobian points under constant extension
preserves that order and prime support.  Thus this prescribed
point remains outside A29's sufficient hypothesis.  In fact A27
gives $\#J_C(\mathbf F_4)=28$ and $\#C(\mathbf F_4)=6$, so for
every one of those six prescribed points the odd primary support
of $\kappa_P$ is contained in $\{7\}$.  A29's detector therefore
fails for all six points, and continues to fail for these same
points after every constant extension, whereas the geometric
cover just constructed becomes split at each prescribed point
after a finite extension.  A27 refuted
universal success of its stated size inequalities at primes with
two primitive; seven does not satisfy that primitive-root
hypothesis.  Neither A27 nor the failure of A29's detector
asserted absence of the present geometric witness.
\end{proof}

The model (TS.3), norm (TS.6) and isomorphism (TS.2) were found
in this run by the root investigator.  The compact certificate
\texttt{verify\_A31\_norm7.py} checks their coefficients directly
over $\mathbf F_4$, including the inverse substitution, both point
values, smoothness at $X=0$ and the signed unit orbit.  It uses no
Jacobian arithmetic package.  The divisor, differential, exact-order
and covering arguments are the written proof above, not an
inference from a numerical search alone.


\section{Effective Galois comparison: exact analytic inputs and constants}
\label{sec:effective-galois}

Throughout this section, $D_L=|\operatorname{disc}(L)|$, $n_L=[L:\mathbb Q]$, and all logarithms are natural.  Write $\mathrm{H}_M$ for the draft assertion that GRH holds for the Hecke $L$-functions of every number field.  Every conclusion using $\mathrm{H}_M$ remains conditional on that assertion.  The trivial Hecke character gives GRH for $\zeta_L$ for every number field $L$.

\paragraph{Source coverage.}
The six supplied \texttt{SORTED\_*.md} files are summaries, with frequently truncated statements.  The later recovery of the complete A1, successor, and LM8 papers permits the row-level scope to be checked against actual proofs.  In particular, LM8 A.32 uses Serre's qualitative open-image theorem for a non-CM elliptic curve only to obtain Zariski density of the commutator image; no effective index is needed there.  LM8 D.25 uses Faltings's finiteness theorem for rational points, not the Faltings--Serre trace-comparison method.  The bounds below are explicit available refinements of finite-quotient arguments, not a claim that these two qualitative inputs were conditional.  The Gaussian and Fermat-cubic CM elliptic curves in other LM8 rows do not satisfy the non-CM hypothesis of the open-image bound below.  The row ledger records the precise finite fields in the recovered proofs.

\subsection{The finite-cover Chebotarev bound}

\paragraph{Cited theorem G-CHEB (Bach--Sorenson).}
If $L/K$ is finite Galois, $L\ne\mathbb Q$, and $C$ is a conjugacy class in $\operatorname{Gal}(L/K)$, then GRH for $\zeta_L$ and $\zeta_{\mathbb Q}$ gives an unramified degree-one prime $\mathfrak p$ of $K$ with Frobenius class $C$ and
\[
 N\mathfrak p\le B(L):=(4\log D_L+\tfrac52 n_L+5)^2.
\]
The exact reference is Bach--Sorenson, \emph{Explicit bounds for primes in residue classes}, Math. Comp. 65 (1996), 1717--1735, Theorem 5.1, p.~1729; the analytic hypotheses are specified in the last paragraph of \S2, p.~1720.  Their Theorem 3.2 and Corollary 3.3, p.~1721, make the reduction to a cyclic extension.  No Artin holomorphy conjecture is assumed.  This is already a conditional theorem of the literature; it is not a new finding of this work.

\paragraph{Finding G01 (a bounded prime outside a prescribed set).}
Let $L/\mathbb Q$ be finite Galois and let $S$ be a finite set of rational primes. Put
\[
 m=\operatorname{rad}\!\left(2\prod_{p\in S}p\right),\qquad
 \widetilde L=L(\sqrt m).
\]
For any nonempty conjugation-stable subset $C\subseteq\operatorname{Gal}(L/\mathbb Q)$, GRH for $\zeta_{\widetilde L}$ and $\zeta_{\mathbb Q}$ gives a rational prime $p\notin S\cup\{2\}$, unramified in $L$, such that
\[
 \operatorname{Frob}_p(L/\mathbb Q)\subseteq C,
 \qquad p\le B(\widetilde L).
\]
Thus $\mathrm{H}_M$ supplies this statement for arbitrary $L$. The stipulated Dirichlet GRH supplies it when $\widetilde L/\mathbb Q$ is abelian, in particular when $L/\mathbb Q$ is abelian.

\begin{proof}
The integer $m$ is even and squarefree, and $m\ge2$, so every prime dividing $m$ ramifies in the quadratic field $\mathbb Q(\sqrt m)$.  Therefore each such prime ramifies in $\widetilde L$.  Restriction $\operatorname{Gal}(\widetilde L/\mathbb Q)\to\operatorname{Gal}(L/\mathbb Q)$ is surjective.  The inverse image of $C$ is nonempty and stable under conjugation; choose one conjugacy class in it and apply G-CHEB.  The resulting prime is unramified in $\widetilde L$, hence does not divide $m$, and restriction of Frobenius proves the claim.  This argument is the squarefree-prime-avoidance construction of Mayle--Wang, Corollary~6, with $2$ explicitly included to remove a possible exceptional prime.
\end{proof}

\paragraph{Finding G02 (finite quotient trace comparison).}
Let $\rho_i:G_{\mathbb Q,S}\to\operatorname{GL}_d(R)$, $i=1,2$, be continuous representations over a finite commutative ring $R$, and let $L$ be the fixed field of the joint kernel. If their trace functions differ, they differ at a prime $p\notin S\cup\{2\}$ satisfying $p\le B(L(\sqrt m))$, under the GRH assumption in G01. The same assertion holds for any conjugation-invariant function of the joint image, including coefficients of characteristic polynomials.

\begin{proof}
The joint image is $\operatorname{Gal}(L/\mathbb Q)$.  The set on which the two trace functions differ is nonempty and stable under conjugation.  G01 applies to it.  A conjugation-invariant function is handled by precisely the same argument.
\end{proof}

\paragraph{Arbitrary number-field bases.}
For a finite Galois extension $L/K$, let $M/\mathbb Q$ be its normal closure and let $S_0$ contain the rational primes below every prime to be excluded.  With $m=\operatorname{rad}(2\prod_{p\in S_0}p)$, any prescribed class in $\operatorname{Gal}(L/K)$ occurs at a degree-one prime $\mathfrak p$ of $K$, above a rational prime outside $S_0\cup\{2\}$, with
\[
 N\mathfrak p\le B(M(\sqrt m))
\]
under GRH for $\zeta_{M(\sqrt m)}$ and $\zeta_{\mathbb Q}$.  To see this, extend a representative $g\in\operatorname{Gal}(L/K)$ to $\gamma\in\operatorname{Gal}(M/K)$ and apply G01 to its conjugacy class in $\operatorname{Gal}(M/\mathbb Q)$.  Choose a prime of $M$ whose Frobenius is precisely $\gamma$.  Since $\gamma$ fixes $K$, the prime below it in $K$ has residue degree one over $\mathbb Q$; its Frobenius in $L/K$ is $g$.  This gives an explicit form for the arbitrary-base finite affine images in the ledger, after their finite extension is fixed.  It does not bound a preceding Hilbert specialization or a depth parameter not specified by the proof.

This test compares trace functions at the chosen finite level.  Equality modulo a fixed power of a coefficient prime does not imply equality of the full adic representations.  To obtain a finite test for their full trace functions, one needs the next construction.

\subsection{A complete finite test for adic semisimplifications}

\paragraph{Finding G03 (an explicit uniform trace bound).}
Let $\mathcal O$ be the ring of integers in a finite extension $F/\mathbb Q_\ell$, with uniformizer $\varpi$ and residue field of cardinality $q$.  Let
\[
 \rho_1,\rho_2:G_{\mathbb Q,S}\longrightarrow\operatorname{GL}_d(\mathcal O)
\]
be continuous.  Set
\[
 N=2(q^{2d^2}-1),\qquad
 m=\operatorname{rad}\!\left(2\prod_{p\in S}p\right),
\]
and
\[
 \mathcal B(d,q,S)=
 \left[4\bigl((N-1)\log m+N\log N\bigr)+\tfrac52 N+5\right]^2.
 \tag{G03.1}
\]
Conditional on $\mathrm{H}_M$, equality of the exact Frobenius traces for every prime $p\nmid m$ with $p\le\mathcal B(d,q,S)$ is equivalent to equality of their trace functions on $G_{\mathbb Q,S}$.  It is therefore equivalent to isomorphism of their semisimplifications over $F$.  In particular it is equivalent to isomorphism when both representations are semisimple.

\begin{proof}
Suppose the trace functions differ.  Let $M$ be the $\mathcal O$-linear span of
\[
 \{(\rho_1(g),\rho_2(g)):g\in G_{\mathbb Q,S}\}
 \subseteq M_d(\mathcal O)\oplus M_d(\mathcal O).
\]
Products of spanning elements are again spanning elements, so $M$ is an $\mathcal O$-algebra with identity. It is a finite free $\mathcal O$-module of rank $r\le2d^2$, because it is a submodule of the finite free module on the right and $\mathcal O$ is a discrete valuation ring.  Reduction in this algebra gives a continuous finite quotient
\[
 \delta:G_{\mathbb Q,S}\longrightarrow H\subseteq(M/\varpi M)^\times.
\]
Continuity follows because $M$ is a finite module with its subspace topology and $\varpi M$ is open in $M$.  The quotient ring has $q^r$ elements; its zero element is not a unit. Hence $|H|\le q^r-1\le q^{2d^2}-1$.

Define the $\mathcal O$-linear functional
\[
 t:M\longrightarrow\mathcal O,\qquad t(A,B)=\operatorname{tr}A-\operatorname{tr}B.
\]
Its image is a nonzero ideal $\varpi^a\mathcal O$. Therefore $\lambda=\varpi^{-a}t$ is integral on $M$ and has nonzero reduction $\overline\lambda$ on $M/\varpi M$. Since the pairs $(\rho_1(g),\rho_2(g))$ span $M$, at least one has nonzero $\overline\lambda$-value. Conjugation by every such pair preserves $t$, hence also $\lambda$. Consequently
\[
 C=\{h\in H:\overline\lambda(h)\ne0\}
\]
is well defined, nonempty, and stable under conjugation. Every $g$ mapping to $C$ has unequal traces.

Let $L/\mathbb Q$ be the finite Galois extension defined by $\ker\delta$. It is unramified outside $S$, and $n_L=|H|\le q^{2d^2}-1$. Thus $\widetilde L=L(\sqrt m)$ has degree $\widetilde n\le N$ and is ramified only at primes dividing $m$. By the discriminant estimate of Serre, \emph{Quelques applications du th\'eor\`eme de densit\'e de Chebotarev}, Proposition~6 (also Mayle--Wang, Lemma~7),
\[
 \log D_{\widetilde L}
 \le(\widetilde n-1)\log\operatorname{rad}(D_{\widetilde L})
       +\widetilde n\log\widetilde n
 \le(N-1)\log m+N\log N.
\]
G01 supplies $p\nmid m$ of unequal traces with $p\le B(\widetilde L)\le\mathcal B(d,q,S)$. This proves the contrapositive of the finite-test assertion. The converse is immediate.  In characteristic zero, the Brauer--Nesbitt character criterion identifies equal trace functions with isomorphic semisimplifications.

The algebra construction is the extension to a coefficient discrete valuation ring of Mayle--Wang, Proposition~12, rather than an assumption of a known congruence depth. Its use of the quotient $M/\varpi M$ is essential: it must not be replaced by entrywise reduction of the ambient pair of matrix algebras.
\end{proof}

The exact analytic functions used by this proof are $\zeta_{\widetilde L}$ and $\zeta_{\mathbb Q}$, where $L$ is the finite quotient just constructed.  Neither the $L$-function of a positive-weight motive underlying $\rho_i$ nor Artin holomorphy is required. The finiteness and characteristic-zero character theory are independent of GRH. An executable numerical test additionally requires a way to compute the indicated exact traces; the theorem by itself does not manufacture that input.

\paragraph{Finding G04 (the precise limit of trace separation).}
Frobenius traces cannot distinguish arbitrary extension classes, even for lisse sheaves on an open subscheme of $\operatorname{Spec}\mathbb Z$.

\begin{proof}
Fix an odd prime $\ell$ and let $c(g)=\log_\ell(\chi_\ell(g)^{\ell-1})$, where $\chi_\ell$ is the cyclotomic character. This is a nonzero continuous additive homomorphism $G_{\mathbb Q,\{\ell\}}\to\mathbb Q_\ell$: its argument lies in $1+\ell\mathbb Z_\ell$, where the $\ell$-adic logarithm is an injective homomorphism, and the cyclotomic image contains this group. The two representations
\[
 \rho(g)=\begin{pmatrix}1&c(g)\\0&1\end{pmatrix},
 \qquad \rho_0(g)=\begin{pmatrix}1&0\\0&1\end{pmatrix}
\]
have trace $2$ and determinant $1$ on every element, including every unramified Frobenius. They are not isomorphic, because one has nontrivial image and the other has trivial image. Their semisimplifications are both the sum of two trivial characters.
\end{proof}

\paragraph{Finding G05 (explicit prime bounds for two LM7 predicates).}
Let $S$ be a finite set of rational primes and put
$m=\operatorname{rad}(2\prod_{p\in S}p)$.  Under the stipulated
Dirichlet GRH there is a prime $q\nmid m$ with $q\equiv2$ or $5\pmod9$
and
\[
 q\le\bigl[72\log3+24\log(4m)+35\bigr]^2. \tag{G05.1}
\]
Under $\mathrm H_M$ there is a prime $q\nmid m$ with $q\equiv4$ or
$7\pmod9$ and $3^{(q-1)/3}\not\equiv1\pmod q$, and
\[
 q\le\bigl[456\log3+72\log(4m)+95\bigr]^2. \tag{G05.2}
\]
These are quantitative refinements of the arithmetic predicates in the
recovered LM7 A.25 and A.16 registers, respectively.  Their source file
\texttt{sections/A\_agent.tex} was not recovered, so this finding does
not verify those rows' separate pro-three group presentations.

\begin{proof}
Put $L_0=\mathbb Q(\zeta_9)$ and $\alpha=\sqrt[3]{3}$.  The extension
$L_0/\mathbb Q$ is abelian of degree six.  The order
$\mathbb Z[\zeta_9]$ has discriminant of absolute value $3^9$:
for $\Phi_9(X)=X^6+X^3+1$,
\[
 \Phi_9'(\zeta_9)=3\zeta_9^2(2\zeta_3+1),\qquad
 N_{\mathbb Q(\zeta_3)/\mathbb Q}(2\zeta_3+1)=3.
\]
Taking the norm from $L_0$, whose degree over $\mathbb Q(\zeta_3)$ is
three, gives $3^9$.  Thus $D_{L_0}\le3^9$.  The two automorphisms
whose residues are $2,5\pmod9$ form the nonempty Frobenius predicate
for (G05.1).

The polynomial $X^3-3$ is Eisenstein at three.  Its real cubic field
$\mathbb Q(\alpha)$ is not Galois, so it cannot be contained in the
abelian extension $L_0$.  No other root can lie in $L_0$, either,
because $\zeta_3\in L_0$.  A cubic polynomial without a root is
irreducible; hence $L=L_0(\alpha)$ has degree eighteen over
$\mathbb Q$.  It is Galois, since $\zeta_3\in L_0$ makes it a
splitting field.  The product basis of
$\mathbb Z[\zeta_9]\mathbb Z[\alpha]$ is a basis of an order in $L$.
The cubic order has discriminant $-3^5$, and the tensor trace matrix
therefore gives
\[
 D_L\le (3^9)^3(3^5)^6=3^{57}. \tag{G05.3}
\]
For each $a=4,7$ and $j=1,2$ there is an automorphism
\[
 \zeta_9\longmapsto\zeta_9^a,\qquad
 \alpha\longmapsto\zeta_3^j\alpha.
\]
These four automorphisms form a nonempty conjugation-stable subset
of $\operatorname{Gal}(L/\mathbb Q)$.  At an unramified prime of residue
$a\pmod9$, the field $\mathbb F_q$ contains $\mu_3$, and $X^3-3$
has a root if and only if Frobenius fixes all its roots.  Thus the
four-element subset describes precisely the noncube predicate in
(G05.2).  Ordinary Chebotarev already gives density $4/18=2/9$;
the first predicate has density $2/6=1/3$.

For either $T=L_0$ or $T=L$, set $n=[T:\mathbb Q]$ and
$\widetilde T=T(\sqrt m)$.  Its degree is at most $2n$, and
\[
 \log D_{\widetilde T}\le2\log D_T+n\log(4m). \tag{G05.4}
\]
Indeed $m$ is positive, even and squarefree, so the quadratic field
has discriminant $4m$.  If it is not contained in $T$, product integral
bases give an order of discriminant $D_T^2(4m)^n$ in the degree-$2n$
compositum, whose field discriminant divides that order discriminant.
If it is contained, $\widetilde T=T$ and the displayed inequality
holds directly.  This covers the possible failure of linear disjointness.

Apply G01 and substitute (G05.4) in G-CHEB.  For $T=L_0$ we have
$[\widetilde T:\mathbb Q]\le12$ and
$\log D_{\widetilde T}\le18\log3+6\log(4m)$, giving (G05.1).
This compositum is abelian over $\mathbb Q$, so the stipulated
Dirichlet GRH supplies its Dedekind-zeta GRH.  For $T=L$ the
corresponding bounds are $36$ and $114\log3+18\log(4m)$, giving
(G05.2) under $\mathrm H_M$.  The chosen primes are unramified in
the relevant fields and avoid every divisor of $m$.  Neither use
requires Artin holomorphy.
\end{proof}

\subsection{Known explicit elliptic-curve consequences}

\paragraph{Cited result G-LOMBARDO (an unconditional bound over number fields).}
Let $E/K$ be an elliptic curve over a number field with
$\operatorname{End}_{\overline K}(E)=\mathbb Z$, let $d_K=[K:\mathbb Q]$,
and let $h_F(E)$ be its stable Faltings height in the convention of the
cited paper.  Lombardo, \emph{Bounds for Serre's open image theorem for
elliptic curves over number fields}, Algebra Number Theory 9 (2015),
2347--2395, Corollary~9.3, p.~2393 \cite{Lombardo2015}, proves
\[
 [\operatorname{GL}_2(\widehat{\mathbb Z}):\rho_{E,\infty}(G_K)]
 < \exp(10^{21483})\,d_K^{\gamma}
       \max\{1,h_F(E),\log d_K\}^{2\gamma},
 \qquad \gamma=2.4\cdot10^{10}. \tag{G-LOMBARDO}
\]
This published result is unconditional. It bounds the adelic image index;
it does not assert surjectivity or identify the exact image. We use the
stated corollary, whose isogeny input is recorded in its Theorem~2.1,
and do not substitute the improvement in Remark~9.4 that invokes the
then-unpublished input of Remarks~2.3 and~2.7. For the actual number
field $L$ in LM8 A.32, it applies directly with $E=F_a$ and $K=L$.

\paragraph{Cited results G-MW1 and G-MW2.}
Mayle--Wang, \emph{On the effective version of Serre's open image theorem}, Bull. Lond. Math. Soc. 56 (2024), 1399--1416, Theorems~1 and~2, give the following conditional bounds. For a non-CM elliptic curve $E/\mathbb Q$, its mod-$\ell$ image equals $\operatorname{GL}_2(\mathbb F_\ell)$ whenever
\[
 \ell>964\log\operatorname{rad}(2N_E)+5760.
 \tag{G-MW1}
\]
For non-CM elliptic curves $E_1,E_2/\mathbb Q$ which are not $\mathbb Q$-isogenous, some prime of good reduction for both has unequal Frobenius traces and satisfies
\[
 p\le\left(482\log\operatorname{rad}(2N_{E_1}N_{E_2})+2880\right)^2.
 \tag{G-MW2}
\]
These become conditional consequences of $\mathrm{H}_M$: the analytic input in their \S3 and the proof of Theorem~2 is G-CHEB for finite auxiliary fields, together with the squarefree quadratic compositum of Corollary~6. The open-image theorem itself and Faltings's isogeny theorem were already unconditional. These numerical bounds are not bounds for arbitrary abelian varieties, arbitrary number fields, or CM elliptic curves. In particular they do not prove Serre's uniformity conjecture.

\paragraph{Finding G06 (the actual LM8 A.32--A.33 elliptic factors).}
The factors used in these rows have the models
\[
 F_a/\mathbb Q:\quad y^2=x^3-3x+a,\qquad a=6,-6.
\]
Under $\mathrm H_M$, their mod-$\ell$ Galois image over $\mathbb Q$ is
$\operatorname{GL}_2(\mathbb F_\ell)$ whenever
\[
 \ell>964\log6+5760. \tag{G06.1}
\]
For the number field $L$ chosen in the recovered proof, the restricted
image satisfies
\[
 [\operatorname{GL}_2(\mathbb F_\ell):
       \overline\rho_{F_a,\ell}(G_L)]\le[L:\mathbb Q]
 \quad\text{for every prime satisfying (G06.1).} \tag{G06.2}
\]
Independently of $\mathrm H_M$, G-LOMBARDO bounds their adelic image
over $L$. Since the models descend to $\mathbb Q$, restriction also gives
\[
 [\operatorname{GL}_2(\widehat{\mathbb Z}):\rho_{F_a,\infty}(G_L)]
 <[L:\mathbb Q]\,\exp(10^{21483})
          \max\{1,h_F(F_a)\}^{4.8\cdot10^{10}}. \tag{G06.3}
\]

\begin{proof}
The models and their role are explicit in recovered LM8 A.32.
Their non-CM hypothesis is supplied by the ordinary-reduction
argument of LM8 A.14: for $F_6$, the traces at $5,11$ are $-1,3$,
giving the distinct imaginary quadratic Frobenius fields with
discriminants $-19,-35$. A geometric characteristic-zero CM field
would inject into both, which is impossible. The other sign is a
quadratic twist. These two point counts were rechecked directly;
no classification of CM $j$-invariants is used.

The discriminant of each displayed integral model is
$-16(4(-3)^3+27\cdot6^2)=-2^9 3^3$.
The model has good reduction outside $2,3$. Its discriminant valuations
$9$ and $3$ are smaller than $12$, so no integral Weierstrass change
can lower them: discriminant valuations of isomorphic models differ
by a multiple of $12$. Both are positive in a minimal model, so both
primes are bad and $\operatorname{rad}(2N_{F_a})=6$.
Substitute this value in G-MW1.

For any homomorphism $\rho$ and subgroup $G_L\le G_{\mathbb Q}$,
the map on coset sets is surjective, giving
$[\rho(G_{\mathbb Q}):\rho(G_L)]\le[L:\mathbb Q]$.
With the mod-$\ell$ surjectivity just obtained this proves (G06.2).
With G-LOMBARDO over $\mathbb Q$, it proves (G06.3).
Full image after restriction to $L$ is not asserted. The original
rows use only qualitative commutator Zariski density, which was
already unconditional; these formulas supply numerical refinements
for their actual elliptic factors.
\end{proof}

\paragraph{Grothendieck framework and locators.}
For $X=\operatorname{Spec}\mathbb Z[1/\prod_{p\in S}p]$, finite quotients of $G_{\mathbb Q,S}$ are connected finite \'etale covers of $X$, and continuous integral adic representations are inverse systems of finite locally constant coefficient sheaves on $X$. The precise Galois-category references are SGA~1, Expos\'e~V, Th\'eor\`eme~4.1 and \S7, together with Propositions~8.1--8.2 (the generic-point description over a normal base). G01--G03 locate a closed point of this arithmetic scheme where a prescribed finite monodromy invariant is visible.  The finite quotient is a weight-zero Artin realization.  It is not a new purity statement for the original positive-weight sheaf.  This section supplies no theorem about crystalline realizations, divided powers, or general motivic $L$-functions; those structures are not used in these proofs.  Cote~149 scans~33--36 concern Kummer point separation and prime-to-characteristic geometric fundamental groups, not the numerical constants of these number-field prime bounds.

\paragraph{Remaining row-level gate.}
To apply a numerical bound to an individual programme row, its complete proof must identify the finite quotient or pair of lisse sheaves, their coefficient ring, rank, ramification support, and the nonempty conjugacy-invariant test set. For an open-image conclusion it must also establish the relevant geometric endomorphism and subgroup hypotheses. None of these inputs can be inferred from a truncated row title alone.


\section{The cyclotomic coordinate of the finite gentle levels}
\label{sec:gt-certificates}

Write $U_N=(\mathbb Z/N\mathbb Z)^\times$, and use natural logarithms.
The word ``level'' in the supplied material can mean a matrix modulus,
a power imposed on a relative kernel, the order of a finite group, or
the modulus of the multiplier itself. These are different data. We
retain the actual quotient maps when making the comparison.

\subsection{The precise Bach input}

\paragraph{Finding T01 (the corrected uniform certificate).}
Under Dirichlet GRH, and hence under Alp\"oge's theorem as stipulated
in this problem, for every integer $N\geq3$ the set
\[
 \mathcal P_3(N)=\{p:p\text{ prime},\ p\nmid N,
                              \ p<3(\log N)^2\}
 \tag{T01}
\]
generates $U_N$. The groups $U_1,U_2$ are trivial and need no generators.
The inequality is strict, and there are no omitted small moduli.
Neither a Hecke GRH over a nonabelian field nor Masullo's claim is needed.

The coefficient $2$ attributed to unit generation in the prompt and
Fable's suggestion~7 is not the statement of Bach's cited theorem.
Bach's Theorem~2, printed page~372, treats a proper subgroup containing
every positive integer below a threshold; integers not coprime to the
modulus are allowed to be outside the subgroup. Its bound is
$2(\log N)^2$. Theorem~3 on the same page imposes coprimality and has
bound $3(\log N)^2$. The consequence for generators is stated on
printed page~373 \cite[Theorems~2--3, pp.~372--373]{Bach1990}.
The distinction is also explicitly recorded in
\cite[\S1.1]{LLS2015}.

\begin{proof}
Only the passage from Bach's integer statement to prime generators
requires proof. Let $H$ be generated by $\mathcal P_3(N)$. Every positive
integer $a<3(\log N)^2$ coprime to $N$ factors into primes in
$\mathcal P_3(N)$, so its residue lies in $H$. If $H$ were proper,
Bach's Theorem~3, applied with $x=3(\log N)^2$, would give
$x<3(\log N)^2$, a contradiction. This proves the assertion. For
$N=1,2$ the assertion is the direct trivial-group calculation.

The distinction between the two theorems cannot be discarded: in
$U_{12}$ the subgroup $\{1,5\}$ omits the positive integer $2$, but
$2$ is not a unit; the least omitted positive unit is $7$. This example
refutes the inference that a witness supplied by Theorem~2 must be a
unit. It does not refute the stronger numerical coefficient-$2$
generation statement, which is not asserted false here.
\end{proof}

The cited theorem's analytic proof has not been independently reproved
in this run, and Bach's original computation through one million moduli
has not been rerun. Its exact hypotheses, constants, strict inequalities,
and small-modulus scope were checked in the original scanned article.
The direct computations below are independent of this analytic input.

\paragraph{Finding T07 (a sharper criterion and a large-modulus consequence).}
Theorem~1.1 of Lamzouri--Li--Soundararajan, in the authors' corrected
arXiv version~v3 of 8~December~2016, applies to every proper subgroup
of $U_N$, for $N\geq3000$, with the same Dirichlet GRH input.
Define
\[
 \begin{split}
 A(N)&=\max\left(0,2\log\log N-\frac85
                    -\sum_{p\mid N}\frac{\log p}{p-1}\right),\\
 B(N)&=\max\left(0,2\log\log N+3
          +\frac{2\omega(N)(\log\log N)^2}{\log N}-2A(N)\right),
 \end{split}
\]
where $\omega(N)$ counts distinct prime divisors. A prime $\ell\nmid N$
outside the subgroup exists with
$\ell\leq(\log N+B(N))^2$
\cite[Theorem~1.1(1), corrected arXiv version~v3]{LLS2015}.
Thus the desired strict coefficient-$2$ generator bound follows
whenever $B(N)<(\sqrt2-1)\log N$, by the same subgroup argument as
T01. Part~(2) gives $\ell\leq(\log N)^2$ under the additional
condition that $N$ has no prime divisor below $(\log N)^2$.
Theorem~1.4 instead concerns a specified coset and depends on its
subgroup index; Corollary~1.1 concerns quadratic nonresidues for prime
moduli. Neither states a uniform coefficient-$2$ generator bound for
arbitrary composite moduli.

The displayed bound in Theorem~1.1(1) is not uniformly below
$2(\log N)^2$ merely for $N\geq3000$. Put
$R(N)=(\log N+B(N))^2/(\log N)^2$. Exact rational interval calculations
give
\[
 \begin{split}
 4.98743&<R(30030)<4.98744,\\
 2.02303&<R(13082761331670030)<2.02304.
 \end{split}
\]
The moduli are the products of the first six and first fourteen primes,
respectively. The executable \texttt{check\_lls\_bound\_scope.py}
certifies these intervals using the positive logarithm series of
(T05.1), also applied to rational endpoints for $\log\log N$;
its two records are in \texttt{LLS\_BOUND\_EXPRESSION\_CHECKS.tsv}.
These are evaluations of the theorem's upper-bound expression,
\emph{not} counterexamples to actual coefficient-$2$ generation.
They show why the displayed LLS theorem plus checking only $N<3000$
does not complete such a uniform claim. The authors' v3 revision notice
identifies corrections to the asymptotic Theorems~1.2--1.3; no constants
from their superseded version are used here.

Under A, the strict coefficient-$2$ generators also suffice for every
integer $N\ge\exp(300)$. This deliberately coarse threshold is a new
deduction of the displayed LLS bound, not a threshold quoted from that
paper; the remaining interval is not asserted checked.

\begin{proof}
The sufficient $B(N)$ criterion follows by the proper-subgroup argument
in T01. The two expression enclosures follow from the exact rational
positive-series estimates and complete calculations just specified.
For the final assertion put $L=\log N\ge300$ and $l=\log L$.
Since $2^{\omega(N)}\le N$ and $\log2>2/3$, the defining formula gives
$B(N)<3l^2+2l+3$. Indeed $A(N)\ge0$ and
$2\omega(N)l^2/L<3l^2$. The function
$F(L)=(3(\log L)^2+2\log L+3)/L$ has derivative
$-(3\log L-1)(\log L-1)/L^2<0$ for $L\ge300$.
Also $\log300<6$, since $e>8/3$ and $(8/3)^6>300$.
Thus $F(L)\le F(300)<123/300=41/100$ and
\[
 (L+B(N))^2<(141/100)^2L^2
             =\frac{19881}{10000}L^2<2L^2.
\]
LLS Theorem~1.1 applies since $N\ge\exp(300)>3000$ and gives a
coprime prime outside any proper subgroup below this last bound.
The same subgroup argument proves generation. The elementary lower
bound $\log2>2/3$ is, for example, the first term plus the positive
remainder in $\log2=2\sum_{j\ge0}(1/3)^{2j+1}/(2j+1)$.
\end{proof}

\subsection{What the finite certificate certifies}

\paragraph{Finding T02 (finite cyclotomic tests and their limits).}
Let $\mathcal G=\widehat{\mathrm{GT}}_0$, with multiplier
$\lambda:\mathcal G\to\widehat{\mathbb Z}^{\times}$, and let
$q:\mathcal G\twoheadrightarrow Q$ be a finite quotient. To speak of
its full multiplier modulo $N$, require an actual homomorphism
$\lambda_N:Q\to U_N$ with
\[
                    \lambda\bmod N=\lambda_N\circ q.
 \tag{T02.1}
\]
Equivalently, $\ker q\subseteq\ker(\lambda\bmod N)$. If this is not
available, the actual finite image of $(q,\lambda\bmod N)$ is a
specified refinement that does retain the multiplier; this must not
silently be identified with $Q$.

Suppose $P$ is any proved generating set of $U_N$, for instance
$\mathcal P_3(N)$ under T01, or a list certified in T05. Then:
\begin{enumerate}
\item a subgroup $J\leq Q$ has $\lambda_N(J)=U_N$ if its multiplier
image contains every residue in $P$;
\item two homomorphisms from $U_N$ agree if they agree on $P$;
\item a nontrivial character of $U_N$ is nontrivial on some member of $P$.
\end{enumerate}
These statements concern subgroups and homomorphisms already shown to
have the indicated algebraic properties. A set of candidate shadows
has not thereby been shown to be a subgroup. The certificate does not
check the two cycle equations, the pentagon, the $f$-coordinate, or
compatible lifting through all finite quotients.

\begin{proof}
Each assertion follows by writing every unit as a word in $P$. For
the third, otherwise the kernel of the character would contain all of
$U_N$. For the first, it is essential that the image be a subgroup,
so it contains all words in its specified generators. The factorization
criterion in (T02.1) is the elementary universal property of a quotient.

For the limitation, retain the supplied LM7\_PUSH E.65 theorem:
there is a continuous character $\Delta:\mathcal G\to\mathbb F_2$
vanishing on the Galois subgroup and taking value $1$ on an actual
element of $\ker\lambda$ \cite[E.58, E.65]{LM7run}. It cannot factor
through $\lambda$ or through any $\lambda\bmod N$, since a homomorphism
of that form is zero on $\ker\lambda$. Moreover $\ker\Delta$ is proper
but has the same full cyclotomic image as $\mathcal G$: it contains
$G_{\mathbb Q}$, whose cyclotomic character is onto. Thus even all of
the cyclotomic finite certificates together cannot distinguish these
two subgroups. This is a concrete counterexample within the supplied
gentle-group framework to inferring equality of groups from equality
of their multiplier images. It uses the supplied E.65 theorem as
background, not a new proof of E.58's profinite construction.
\end{proof}

For clarity, the arithmetic inclusion already gives
$\lambda(\mathcal G)=\widehat{\mathbb Z}^{\times}$, independently of
GRH. When an actual finite quotient retains $\lambda\bmod N$, its
full multiplier image is therefore already $U_N$. Bach's input gives
a bounded list of witnesses or tests, not the original surjectivity
theorem. Also, an arithmetic Frobenius at $p\nmid N$ has multiplier
$p$ modulo $N$; no claim is made that its full profinite multiplier
is the nonunit integer $p$ in $\widehat{\mathbb Z}$.

\paragraph{The scheme and fibre-functor formulation.}
Put $S_N=\operatorname{Spec}\mathbb Z[1/N]$ and
$X_N=\operatorname{Spec}\mathbb Z[1/N,\zeta_N]$. This is the connected
finite \'etale Galois cover with group $U_N$. Indeed it is the primitive
root component of the \'etale scheme of roots of $T^N-1$ after $N$ is
inverted; the cyclotomic field identifies its generic Galois group.
Arithmetic Frobenius sends $\zeta_N$ to $\zeta_N^p$. T02 is consequently
a certificate on a finite quotient of the fibre-functor automorphism
group of $S_N$, in SGA~1, Expos\'e~V, Th\'eor\`eme~4.1 and \S7
\cite{SGA1}. This is the applicable Grothendieck framework. It supplies
no additional purity, crystalline, divided-power, or motivic assertion.
For each prime $\ell$ the compatible roots of unity are the Tate module
$\mathbb Z_\ell(1)$; finite level records its cyclotomic scalar action.

\subsection{Elementary certificates for the two- and three-primary towers}

\paragraph{Finding T03 (strict coefficient-$2$ certificates without GRH).}
For integers $a,b\geq0$ and $N=2^a3^b>1$, the following prime lists generate $U_N$, and every
listed prime is coprime to $N$ and strictly less than $2(\log N)^2$:
\[
\begin{array}{c|c}
 (a,b)&\text{prime generators}\\\hline
 (1,0)&\varnothing\\
 (2,0)&3\\
 a\geq3,\ b=0&3,5\\
 a=0,\ b\geq1&2\\
 a=1,\ b\geq1&5\\
 a=2,\ b\geq1&5,7\\
 a\geq3,\ b\geq1&5,7,13
\end{array}
 \tag{T03}
\]
Also $2$ generates $U_{5^b}$ for every $b\geq1$ and satisfies the same
strict bound. These are elementary unconditional statements, requiring
neither of the analytic inputs in the prompt.

\begin{proof}
For $a\geq3$, the subgroup of units congruent to $1$ modulo $4$ has
order $2^{a-2}$ and is generated by $5$. The elementary valuation
identity
\[
                  v_2(5^{2^j}-1)=j+2
 \tag{T03.1}
\]
proves the order: the base case is $5-1=4$, and on each squaring the
factor $5^{2^j}+1$ has valuation one. Since $3\equiv-1\pmod4$, the
two primes $3,5$ generate all of $U_{2^a}$. The case $a=2$ is direct.

Both $2$ and $5$ have order $2\cdot3^{b-1}$ modulo $3^b$. To see this,
they are $-1$ modulo three, and
\[
 v_3(2^2-1)=v_3(5^2-1)=1,
 \qquad v_3(z^3-1)=v_3(z-1)+1\quad(z\equiv1\pmod3).
 \tag{T03.2}
\]
The latter identity follows by expanding $(1+3^s u)^3-1$ for
$s\geq1$ and $3\nmid u$: its first term has uniquely least
valuation $s+1$. Repeated use gives the asserted orders, equal to
$|U_{3^b}|$. This proves the first two ternary cases.

For $a=2,b\geq1$, the quotient of $U_N$ by its squares is canonically
$U_{12}$, generated by $5,7$: their products give $1,5,7,11$.
The projection of $5$ generates the ternary primary component, by (T03.2).
Thus the subgroup they generate is all of $U_N$. Here is a proof
that these two checks suffice: if $H$ is proper in a finite abelian
group $G$, then $G/H$ has a quotient of prime order; the only
possible primes dividing $|U_N|$ are two and three. Such a quotient
would contradict one of the two checks.

For $a\geq3,b\geq1$, the quotient by squares is instead $U_{24}$.
Its eight distinct elements are the products of $5,7,13$:
\[
                    1,5,7,13,11,17,19,23.
 \tag{T03.3}
\]
Reduction from $U_{2^a}$ to $U_8$ kills exactly its squares, using
(T03.1); reduction from $U_{3^b}$ to $U_3$ does the same for its
two-primary quotient, using (T03.2). This proves the assertion about
the quotient by squares. The projection of $5$ again generates the
three-primary component. The preceding prime-order-quotient argument
finishes the mixed case. Empty or absent primary components cause
no exception.

For $5^b$, $2$ has order four modulo five, and $v_5(2^4-1)=1$.
The same binomial argument, now with fifth powers, shows
$v_5(2^{4\cdot5^j}-1)=j+1$, so its order is $4\cdot5^{b-1}$.

It remains to verify the strict bounds for these entire families,
rather than just for a finite sample. The elementary series
$\log2=2(1/3+(1/3)^3/3+\cdots)$ gives $\log2>2/3$; also
$e<3$, by its factorial series, gives $\log3>1$. For the minimal
moduli $3,4,8,6,12,24$ in the nonempty rows, respectively,
these inequalities give
\[
 2(\log N)^2>2,\quad 32/9,\quad8,\quad50/9,\quad98/9,\quad18.
\]
These exceed respectively the largest proposed primes $2,3,5,5,7,13$.
The bound increases with $N$. For $5^b$ it suffices that $5>3$.
This proves strictness in every case, including all boundary cases.
\end{proof}

\subsection{Identifying the actual scalar moduli}

For the LM7 homology tower, retain the actual maps
\[
 B=\langle s,r\mid s^2=r^3=1\rangle,
 \quad t=sr,\quad
 q:B\longrightarrow\operatorname{SL}_2(\mathbb Z/12)/\{\pm I\},
\]
where $q(s)$ and $q(r)$ are the classes of
$S=\left(\begin{smallmatrix}0&1\\-1&0\end{smallmatrix}\right)$ and
$R=\left(\begin{smallmatrix}0&-1\\1&-1\end{smallmatrix}\right)$.
Set $N_B=\ker q$, $C=N_B^{\rm ab}$, and
$G_j=B/(N_B'N_B^j)$. These are the discrete source definitions of
LM7 E.1/E.4; the maps to their finite targets extend to completions.

\paragraph{Finding T04 (source-to-modulus comparison).}
In the recovered LM7\_PUSH E.1/E.4 homology tower, the cusp has order
$12j$ in $G_j=B/(N_B'N_B^j)$. The specifically reported coefficient
checks $j=2,4,8,16,3,9,27,81$, together with the retained $j=6$, therefore
have cusp scalar moduli
\[
        24,48,96,192,36,108,324,972,72.
 \tag{T04.1}
\]
For the correcting characters of LM8\_PUSH C.33--C.38,
\[
 c_{p,k}(u)=\frac{\log(u^{p-1})}{p}\pmod{p^k},
 \qquad p\geq5,\quad k\geq1,
 \tag{T04.2}
\]
the exact scalar conductor is $p^{k+1}$. At $k=1$ this is
$(u^{p-1}-1)/p\bmod p$. The T01 cutoff at that modulus is
$3(k+1)^2(\log p)^2$.

\begin{proof}
E.1 identifies $C=N_B^{\rm ab}\simeq\mathbb Z^{97}$ and records that
$c=[t^{12}]$ is primitive \cite[E.1, proof; E.4]{LM7run}.
The matrix $SR=\left(\begin{smallmatrix}1&-1\\0&1\end{smallmatrix}\right)$
has order twelve even modulo the single sign, since its diagonal
entries cannot become $-1$ modulo twelve. Thus a power $t^r$ can be trivial
in $G_j$ only if $r=12s$. Its value in the kernel $C/jC$ is then
$sc$, zero exactly when $j\mid s$. Hence $\operatorname{ord}_{G_j}(t)=12j$.
This identifies the cusp scalar, not the whole order of $G_j$.
The pure peripheral element $t^2$ has order $6j$, so the half-multiplier
$(\lambda-1)/2$ in the cycle equations is also determined by
$\lambda\bmod12j$.

For (T04.2), if $u/v\in1+p^{k+1}\mathbb Z_p$, then
$\log((u/v)^{p-1})\in p^{k+1}\mathbb Z_p$; division by $p$ makes
the difference zero modulo $p^k$. Thus the character factors through
$U_{p^{k+1}}$. On $u=1+p^k$, expansion of the logarithm gives
\[
 c_{p,k}(1+p^k)\equiv(p-1)p^{k-1}\not\equiv0\pmod{p^k}.
\]
The omitted terms have valuation at least $k$ after division by $p$,
including $k=1$ since $p\geq5$. Therefore it does not factor through
$U_{p^k}$, proving exactness. Reduction of
$\log(1+pa)/p$ modulo $p$ is $a$, which proves the $k=1$ formula.
\end{proof}

More generally, in a finite characteristic free-group target $H$ of
exponent $e$, reduction of $\lambda$ modulo $2e$ determines every
exponentiation by $\lambda$ and by $(\lambda-1)/2$ in the two cycle equations.
When $e$ is odd, reduction modulo $e$ suffices, since two is invertible.
This gives a safe finite scalar parameter once the actual target is
specified; it is not a claim that $2e$ is minimal. For example LM8
C.14 has exponent twelve and explicitly uses multiplier modulus
twenty-four, whereas LM7 D.7 has exponent $101$ in its pure target
and uses $m=0\bmod101$, equivalently $\lambda=1\bmod101$.
The redundant odd lift modulo $202$ gives the same unit group.

The complete recovered-source inventory is
\texttt{GT\_LEVEL\_SCOPE.tsv}. The following are its principal entries;
``frame'' and ``coefficient'' entries are deliberately not promoted
to unspecified scalar quotients.
\begin{center}
\begin{tabular}{p{0.24\linewidth}p{0.23\linewidth}p{0.45\linewidth}}
\hline
Source & Integer or family & Actual meaning and scope\\\hline
LM7 E.1,E.4 & $12j$ & Cusp order in the specified homology quotient.\\
LM7 E.5 & $360$ & Cusp order in $G_6\times A_5$; constructed multiplier one.\\
LM7 E.13,E.14 & $288,576$ & Explicit multiplier-one congruences in the full register.\\
LM7 E.22,E.24 & $72,144$ & Cusp order and the branches $1,73\bmod144$ over retained points.\\
LM7 D.7 & $101$ & Odd-prime multiplier parameter; the tested stratum has multiplier one.\\
LM7 E.50,E.51 & $20,28$ & Actual projective congruence quotients; symplectic tests have multiplier one.\\
LM7 E.52--E.56 & $m,2m,4m$, $(m,6)=1$ & Infinite congruence and exponent families, with specified maps.\\
LM7 E.65 & $12$ & Signed CM frame modulus on $\ker\lambda$, not a cyclotomic character.\\
LM8 C.12 & Every $N$ & Normalized implementer $\left(\begin{smallmatrix}\lambda&t\\0&1\end{smallmatrix}\right)$; determinant is $\lambda$.\\
LM8 C.14 & $24$ & Explicit multiplier modulus for the exponent-twelve $S_4$ envelope.\\
LM8 C.18 & $15$ & Exact domain character $(\lambda/5)(\lambda\bmod3)$.\\
LM8 C.33--C.38 & $p^{k+1}$ & Scalar correcting character (T04.2), not the whole corrected character.\\
LM8 C.49 & $24,60$ & Framed images restricted to multiplier one; no full arbitrary-multiplier image claimed.\\
GT\_TOWER summary & $4,12;\ m=2,3,6;\ 6m=12,18,36$ & Named matrix and power levels only; the full register is unavailable.\\\hline
\end{tabular}
\end{center}

For C.12 the determinant of an arbitrary projective matrix is defined
only modulo squares. It is the displayed unique normalized
representative that retains the full multiplier. For C.49, the orders
$96,48,2880,1440$ of the two framed/arithmetic images are not
cyclotomic moduli. Likewise $T_{24}$ and $T_{48}$ in LM7 E.59/E.62
name targets built from matrix groups of those orders; their subscripts
are not established multiplier moduli. The finite coefficient depths
in the relative nilpotent towers have the same distinction.

\subsection{Direct finite certificates and verification}

\paragraph{Finding T05 (53 exact finite checks).}
For all 53 integers in \texttt{GT\_UNIT\_CERTIFICATES.tsv}, the primes
strictly below $2(\log N)^2$ and coprime to $N$ generate $U_N$.
This assertion is unconditional and finite. A literal word in the
listed prime generators is supplied for every unit in
\texttt{GT\_UNIT\_WORDS.tsv}; there are 4,825 such words in total.
The table also verifies the coefficient-$3$ bound, though no analytic
theorem is needed for these individual computations.

\begin{proof}
The executable \texttt{verify\_gt\_units.py} constructs a subgroup by
successive multiplication and records a nonnegative prime exponent
vector for every reached residue. All arithmetic of group elements is
integer modular arithmetic. The separate read-only
\texttt{check\_unit\_certificates.py} reads the saved tables, without
performing any subgroup search, and re-evaluates each word. It verifies
that each result is a unit, no residue is repeated, and the number
of distinct residues is
$\varphi(N)=N\prod_{p\mid N}(1-1/p)$. It also compares them with
every integer residue coprime to $N$. Thus the displayed products
constitute a certificate of generation, independently of the method
used to discover them. Both executables were run successfully.

Strict bounds are certified with rational intervals, rather than
rounded logarithms. Write $N=2^h r$ with $1\leq r<2$. For rational
$a/b\in[1,2]$, set $t=(a-b)/(a+b)$ and use
\[
 2\sum_{j=0}^{K-1}\frac{t^{2j+1}}{2j+1}
 \ \leq\ \log(a/b)\ \leq\
 2\sum_{j=0}^{K-1}\frac{t^{2j+1}}{2j+1}
       +\frac{2t^{2K+1}}{(2K+1)(1-t^2)}.
 \tag{T05.1}
\]
The error bound follows by replacing every denominator of the
remaining positive terms by $2K+1$, which can only enlarge the sum,
then summing the resulting geometric series. Combining the
intervals for $h\log2$ and $\log r$, and squaring their nonnegative
endpoints, gives rational bounds for $c(\log N)^2$. The program
increases $K$ until both endpoints have the same integer part and
the lower endpoint exceeds that integer. Every proposed prime below
or equal to that integer is therefore strictly below the true bound.
This test terminated for all reported moduli at both $c=2$ and $c=3$.
The printed decimals are only a convenient display of these checks.
\end{proof}

\begin{longtable}{rrrrl}
\caption{Exact finite unit-generation certificates. Decimal bounds are informational; the cutoff column is certified with rational intervals. The lists are generating lists, not claims of minimal cardinality.}\\
\hline
$N$&$\varphi(N)$&$2(\log N)^2$ (display)&Exact cutoff&Prime generators\\\hline
\endfirsthead
\hline
$N$&$\varphi(N)$&$2(\log N)^2$ (display)&Exact cutoff&Prime generators\\\hline
\endhead
2 & 1 & 0.960906 & 0 & $\varnothing$ \\
3 & 2 & 2.413898 & 2 & $2$ \\
4 & 2 & 3.843624 & 3 & $3$ \\
5 & 4 & 5.180581 & 5 & $2$ \\
6 & 2 & 6.420804 & 6 & $5$ \\
7 & 6 & 7.573133 & 7 & $2,3$ \\
8 & 4 & 8.648154 & 8 & $3,5$ \\
9 & 6 & 9.655592 & 9 & $2$ \\
10 & 4 & 10.603796 & 10 & $3$ \\
11 & 10 & 11.499803 & 11 & $2$ \\
12 & 4 & 12.349522 & 12 & $5,7$ \\
13 & 12 & 13.157930 & 13 & $2$ \\
14 & 6 & 13.929247 & 13 & $3$ \\
15 & 8 & 14.667072 & 14 & $2,7$ \\
16 & 8 & 15.374496 & 15 & $3,5$ \\
17 & 16 & 16.054196 & 16 & $2,3$ \\
18 & 6 & 16.708498 & 16 & $5$ \\
19 & 18 & 17.339442 & 17 & $2$ \\
20 & 8 & 17.948824 & 17 & $3,11$ \\
23 & 22 & 19.662648 & 19 & $2,5$ \\
24 & 8 & 20.200052 & 20 & $5,7,13$ \\
25 & 20 & 20.722323 & 20 & $2$ \\
27 & 18 & 21.725081 & 21 & $2$ \\
28 & 12 & 22.207174 & 22 & $3,5$ \\
29 & 28 & 22.677362 & 22 & $2$ \\
30 & 8 & 23.136287 & 23 & $7,11$ \\
31 & 30 & 23.584536 & 23 & $2,3$ \\
36 & 12 & 25.683216 & 25 & $5,7$ \\
48 & 16 & 29.972395 & 29 & $5,7,13$ \\
49 & 42 & 30.292530 & 30 & $2,3$ \\
60 & 16 & 33.527315 & 33 & $7,11,13$ \\
64 & 32 & 34.592617 & 34 & $3,5$ \\
72 & 24 & 36.579746 & 36 & $5,7,13$ \\
81 & 54 & 38.622367 & 38 & $2$ \\
96 & 32 & 41.666549 & 41 & $5,7,13$ \\
101 & 100 & 42.598675 & 42 & $2$ \\
108 & 36 & 43.844706 & 43 & $5,7$ \\
120 & 32 & 45.840154 & 45 & $7,11,13,17$ \\
121 & 110 & 45.999214 & 45 & $2$ \\
144 & 48 & 49.398088 & 49 & $5,7,13$ \\
169 & 156 & 52.631722 & 52 & $2$ \\
192 & 64 & 55.282515 & 55 & $5,7,13$ \\
202 & 100 & 56.355412 & 56 & $3$ \\
288 & 96 & 64.138243 & 64 & $5,7,13$ \\
289 & 272 & 64.216783 & 64 & $2,3$ \\
324 & 108 & 66.833991 & 66 & $5,7$ \\
360 & 96 & 69.292441 & 69 & $7,11,13,17$ \\
361 & 342 & 69.357767 & 69 & $2$ \\
529 & 506 & 78.650592 & 78 & $2,5$ \\
576 & 192 & 80.800209 & 80 & $5,7,13$ \\
841 & 812 & 90.709450 & 90 & $2$ \\
961 & 930 & 94.338145 & 94 & $2,3$ \\
972 & 324 & 94.651073 & 94 & $5,7$ \\
\hline
\end{longtable}


\paragraph{Finding T06 (the remaining evidence boundary).}
The six attached \texttt{SORTED\_*.md} files are summaries, not the
complete row registers their filenames are said to accompany. The
GT\_TOWER\_CHECK summary names a register with 180 rows, but does not
contain that register or the exact definitions of every finite target.
Search and recovery in this run obtained LM8's full 304-row register
and source, and LM7's full 218-row register with only part of its proof
sources. They did not obtain GT\_TOWER\_CHECK's full register or paper.
Consequently T05 proves every listed arithmetic certificate, but is
not described as an exhaustive computation of all finite levels in
the unrecovered 180-row run. The precise remaining gate is to obtain
its finite target definitions and identify their multiplier factors;
the reusable verifier then checks any additional explicit moduli.
For the genuinely infinite families already present in the recovered
sources, T01 applies at every identified modulus, and T03 gives the
stronger unconditional coefficient-$2$ result for the indicated towers.

\begin{proof}
This is a source-coverage finding. The attached summary's own register
count is 180; its displayed statements are truncated and list only
selected rows. The recovered LM7 and LM8 CSV files have 218 and 304
rows respectively. The source inventory records, for each use, whether
the underlying proof, full statement, or only summary was available.
Thus neither a missing row nor an unknown quotient map has been
silently assigned an integer modulus. No analytic assumption can
supply those omitted definitions.
\end{proof}


\section{Audit of the available rows and their analytic hypotheses}
\label{sec:row-audit}

\paragraph{Finding R01 (exact source coverage).}
The six attached \texttt{SORTED\_*.md} files are summaries, not the
complete row lists named in them. Their statements are frequently
truncated. They report $1190$ register rows. Subsequent recovery gives
the following exact inventory; ``proof blocks'' means available numbered
source blocks, not independently verified proofs.
\begin{center}
\begin{tabular}{lrrr}
\hline
Run & Reported rows & Recovered rows & Available proof blocks\\
\hline
\texttt{ANABELIAN\_A1} &121&121&121\\
\texttt{ANABELIAN\_A1\_SUCCESSOR}&210&0&0\\
\texttt{ANABELIAN\_SUCCESSOR}&157&157&157\\
\texttt{GT\_TOWER\_CHECK}&180&0&0\\
\texttt{LM7\_PUSH}&218&218&9\\
\texttt{LM8\_PUSH}&304&304&304\\
\hline
Total&1190&800&591\\
\hline
\end{tabular}
\end{center}
The nine LM7 blocks are E.1--E.7, E.58 and E.65. The $582$ blocks in
the complete A1, successor and LM8 main sources were available for
text screening. The targeted analytic passages were read; availability
does not mean that every mathematical proof was re-proved.

The flag file has $575$ records, representing $525$ distinct run/statement
pairs in $23$ runs. Its six-run subset has $166$ records and $152$
distinct pairs. The full audit table preserves all $800$ recovered rows,
and the curated table has $108$ entries, including the two other flagged
run examples, three intake items, the manuscript Artin passage and
Finding A06. Declared register dependencies give $201$ rows with a path
to one of $38$ curated qualitative prime-selection inputs. These paths
are recorded dependencies, not a newly certified proof graph.

For the two missing full registers, exactly $47$ A1-successor and $50$
GT identifiers are visible in the supplied summaries or flags. They
are listed separately; no unseen rows are invented. The remaining
coverage gap is therefore $390$ full register rows, together with $209$
LM7 proof blocks whose statements were recovered.

\paragraph{Finding R02 (the RH/GRH flags are false positives for conditional dependence).}
There are just two flags labelled RH/GRH. In LM8 A.69, the abbreviation
``DS/RH'' is resolved by the full proof as Deuring--Shafarevich and
Riemann--Hurwitz, not the Riemann Hypothesis. In
\texttt{S6\_SERIES\_GRIND}, row \texttt{maps:src:F.2.36}, the attached
context explicitly says that the class and unit groups were certified
with \texttt{GRH=false}. Its underlying computation was not recovered,
so this is an audit of the reported scope, not a new certification.
Neither flag establishes a formerly conditional theorem that is now
unconditional. Conversely, their absence from other rows does not show
analytic independence: the recovered LM8 proofs contain many
qualitative Chebotarev arguments.

\subsection{Qualitative prime choice and its quantitative refinements}

The finite-Galois qualitative Chebotarev theorem is unconditional:
a conjugacy class $C$ has density $|C|/|G|$. The exact source statement
checked here is Milne, \emph{Class Field Theory}, v4.03,
Chapter V, Theorem 3.23, printed p.~164 \cite{MilneCFT}.
The existing existence and density uses below therefore require neither
stipulated GRH hypothesis. A new least-prime version follows from
Finding G01 after the actual finite field and excluded primes have
been specified. Its exact analytic hypotheses are GRH for
\[
 \zeta_{\widetilde L}(s),\quad\zeta_{\mathbb Q}(s),\qquad
 \widetilde L=L\!\left(\sqrt{\operatorname{rad}
                     (2\prod_{p\in S}p)}\right).
\]
For a number-field base one uses the normal closure over $\mathbb Q$
and the rational primes below the excluded places. The stipulated
all-number-field Hecke hypothesis supplies these zeta functions as
trivial Hecke characters. Dirichlet GRH supplies this particular
argument when $L/\mathbb Q$ is abelian. No general Artin holomorphy
assumption is introduced.

\begin{longtable}{p{.22\linewidth}p{.43\linewidth}p{.26\linewidth}}
\hline
Rows & Actual finite field or predicate & Quantitative family\\
\hline
\endhead
LM8 A.16 & $M/\mathbb Q$ finite Galois containing the normalizer field;
complex-conjugation class. & $\zeta_M$, with G01 exclusions.\\[3pt]
LM8 B.20, B.22 & Normal closure of the specialized finite torsor fibre,
after the quotient, deck-constant field $K$, and Hilbert point $a$
are chosen. & Zeta of this normal closure.\\[3pt]
LM8 C.30; D.1; D.58 & Respectively
$\mathbb Q(\zeta_{12})$ with class $5$,
$\mathbb Q(\zeta_{5^r})$ with class $1$,
$\mathbb Q(\zeta_{30})$ with class $17$. & Dirichlet $L$-functions
modulo $12$, $5^r$, $30$, with finite exclusions.\\[3pt]
LM8 D.8 & $\mathbb Q(E_c[5],\mu_5)$; completely split primes.
& Zeta of this division field.\\[3pt]
LM8 D.11 & Kernel field of
$\rho_{M(-1)}\bmod\ell^{k+1}$ over the stated number field.
& Zeta of its $\mathbb Q$-normal closure.\\[3pt]
LM8 D.3, D.28 & Kernel field of
$g\mapsto(s_0(g)\bmod5,\chi_5(g)\bmod5^{r+1})$;
target $(\exp C,-1+5^r)$ in the actual Heisenberg image.
& Zeta of this actual image field.\\[3pt]
LM8 D.44, D.48, D.62, D.63, D.65 & Joint kernel field of the actual
finite Alexander action on $W_r$ and $\chi_5\bmod5^r$;
D.63 uses $r+1$. & Zeta of the joint image field.\\[3pt]
LM8 D.60, D.61, D.66 & Kernel field of the finite affine homomorphism
$g\mapsto(z(g),\rho_r(g))$ with additive coefficient module $W_r$.
& Zeta of the actual affine image field.\\[3pt]
LM8 D.67, D.69--D.71, D.73 & The same construction with finite
nonabelian coefficient group $H_r$ and its actual action;
arbitrary number field in D.70. & Zeta of the $\mathbb Q$-normal closure.\\[3pt]
LM8 D.78--D.87 & At each original finite Cartier-dual coefficient
module, the actual finite affine cocycle/action field.
& Zeta of that finite field.\\[3pt]
LM8 E.1, E.2 & $K=\mathbb Q(i)$ and $K(P)/K$ for a generator
$P$ of $E_6[12]$; degree $16$ in the source calculation.
& Finite-order Hecke $L_K(s,\chi)$, or $\zeta_{K(P)}$.\\[3pt]
LM7 A.16 & $\mathbb Q(\zeta_9,\sqrt[3]3)$;
$q\equiv4,7\pmod9$ and $3$ noncubic modulo $q$.
& Zeta of this degree-$18$ field.\\[3pt]
LM7 A.25 & $\mathbb Q(\zeta_9)$; classes $2,5$.
& Dirichlet $L(s,\chi)$ modulo $9$.\\
\bottomrule[0.4pt]
\end{longtable}

The LM7 proof sources for A.16 and A.25 are missing. Their fields are
reconstructed from the full register predicates, not from an asserted
recheck of their group presentations. Over $\mathbb Q(\zeta_3)$, the
extensions adjoining $\zeta_9$ and $\sqrt[3]3$ are independent cubic
extensions; four of the $18$ elements of the full Galois group satisfy
the conjugacy-stable A.16 predicate. This gives density $2/9$.
A.25's two classes among the six units modulo $9$ give density $1/3$.

For an additive or nonabelian cocycle, the finite homomorphism is the
actual graph $g\mapsto(z(g),\rho(g))$. A witness $\delta$ in the
action kernel with $z(\delta)\ne0$, respectively $\ne1$, gives a
nonempty conjugacy class in that actual image. One cannot replace it
by an arbitrarily chosen element of a larger formal affine group.
The recovered LM8 proofs make this distinction explicitly.

LM8 A.20 provides a separate logical warning: for a nontrivial quadratic
extension $L/K$, the inert primes have positive density and are
disjoint from the split primes, although the split set becomes
density one after pullback to $L$. This unconditional counterexample
does not refute the stronger arithmetic theorem for which a stronger
density hypothesis might suffice.

Stagewise effective prime choice does not bound a preceding Hilbert
specialization, a separating finite quotient, or the first obstruction
depth of an inverse system. LM7 A.43 and B.6 already give finite
field enumerations at each fixed depth from degree and discriminant
bounds; they expressly assert no uniform stopping rule for every
positive infinite-tower instance.

\subsection{Primitive roots, geometric existence and finite-field weights}

\paragraph{Finding R03 (fixed-prime criteria are not Artin-existence theorems).}
A1 B.3 and successor A.8, A.10, A.17 and A.35 use a hypothesis such as
$\operatorname{ord}_\ell(p)=\ell-1$ at a specified prime. This finite
condition makes $\Phi_\ell$ irreducible over $\mathbb F_p$ and its
nontrivial character values one Frobenius orbit. It does not produce
a prime satisfying the additional geometric or numerical conditions.
Successor A.12 supplies a genus-three counterexample to extending a
genus-two point-only test, even with $2$ primitive modulo $3$.

For the separate assertion of infinitely many primitive-root primes,
Hooley requires the family
\[
 K_n=\mathbb Q(\zeta_n,p^{1/n}),\qquad n\text{ squarefree},
 \qquad \zeta_{K_n}(s).
\]
The theorem was checked in \S7, printed p.~47, of the Russian
translation of his 1967 paper. The fields are generally nonabelian
over $\mathbb Q$, so the stipulated Dirichlet hypothesis does not
supply this family. The stipulated all-field Hecke hypothesis does.
Its arithmetic conclusion gives no theorem about intersection with a
curve-dependent theta-support set. The universal primitive-root theta
condition is in fact refuted by Finding A06. That example has
$p$-rank zero, and leaves the positive-$p$-rank restricted problem
unresolved. It is not a counterexample to A1.

The elementary density obstruction also remains useful: when
$\ell\equiv1\pmod8$, both $2$ and $-1$ are squares, so
$\langle2,-1\rangle$ is a proper subgroup of
$\mathbb F_\ell^\times$. A positive-density support set can therefore
be disjoint from the desired arithmetic condition. For torsion order
$\ell^a$, primitivity modulo $\ell$ also does not automatically
give primitivity modulo $\ell^a$.

A1 B.3 and successor A.11 instead use finite-field $L$-polynomials:
\[
 L(C,\chi,T)=\det(1-\operatorname{Frob}_qT
                  \mid H^1(C_{\bar k},\mathcal L_\chi)).
\]
Here $C/k$ is proper and pointed by $P\in C(k)$, $k=\mathbb F_q$,
$a:J_C(k)\twoheadrightarrow A$ is a finite prime-to-$p$ quotient,
$N=\#C(k)$, and
$N_b=\#\{R\in C(k):a([R-P])=b\}$.
Write $S_\chi=\sum_{R\in C(k)}\chi(a([R-P]))$.
The nontrivial Lang eigensheaves remain geometrically nontrivial
on $C$, since its Abel map surjects on the geometric abelian
prime-to-$p$ fundamental group. These eigensheaves have $2g-2$
eigenvalues of absolute value $\sqrt q$. The resulting pair estimate is
\[
 \sum_bN_b(N_b-1)
 \le {N^2\over|A|}
   +{|A|-1\over|A|}(2g-2)^2q-N.
\]
Deligne's \emph{Weil I}, formula (1.5.1) and Theorem (1.6), are
unconditional finite-field inputs. Neither number-field GRH hypothesis
supplies the still missing quotient or prescribed-point existence in
A1 B.5 and successor A.15. Successor A.38 explicitly distinguishes
closed-point Chebotarev from prime-order theta avoidance on an entire
nonzero cyclic orbit.

The least-prime results A1 B.27, B.30 and B.34, and successor
A.19, A.24, A.31, A.72, A.74, A.91--A.93 and A.95, concern specified
curves or explicitly bounded finite censuses. They use finite
Frobenius/Cartier calculations and finite exclusions, not analytic
least-prime theorems. A1 B.34 is a \emph{field-first} minimum, not a
least prime over arbitrary extensions. A1 B.32 and successor A.96
retain unresolved cases. A1 B.4 decides a fixed prime over all
constant extensions; it does not bound the search over all primes.
Successor A.16 rules out a bound depending only on $(p,g)$.

\subsection{Other false positives and verification scope}

LM8 A.32--A.33 use Serre's qualitative non-CM open-image theorem;
their commutator Zariski-density conclusion was already unconditional.
The recovered proofs explicitly identify
$F_a/\mathbb Q:y^2=x^3-3x+a$, for $a=6,-6$, with $j=-216$;
the later choice of a number field $L$ is a restriction of this
representation. Thus there is no unresolved descent condition.
Finding G06 applies to these actual factors, for which
$\operatorname{rad}(2N_{F_a})=6$. Under $\mathrm H_M$, the mod-$\ell$
image over $\mathbb Q$ is $\operatorname{GL}_2(\mathbb F_\ell)$ for
$\ell>964\log6+5760$, and its restricted image over $L$ has index
at most $[L:\mathbb Q]$. Full image over $L$ is not asserted.
Independently of $\mathrm H_M$, the cited Lombardo theorem over
$\mathbb Q$ and the same restriction argument give
\[
 [\operatorname{GL}_2(\widehat{\mathbb Z}):\rho_{F_a,\infty}(G_L)]
 <[L:\mathbb Q]\,\exp(10^{21483})
       \max\{1,h_F(F_a)\}^{4.8\cdot10^{10}}.
\]
Here $h_F$ is the stable Faltings height in the cited theorem's
convention. These refinements do not change the original rows' status.

LM8 A.46--A.48 use Faltings's crystalline comparison; D.25 uses
Faltings's finiteness of rational points. Neither is a finite
Faltings--Serre trace-comparison test. The flagged SA07 F.3.6 excerpt
instead refers to a Hom theorem, and its full proof is missing.

The primitive roots in LM8 A.54 are generators of the two fixed
finite fields of orders $41$ and $61$; A.55 chooses a root of unity;
D.93 uses that $2$ generates $\mathbb Z_5^\times$. D.31 uses the
Artin braid action. The phrase ``vari\'et\'e \'el\'ementaire d'Artin''
in scan 037 is geometric, not a primitive-root conjecture.
GT C.1.21 takes the least \emph{index} with
$v_\ell((J+4)!)\ge N$. Its displayed definition is not a least-prime
problem, and its full proof remains unavailable.

The retained B.23 conductor enumeration and independent Bezout verifier
were rerun successfully. The latter recomputed the division-polynomial
compositions and the complete Bezout identities at conductors $191$
and $229$, after a missing runtime dependency was supplied. This is a
rerun of an existing independent implementation, not a newly written
one. A separate new elementary implementation independently checked
the six primes, multiplicative orders and all finite differences used
in A06. These events are labelled ``computation or formal certificate'';
targeted source reading is labelled ``used and rechecked in this run''
with the precise scope stated. All events are dated 2026-11-21 and name
GPT-6 Astra Pro; no unavailable effort setting is invented.

\paragraph{Unrecovered sources.}
Exact checkpoint and register names, timestamp variants, and broad
title searches were attempted for the A1-successor and GT returns.
No matching full source was recovered, and no nested matching archive
was present in the recovered bundles. The documented Personal Context
route returned ``Personal context is disabled''; it was blocked,
not a zero-result search. The remaining concrete route is to provide
the two original checkpoint archives or their full registers and papers,
or an exact identity for a differently titled upload.
\texttt{SOURCE\_RECOVERY\_MANIFEST.tsv} retains the exact recovered
identities, creation dates, sizes and SHA-256 hashes without signed
download URLs. The ledger's coverage statement must remain
$800/1190$ registers and $591/1190$ available source blocks until those
missing sources are supplied.

\paragraph{Proof of the new scope assertions R01--R03.}
For R01, the retained inventory parses all original flag records and
all recovered CSV registers, counts distinct run/statement identifiers,
and locates the numbered source blocks. Every claimed recovery has a
source hash; the nine LM7 block identifiers are separately listed.
This proves availability counts, not correctness of every source proof.
For R02 the full A.69 text spells out both genus and rank formulas,
while the other flag literally reports the disabled GRH option.
For the finite-prime claim in R03, Frobenius acts on primitive
$\ell$th roots as multiplication by $p$, so their minimal-polynomial
degrees are exactly $\operatorname{ord}_\ell(p)$. It follows that
$\Phi_\ell$ is irreducible precisely in the stated case, without
producing such a prime. For the displayed pair estimate, with the explicitly defined Lang
eigensheaves geometrically nontrivial, character
orthogonality gives $\sum_bN_b^2=|A|^{-1}\sum_\chi|S_\chi|^2$.
The trivial character contributes $N^2$ and every nontrivial one
contributes at most $(2g-2)^2q$ by the cited finite-field trace bound.
Subtracting $\sum_bN_b=N$ gives the formula. The stronger prime
existence and geometric implications are treated, with full proofs
or exact cited hypotheses, in A02--A06 and G01--G05.
\clearpage
\part{Unchecked inputs and remaining gates}
\section{Unchecked inputs and remaining gates}
\label{sec:remaining}

The following entries are deliberately outside the checked-results
part. They are not assigned a proof status by analogy with a nearby
proved result.

\paragraph{U01: the requested Hooley-based proof.}
A02 proves the arithmetic consequence of M. A03 proves the exact
prime-degree geometric implication, and A11 weakens its sufficient
condition to an odd prime power $n$ with
$\langle2,-1\rangle=(\mathbb Z/n)^\times$ and an exact-order-$n$
theta point. A15 and A18 weaken this further: in characteristic two and every
genus at least two, a cyclic cover of any odd degree with primitive
factor of geometric $2$-rank zero suffices. In genus two this is
equivalent to theta meeting every primitive doubling orbit of the
defining cyclic torsion subgroup. A18 removes the exceptional
elliptic double quotient by an exact Deuring--Shafarevich and
conductor-rank calculation.
A04 refutes the inference from
two independent density statements, and A06 refutes the universal
joint primitive-root theta assertion itself. A05 proves A1 by a
different route. No claim is made here that a proof of A1 through
Hooley has been found. A strategy restricted to positive-two-rank
genus-two curves still requires a curve-dependent incidence theorem.
A07 refutes the positive-rank version over finitely generated fields
with transcendental moduli, not over finite fields. A08 shows that
no fixed finite prime list works over all finite fields; in the
ordinary case it excludes any fixed finite list of odd degrees.
A09 supplies an exact degree-$13$ rank-one witness. A13 supplies an
ordinary example of prime degree $4159$, with $2$ of index two in
the unit group. A16 supplies a rank-one example of degree $15$
whose full new factor has positive $2$-rank but whose primitive
factor has rank zero. A10 excludes order-three theta incidence on every nonordinary
genus-two base. None of
these results interchanges the finite-list and all-primes quantifiers;
in higher genus it requires nilpotence of every relevant Frobenius
orbit composite, not merely one nonordinary arrow.
A18 now permits any odd conductor with geometric primitive rank zero.
A19 makes all exponents for each specified finite odd prime support
an exact finite search. A21--A22 show that neither generic
positive rank nor any fixed finite support forces success, even
when mixed orders are allowed. A23 refutes universal success of
the older large-prime quotient inequality over constant extensions.
The remaining existence assertion must therefore allow the prime
support to depend on the particular finite-field curve, or use a
separating construction outside these sufficient tests.

A26 weakens the separate rational quotient criterion using even
augmentation. A27 nevertheless gives a rank-one curve over $\mathbf F_4$
where its numerical hypothesis fails for every extension and every
primitive-two prime. A28 recalls the recovered canonical-class criterion,
and A29 strictly improves its prime-level detection: a nonzero
$r$-primary component of $[K_C-2P]$ suffices if
$\operatorname{ord}_r(p)$ is even, with a cyclic $r^a$ cover and a
prime-to-$p$ rational top Prym. This is unconditional and needs no
prime-distribution theorem. The actual-curve strictness example in A29
shows why passage to prime powers matters even when all prime-level
canonical detectors fail. The residual points have canonical class
supported on the characteristic prime and odd primes of odd
multiplicative order, with the prime-two issue separate in odd
characteristic. In characteristic two this includes every Weierstrass
point and the purely two-primary non-Weierstrass examples of A09.
Constant extensions preserve that primary support. Other covers may
still work; no absence conclusion follows from this residual condition.

A30 refutes a proposed direct higher-genus extension of the
``every detecting quotient is good'' conclusion: a genus-three
rank-one curve has a detecting constant cyclic cubic quotient whose
relative class number is even. It does not exclude other good
quotients. The proposed unramified double-cover repair is unavailable
to the geometric prime-to-two fundamental group and does not preserve
the descending detector. A31 shows that the rank-one curve of A27
nevertheless has an exact signed-doubling theta witness of order seven.
All six of its original rational points permanently miss A29's
canonical detector, yet the geometric witness can be pointed after
a finite extension. Thus the sufficient tests have distinct limits;
their failure is not a nonexistence result.

The approaches actually tested were: fixed-base Hooley density;
Goldfeld's average over bases; Lenstra's fixed finite Chebotarev
constraints and positivity; prime-to-characteristic torsion orders
from Pop--Sa\"idi; a single singular Frobenius arrow in higher genus;
scalar-Frobenius packet counts; generic ordinary and fixed-rank
Prym theorems with finite-type specialization; a direct degree-$13$
theta certificate; prime-power character blocks; primary torsion
projections on theta; the Lang quotient and Weil bound; growth of
fixed-prime primary torsion; the composite primitive projector;
an exact order-$15$ search and certificate; the Hooley--Lenstra
family $n=3q$; the elliptic-exception rank calculation; finite
prime-support torsion via Frobenius translation orbits; the
all-three-primary obstruction; constant twists with an unspecified
split base point; the parity augmentation test; its all-extension rank-one obstruction;
inverse-character canonical detection and its prime-power refinement;
and the independent cuspidal route.
Their respective endpoints are recorded in A02--A31 and the
retained B.28--B.29 pointing audit. Poonen's
projection theorem supplies an order $\ell^r m$, with uncontrolled
prime factors of $m$; multiplying by $m$ moves theta, and passing to
an intermediate cover leaves proper-conductor character blocks.
Those are failures of those two strategies. A15 bypasses the
proper-conductor factors by a product projector on its stated
geometric locus, and A18 removes its geometric exception in
characteristic two. What remains is the same-subgroup orbit-coverage
condition, or the stronger single-point unit-group condition.
A17 gives density $\mathcal A/2$ to the candidate primes
$q\equiv5\pmod8$ for $n=3q$, conditionally on M, without proving
theta incidence for any particular curve. A12 gives an unconditional
curve-dependent primary-torsion test; A14 proves that enlarging the
constant field at one fixed prime cannot eventually make that
particular sufficient test succeed. Neither statement is a theorem
of nonexistence of theta torsion. The recovered Successor A.10
rational prime-cover criterion is separately distinguished in A14.
A19 now makes all primary exponents a finite exact search for
each prescribed finite odd prime support. A20 proves that even
all powers of three fail on the rank-one curve of A09, which
nevertheless has its order-thirteen witness. The prime supports
remain unbounded; no theorem proves that one of these finite
searches succeeds for every positive-rank curve.
The B.28--B.29 rational-twist alternative supplies a good split
point without prescribing it. Its good-point density estimate
does not force such a point into an arbitrarily prescribed fibre,
and extending a good twist can destroy its rational torsion
condition. These failures are recorded separately from geometric
primitive-rank-zero incidence.
The cuspidal route
settles point separation, while the scalar route refutes the
universal auxiliary existence assertion. None of the other routes
silently fills the curve-dependent incidence gap.

\paragraph{U02: the literal primitive-cover problem.}
ANABELIAN\_A1 B.24 allows the degree-one cover with convention
$P_1=J_C$. The B.23 counterexample has $p$-rank zero and is therefore
covered by that convention. The problem for positive $p$-rank
remains open in the recovered run. Similarly ANABELIAN\_A1 B.5 and
B.32 ask for prescribed-point rational class-number witnesses. A29
now settles the stated nonzero canonical-primary component cases,
while their universal existence remains open.
A bound for a prime in a number-field Frobenius class does not
produce the required covering geometry or local unit conditions.
No assertion is made that these separate questions follow from A05.

\paragraph{U03: complete source coverage.}
The two missing registers account for 390 full statements. Only
47 distinct A1-successor IDs and 50 distinct GT-tower IDs are visible
in the supplied summaries and flag excerpts. The missing LM7
proof files account for another 209 proof blocks. Exact-title,
variant-title and content searches recovered four relevant archives,
but not the two missing sources. A final prior-context recovery
attempt was unavailable because that retrieval service was disabled.
The next concrete input is the complete
\texttt{ANABELIAN\_A1\_SUCCESSOR} register/paper and the complete
\texttt{GT\_TOWER\_CHECK} register/paper, together with the omitted
LM7 proof sections. Their absence is an evidentiary gate, not a
mathematical theorem beyond GRH.

\paragraph{U04: infinite towers and unspecified finite data.}
For a recovered row with finite affine monodromy, G01 gives a
prime bound after the actual finite image, cocycle, precision,
and excluded primes are fixed. It does not bound the precision
needed to separate an infinite object, the height of a Hilbert
specialization, or a Baire-category construction. A finite test of
full adic semisimplifications is available under the hypotheses of
G03, with exact computable traces as an additional algorithmic input.
Its applicability to an unspecified row must be checked separately.

\paragraph{U05: scope of the effective open-image bounds.}
The former descent gate for LM8 A.32--A.33 is closed by their actual
models $F_{\pm6}/\mathbb Q$ and the checked refinement G06.
The conditional Mayle--Wang bound applies over $\mathbb Q$; after
restriction to the proof's number field $L$, G06 proves an index
bound $[L:\mathbb Q]$, not necessarily full image.
G-LOMBARDO is an unconditional explicit adelic index bound over
arbitrary number fields. The hypotheses must still be checked for
any other factor; CM factors have CM monodromy rather than full
$\operatorname{GL}_2$ monodromy. No corresponding general abelian-variety
bound has been asserted.

\paragraph{U06: boundaries of verification.}
The complete new proofs were read, and every included computational
certificate reports the exact execution actually performed. No
formal proof-assistant verification of this paper, reconstruction
of all earlier profinite elements, replay of every historical
computer search, or verification of Masullo's entire draft is
claimed. The check table distinguishes whole-proof reading,
independent arithmetic reimplementation, retained-verifier replay,
source-statement checks and carried background. A same-model second
or third reader is not presented as verification by a different
model or by a theorem prover.

\clearpage
\begin{thebibliography}{99}

\bibitem{Alpoge}
L.~Alp\"oge, \emph{Remark}, supplied four-page note dated 20 November
2026, especially step~20, p.~4; supplied
\texttt{ALPOGE\_Main.lean}, \texttt{ALPOGE\_Statements.lean}, and
\texttt{ALPOGE\_LEAN\_SCAN.tsv}. \url{https://alpo.ge/remark.pdf}.
Accepted as a cited theorem by the task's explicit convention.

\bibitem{Masullo}
A.~M.~Masullo, \emph{The Music of the Primes}, supplied 940-page draft,
20 November 2026. Claimed scope: \S1.2; trace formulas:
Theorems~12.5.1.1 and 13.3.1.1; unfinished intended main-statement
heading: Theorem~13.3.1.2, p.~522. All consequences in this report
are conditional on the draft's stated claim M.

\bibitem{Cote149}
A.~Grothendieck, anabelian manuscript, cote~149, supplied
transcriptions \texttt{ARCHIVE\_149\_scan033.tex} through
\texttt{ARCHIVE\_149\_scan038.tex}. No original manuscript images
were supplied or inspected in this run.

\bibitem{Intake}
\emph{INTAKE\_ANABELIAN\_149\_FINAL\_20261114}, supplied intake,
14 November 2026; especially the A.C2.5/A.C2.7/A.C2.19/A.C2.20
review, \S3.3(3) and \S5(a), and the D.30--D.31 dependency discussion.

\bibitem{A1run}
\emph{ANABELIAN\_A1}, checkpoint~03, recovered complete manuscript
\texttt{ANABELIAN\_A1\_paper.tex}, 121-row register and supporting
sources. Exact archive identity and SHA-256 hashes appear in
\texttt{SOURCE\_RECOVERY\_MANIFEST.tsv}. The proof of A1 uses
A.2--A.5, A.14--A.19 and A.30--A.33; the explicit curve is B.23.
Source labels alone are not treated as proof checks.

\bibitem{ASrun}
\emph{ANABELIAN\_SUCCESSOR}, recovered final manuscript
\texttt{ANABELIAN\_SUCCESSOR\_paper.tex} and 157-row register;
especially A.17, A.38 and the prime-to-characteristic covering
criteria. Archive identity and source hashes are retained.

\bibitem{LM7run}
\emph{LM7\_PUSH}, checkpoint~09, 13 November 2026, 218-row register.
Recovered from the user's related
\texttt{DELTA\_2-3\_reciprocity\_of\_frames\_on\_GT} bundle,
with only the proof blocks E.1--E.7, E.58 and E.65.
The complete register is available; other proof blocks are not.

\bibitem{LM8run}
\emph{LM8\_PUSH}, checkpoint~17, 16 November 2026,
304-row register and complete main proof sources. Exact section
filenames and statement IDs are recorded in the row audit and
\texttt{GT\_LEVEL\_SCOPE.tsv}. Earlier source computations are not
claimed rerun unless an event records it.

\bibitem{SGA1}
A.~Grothendieck et al., \emph{Rev\^etements etales et groupe fondamental
(SGA~1)}, Lecture Notes in Mathematics 224, Springer, 1971;
recomposed edition, Documents Mathematiques 3, SMF, 2003.
Exact Expos\'e/statement locators are given in the text.
\url{https://arxiv.org/abs/math/0206203}.

\bibitem{WeilI}
P.~Deligne, \emph{La conjecture de Weil. I}, Publications
Mathematiques de l'IHES \textbf{43} (1974), 273--307,
formula~(1.5.1), Theorem~1.6 and Frobenius convention~(1.15.1).
\url{https://www.numdam.org/item/PMIHES_1974__43__273_0.pdf}.

\bibitem{MilneLEC}
J.~S.~Milne, \emph{Lectures on Etale Cohomology}, version~2.21,
22 March 2013; Example~7.9(a)--(b), Remark~7.10,
Proposition~14.2, Theorem~14.9 and Remark~14.10.
\url{https://www.jmilne.org/math/CourseNotes/LEC.pdf}.

\bibitem{Conrad}
B.~Conrad, \emph{Chow's $K/k$-image and $K/k$-trace, and the
Lang--Neron theorem}, author version dated 17 June 2006,
Theorem~2.1. \url{https://math.stanford.edu/~conrad/papers/Kktrace.pdf}.

\bibitem{ST2009}
M.~Sa\"idi and A.~Tamagawa, \emph{A prime-to-$p$ version of
Grothendieck's anabelian conjecture for hyperbolic curves over finite
fields of characteristic $p>0$}, Publications of RIMS
\textbf{45} (2009), 135--186; Remark~1.2,
Definitions~1.3, 1.5, 1.10, Lemma~1.7 and Proposition~1.8(vi).
\url{https://doi.org/10.2977/PRIMS/1234361157}.
Primary published-version PDF checked at
\url{https://empslocal.ex.ac.uk/people/staff/ms220/Site/Publications_files/publRIMS.pdf}.

\bibitem{PopSaidi}
F.~Pop and M.~Sa\"idi, \emph{On the specialization homomorphism
of fundamental groups of curves in positive characteristic},
in \emph{Galois Groups and Fundamental Groups}, MSRI Publications
\textbf{41} (2003), 107--118; Proposition~2.1,
Propositions~2.3--2.4.
\url{https://library.slmath.org/books/Book41/files/pop.pdf}.

\bibitem{Saidi2005}
M.~Sa\"idi, \emph{On complete families of curves with a given
fundamental group in positive characteristic}, Manuscripta
Mathematica (2005), primary author PDF, Proposition~2.2 and
Theorem~2.3. \url{https://empslocal.ex.ac.uk/people/staff/ms220/Site/Publications_files/manuscripta.pdf}.
These statements concern theta avoidance and ordinarity of covers;
they do not supply the joint primitive-root theta incidence in A03.

\bibitem{Hooley}
C.~Hooley, \emph{On Artin's conjecture}, Journal f\"ur die reine
und angewandte Mathematik \textbf{225} (1967), 209--220,
\S7, concluding theorem.
\url{https://doi.org/10.1515/crll.1967.225.209}.
The original English publisher PDF was unavailable; the theorem and
field hypotheses were checked in the authorized Russian translation,
\emph{Matematika} \textbf{12}:1 (1968), 37--47,
\S7, p.~47, and in the modern numbered form below.
\url{https://www.mathnet.ru/eng/mat455}.

\bibitem{FanPollack}
K.~(S.)~Fan and P.~Pollack, \emph{Counting primes with a given
primitive root, uniformly}, arXiv:2505.05601v2, 14 September 2025,
Theorem~1.1, specialized to the fixed base~2.
\url{https://arxiv.org/abs/2505.05601}.

\bibitem{Lenstra1977}
H.~W.~Lenstra, Jr., \emph{Artin's conjecture on primes with
prescribed primitive roots}, Seminaire Delange--Pisot--Poitou,
Theorie des nombres \textbf{18}:1 (1976--1977), expose~14, 1--8;
Theorems~3.1, 4.1--4.2 and~5.1.
\url{https://www.numdam.org/item/SDPP_1976-1977__18_1_A11_0.pdf}.

\bibitem{Goldfeld}
D.~Goldfeld, \emph{Artin's conjecture on the average},
Mathematika \textbf{15} (1968), 223--226; theorem, p.~223.
\url{https://www.math.columbia.edu/~goldfeld/ArtinsConjecture.pdf}.

\bibitem{Bach1990}
E.~Bach, \emph{Explicit bounds for primality testing and related
problems}, Mathematics of Computation \textbf{55} (1990),
355--380; Theorems~2--3, p.~372, and generator consequence, p.~373;
Dirichlet scope, \S6, pp.~369--370.
\url{https://doi.org/10.1090/S0025-5718-1990-1023756-8}.

\bibitem{LLS2015}
Youness Lamzouri, Xiannan Li, and Kannan Soundararajan,
\emph{Conditional bounds for the least quadratic non-residue and related problems},
Mathematics of Computation \textbf{84} (2015), no.~295, 2391--2412.
\S1.1 explicitly distinguishes Bach's constants two and three.
\url{https://doi.org/10.1090/S0025-5718-2015-02925-1}.
For the general-subgroup bound, the corrected author version checked
is arXiv:1309.3595v3, 8~December~2016:
\url{https://arxiv.org/pdf/1309.3595v3}.
The authors' revision notice at
\url{https://arxiv.org/abs/1309.3595} identifies the corrections to
asymptotic Theorems~1.2--1.3. Theorem~1.1's displayed formula is unchanged.

\bibitem{BachSorenson}
E.~Bach and J.~Sorenson, \emph{Explicit bounds for primes in
residue classes}, Mathematics of Computation \textbf{65} (1996),
1717--1735; Theorem~5.1, p.~1729; analytic hypothesis in the
last paragraph of \S2, p.~1720; Theorem~3.2 and Corollary~3.3.
\url{https://doi.org/10.1090/S0025-5718-96-00763-6}.

\bibitem{Serre1981}
J.-P.~Serre, \emph{Quelques applications du theoreme de densite
de Chebotarev}, Publications Mathematiques de l'IHES
\textbf{54} (1981), 123--201; Proposition~6 (discriminant bound),
\S2.1, Theorem~1 (qualitative Chebotarev).
\url{https://www.numdam.org/item/PMIHES_1981__54__123_0.pdf}.

\bibitem{MayleWang}
J.~Mayle and T.~Wang, \emph{On the effective version of Serre's
open image theorem}, Bulletin of the London Mathematical Society
\textbf{56} (2024), 1399--1416; Theorems~1, 2 and~5,
Corollary~6, Lemma~7 and Proposition~12.
\url{https://doi.org/10.1112/blms.13002}.

\bibitem{MilneCFT}
J.~S.~Milne, \emph{Class Field Theory}, version~4.03,
6 August 2020; Chapter~V, Theorem~3.23, printed p.~164
(qualitative Chebotarev).
\url{https://www.jmilne.org/math/CourseNotes/CFT.pdf}.

\bibitem{Lombardo2015}
D.~Lombardo, \emph{Bounds for Serre's open image theorem for elliptic
curves over number fields}, Algebra \& Number Theory \textbf{9}
(2015), no.~10, 2347--2395; Corollary~9.3, p.~2393;
isogeny input Theorem~2.1. The alternative constants of Remark~9.4
are not used.
\url{https://doi.org/10.2140/ant.2015.9.2347}.
\url{https://msp.org/ant/2015/9-10/ant-v9-n10-p06-p.pdf}.

\bibitem{OzmanPries2019}
E.~Ozman and R.~Pries,
\emph{Ordinary and almost ordinary Prym varieties},
Asian J. Math. \textbf{23} (2019), no.~3, 455--478.
DOI: \href{https://doi.org/10.4310/AJM.2019.v23.n3.a5}{10.4310/AJM.2019.v23.n3.a5}.
Theorem~1.1(2), p.~456; Theorem~8.3, p.~473.

\bibitem{Poonen2005}
B.~Poonen,
\emph{Multiples of subvarieties in algebraic groups over finite fields},
Int. Math. Res. Not. 2005, no.~24, 1487--1498.
Theorem~1.8.
\url{https://math.mit.edu/~poonen/papers/multiples.pdf}.

\bibitem{Yang2021}
Y.~Yang,
\emph{Raynaud--Tamagawa theta divisors and new-ordinariness of
ramified coverings of curves},
J. Algebra \textbf{587} (2021), 263--294.
DOI: \href{https://doi.org/10.1016/j.jalgebra.2021.07.032}{10.1016/j.jalgebra.2021.07.032}.
Theorem~1.1; Definition~3.3 and Theorem~3.5;
Section~3.6.2 and Corollary~3.6.

\bibitem{PriesZhu2012}
R.~Pries and H.~J.~Zhu, \emph{The $p$-rank stratification of
Artin--Schreier curves}, Ann. Inst. Fourier \textbf{62} (2012),
707--726, Lemma~2.6, p.~712.
\url{https://www.numdam.org/item/10.5802/aif.2692.pdf}.

\bibitem{EGAIII2}
A.~Grothendieck, with J.~Dieudonn\'e, \emph{Elements de geometrie
algebrique. III: Etude cohomologique des faisceaux coherents,
seconde partie}, Publ. Math. IHES \textbf{17} (1963), 5--91.
Theorem~7.7.5(I), p.~67; Theorem~7.9.4, pp.~76--77.
\url{https://www.numdam.org/item/PMIHES_1963__17__5_0.pdf}.

\bibitem{EGAIV3}
A.~Grothendieck, with J.~Dieudonn\'e, \emph{Elements de geometrie
algebrique. IV: Etude locale des schemas et des morphismes de
schemas, troisieme partie}, Publ. Math. IHES \textbf{28} (1966),
5--255. Theorem~8.8.2(i)--(ii), p.~28; Theorem~9.7.7(ii), p.~79.
\url{https://www.numdam.org/item/PMIHES_1966__28__5_0.pdf}.

\bibitem{Stacks}
The Stacks Project, Lemma~32.10.1(1)--(3), Tag~\texttt{01ZM};
Lemma~36.32.1, Tag~\texttt{0BDN}; Lemma~30.22.1,
Tag~\texttt{07VK}. Checked 21 November 2026.
\url{https://stacks.math.columbia.edu/}.

\bibitem{Poonen2007Gonality}
B.~Poonen,
\emph{Gonality of modular curves in characteristic $p$},
Math. Res. Lett. \textbf{14} (2007), no.~4, 691--701.
DOI: \href{https://doi.org/10.4310/MRL.2007.v14.n4.a14}{10.4310/MRL.2007.v14.n4.a14}.
Proposition~2.1, Lemma~2.2, and Corollary~2.3, author manuscript p.~3.
\url{https://math.mit.edu/~poonen/papers/gonality.pdf}.

\bibitem{AnbarBeelenNguyenP}
N.~Anbar, P.~Beelen, and N.~Nguyen,
\emph{A new tower with good $p$-rank meeting Zink's bound},
Acta Arith. \textbf{177} (2017), no.~4, 347--374.
DOI: \href{https://doi.org/10.4064/aa8388-6-2016}{10.4064/aa8388-6-2016}.
Theorem~1, equation~(1), author manuscript p.~3.
\url{https://backend.orbit.dtu.dk/ws/files/123787793/p_rank_abn.pdf}.

\bibitem{MilneJV}
J.~S.~Milne, \emph{Jacobian Varieties}, author version of
12 June 2021 of the chapter in \emph{Arithmetic Geometry}
(G.~Cornell and J.~H.~Silverman, eds.), Springer, 1986, 167--212.
Section~9, Proposition~9.1, p.~30, with the characteristic-primary
completion on pp.~31--32; Section~11, Theorem~11.1 and
Corollary~11.4, pp.~35--37.
\url{https://jmilne.org/math/xnotes/JVs.pdf}.

\bibitem{ScanlonVoloch1999}
T.~Scanlon and J.~F.~Voloch,
\emph{Difference algebraic subgroups of commutative algebraic
groups over finite fields},
Manuscripta Math. \textbf{99} (1999), no.~3, 329--339.
Introduction, Section~3, and reference \textnormal{[Box]}.
\url{https://math.berkeley.edu/~scanlon/papers/dag24n.pdf}.
The citation supplies literature context; the specialized
finite-support theorem in A19 is proved directly here.


\bibitem{MilneAV}
J.~S.~Milne, \emph{Abelian Varieties}, author version of
2 January 2022 of the chapter in \emph{Arithmetic Geometry},
Springer, 1986, 103--150.
Section~19, Theorem~19.1(a)--(c) and its Frobenius explanation,
pp.~40--41; Section~12, Proposition~12.9.
\url{https://jmilne.org/math/xnotes/AVs.pdf}.


\bibitem{AchterPries2008}
J.~D.~Achter and R.~Pries,
\emph{Monodromy of the $p$-rank strata of the moduli space of curves},
International Mathematics Research Notices 2008,
Art.~rnn053, 25~pp.
DOI: \url{https://doi.org/10.1093/imrn/rnn053}.
Author version:
\url{https://www.math.colostate.edu/~achter/math/prm0508.pdf}.
Theorem~4.5, author p.~14; Section~5.1 and Lemmas~5.1--5.2,
author pp.~17--18.

\bibitem{LSTX2017}
A.~Landesman, A.~A.~Swaminathan, J.~Tao, and Y.~Xu,
\emph{Lifting subgroups of symplectic groups over
$\mathbf Z/\ell\mathbf Z$},
Research in Number Theory \textbf{3} (2017), Art.~14.
DOI and published primary text:
\url{https://doi.org/10.1007/s40993-017-0078-6}.
Theorem~1 and Proposition~1(b), using the published numbering.

\bibitem{Hall2015}
J.~I.~Hall, \emph{Classical Groups and Geometry},
author notes, version 6 May 2015.
\url{https://users.math.msu.edu/users/halljo/classes/MTH913_S15/cgg_2015-05-06.pdf}.
Theorem~12.26, printed p.~173 (PDF page~179).

\bibitem{Poonen2017RPV}
B.~Poonen, \emph{Rational points on varieties},
Graduate Studies in Mathematics \textbf{186},
American Mathematical Society, Providence, 2017.
Author version: \url{https://math.mit.edu/~poonen/papers/Qpoints.pdf}.
Theorem~7.7.1(ii)--(iii), printed p.~222;
Theorem~4.3.5 and Corollary~4.4.6, printed pp.~101 and~104.

\bibitem{Stacks0245}
The Stacks Project Authors,
\emph{The Stacks Project}, Lemma~35.37.1,
Tag~\texttt{0245}, effective descent for affine morphisms.
\url{https://stacks.math.columbia.edu/tag/0245}.

\end{thebibliography}

\end{document}
